%!TEX TS-program = xelatex
%!TEX root = ../Topo1_WS14/topologie_1.tex
\RequirePackage[thinlines]{easybmat} %-- muss aufgrund eines Bugs vor etex und tikz geladen werden
\documentclass[a4paper, twoside, headsepline, index=totoc,toc=listof, fontsize=10pt, cleardoublepage=empty, headinclude, DIV=12, BCOR=10mm, titlepage]{scrartcl}
\usepackage{scrtime} % KOMA, Uhrzeit ermoeglicht
\usepackage{scrpage2} % wie fancyhdr, nur optimiert auf KOMA-Skript, leicht andere Syntax
\usepackage{etoolbox}
\usepackage{letltxmacro}
\usepackage{ifthen}

%-- Compiler-abhaengige IFs
\usepackage{ifxetex,ifluatex}
\newif\ifxetexorluatex
\ifxetex
  \xetexorluatextrue
\else
  \ifluatex
    \xetexorluatextrue
  \else
    \xetexorluatexfalse
  \fi
\fi

%--Farbdefinitionen
%-- muss vor tikz geladen werden
\usepackage[usenames, table, x11names]{xcolor}
\definecolor{dark_gray}{gray}{0.45}
\definecolor{light_gray}{gray}{0.6}

%--Zum Zeichnen
%-- muss vor polyglossia bzw. babel geladen werden
\usepackage{tikz}
\usepackage{tikz-cd}
\usetikzlibrary{external}
\tikzset{>=latex}
\usetikzlibrary{shapes,arrows,intersections}
\usetikzlibrary{calc,3d}
\usetikzlibrary{decorations.pathreplacing,decorations.markings, decorations.pathmorphing}
\usetikzlibrary{angles}

%-- Konfiguration von tikzexternalize
\tikzexternalize[prefix=tikz/, up to date check=diff]

\ifxetexorluatex
\tikzset{external/system call={xelatex \tikzexternalcheckshellescape %-- verwende LuaLaTeX, wegen dynamischer Speicherallokation
    -halt-on-error -interaction=batchmode -jobname "\image" "\texsource"}}
\else
\tikzset{external/system call={pdflatex \tikzexternalcheckshellescape 
    -halt-on-error -interaction=batchmode -jobname "\image" "\texsource"}}
\fi

%-- tikzexternalize fuer tikzcd deaktivieren, da inkompatibel
\AtBeginEnvironment{tikzcd}{\tikzexternaldisable}
\AtEndEnvironment{tikzcd}{\tikzexternalenable}

%-- um Inkompatibilitaeten von quotes und polyglossia bzw. babel zu vermeiden
\tikzset{
  every picture/.append style={
    execute at begin picture={\shorthandoff{"}},
    execute at end picture={\shorthandon{"}}
  }
}
\usetikzlibrary{quotes}
\usepackage{pgfplots}
\usepgfplotslibrary{colormaps}
\pgfplotsset{compat=1.10}



%-- Mathesymbole etc.
\usepackage{mathtools} % beinhaltet amsmath
\mathtoolsset{showonlyrefs, centercolon}
\newtagform{brackets}[\textbf]{[}{]}
\usetagform{brackets}
\usepackage{wasysym} % zusätzliche Symbole
\usepackage{amssymb} % zusätzliche Symbole
\usepackage{latexsym} % zusätzliche Symbole
\usepackage{stmaryrd} % für Blitz
\usepackage{nicefrac} % schräge Brüche
\usepackage{cancel} % Befehle zum Durchstreichen
\usepackage{mathdots} % Verbesserung von Punkten wie zB \ldots
\usepackage{mathrsfs} % das X wurde benötigt

%-- einzelne Symbole aus dem mathx-package
\DeclareFontFamily{U}{mathx}{\hyphenchar\font45}
\DeclareFontShape{U}{mathx}{m}{n}{<-> mathx10}{}
\DeclareSymbolFont{mathx}{U}{mathx}{m}{n}
\DeclareMathSymbol{\bigplus}{1}{mathx}{"90}

%-- 
\DeclareSymbolFont{bbold}{U}{bbold}{m}{n}
\DeclareSymbolFontAlphabet{\mathbbold}{bbold}
\newcommand{\mathds}[1]{\mathbb{#1}} % Um Kompatibilitaet mit frueheren Benutzung von dsfont herzustellen
\newcommand{\ind}{\mathbbold{1}} % charakteristische-Funktion-Eins

\def\mathul#1#2{\color{#1}\underline{{\color{black}#2}}\color{black}} %farbiges Untersteichen im Mathe-Modus
\renewcommand{\le}{\leqslant}
\renewcommand{\ge}{\geqslant}

%-- Underbrace als Befehl in LaTeX-Syntax (und ohne Spacing probleme mit nachfolgenden Operatoren...)
\newcommand{\Underbrace}[2]{{\underbrace{#1}_{#2}}}

%-- Alles was mit Schrift und XeTeX zu tun hat
\ifxetexorluatex
    % XeLaTeX or LuaTeX
	% \usepackage{eulervm}
	\usepackage[no-math]{fontspec}
    \usepackage{polyglossia} % moderner babel-ersatz
    \setmainlanguage[spelling=new,babelshorthands=true]{german}
	\shorthandoff{"}
	\setotherlanguage{english}
	% \newcommand\glqq{"}
	% \newcommand\grqq{"}
	\defaultfontfeatures{Mapping=tex-text, WordSpace={1.2},Ligatures={Required,Common,Contextual}} %
	\setmainfont{FiraSans}[%
		Extension=.otf,
		Scale=.94,
		UprightFont=*-Book,
		BoldFont=*-Bold,
		ItalicFont=*-LightItalic,
		BoldItalicFont=*-MediumItalic]
	% \setmainfont{SourceSansPro}[%
	% 	Extension=.otf,
	% 	UprightFont=*-Regular,
	% 	BoldFont=*-Semibold,
	% 	ItalicFont=*-LightIt,
	% 	BoldItalicFont=*-SemiboldIt,
	% 	SmallCapsFont={LinLibertine_R.otf},
	% 	SmallCapsFeatures={Letters=SmallCaps}]
	% \setsansfont{Kurier}[Scale=MatchUppercase, Extension=.otf, UprightFont=*Medium-Regular, BoldFont=*Heavy-Regular, ItalicFont=*Medium-Italic,BoldItalicFont=*Heavy-Italic]
	\setsansfont{FiraSans}[%
		Extension=.otf,
		UprightFont=*-Book,
		BoldFont=*-Bold,
		ItalicFont=*-LightItalic,
		BoldItalicFont=*-MediumItalic]
	\setmonofont{FiraMono-Regular.otf}[Scale=MatchLowercase, ItalicFont={*}]
	\newfontfamily\medium[BoldFont=FiraSans-Medium.otf]{FiraSans-Medium.otf}
	% \renewcommand{\familydefault}{\sfdefault} 
	\usepackage{xltxtra}
	\usepackage{microtype}
	\linespread{1.05}
\else
    % default: pdfLaTeX
    \usepackage[ngerman]{babel}
    \usepackage[T1]{fontenc}
    \usepackage[utf8]{inputenc}
	\usepackage{textcomp} %verhindert ein paar Fehler bei den Fonts
    \usepackage[babel=true, tracking=true]{microtype}
	\usepackage{lmodern}
	\renewcommand{\familydefault}{\sfdefault} 
	\providecommand{\medium}{\bfseries}
\fi

\usepackage{fontawesome}

%--Mixed
\usepackage[shortlabels,inline]{enumitem}
\setlist[itemize,1]{label=$\bullet$}
\setlist[enumerate]{font=\medium}
\setlist[description]{font=\normalfont\medium\small}
% \usepackage[neverdecrease]{paralist}
\usepackage[german=quotes]{csquotes}
\usepackage{makeidx}
\usepackage{booktabs}
\usepackage{wrapfig}
\usepackage{float}
\usepackage[margin=10pt, font=small, labelfont={sf, bf}, format=plain, indention=1em]{caption}
\captionsetup[wrapfigure]{name=Abb. }
\usepackage{stackrel}
\usepackage{multicol}

\flushbottom
\KOMAoptions{DIV=last}

%--Unterstreichung
\usepackage[normalem]{ulem}
\setlength{\ULdepth}{1.8pt}

%--Indexverarbeitung
\newcommand{\bet}[1]{\textbf{\emph{#1}}}
\newcommand{\Index}[1]{\textbf{\emph{#1}}\index{#1}}
\makeindex
\setindexpreamble{{\noindent \itshape Die \emph{Seitenzahlen} sind mit Hyperlinks zu den entsprechenden Seiten versehen, also anklickbar}\par \bigskip}
\renewcommand{\indexpagestyle}{scrheadings}

%--Titel abändern
\RedeclareSectionCommand[%
	beforeskip=-10pt plus -2pt minus -1pt,%
	afterskip=-1em plus -2pt minus -1pt,%
	font={\normalsize}%
]{subsection}
\RedeclareSectionCommand[%
	beforeskip=-10pt plus -2pt minus -1pt,%
	afterskip=-1em plus -2pt minus -1pt,%
	font={\normalsize}%
]{subsubsection}
\RedeclareSectionCommand[%
	beforeskip=3.25ex plus 1ex minus .2ex,%
	afterskip=-1em plus -2pt minus -1pt,%
	font={\normalfont\medium\small\scshape}%
]{paragraph}
\renewcommand{\minisec}[1]{\paragraph{#1:}}

%--Marginnote & todonotes
\deffootnote[1.5em]{1.5em}{1.5em}{\textsuperscript{\thefootnotemark}\ }
\usepackage[fulladjust]{marginnote}
\renewcommand*{\marginfont}{\color{dark_gray} \itshape \footnotesize}
\usepackage[textsize=small]{todonotes}
\usepackage{ragged2e}
\renewcommand*{\raggedleftmarginnote}{\RaggedLeft}
\renewcommand*{\raggedrightmarginnote}{\RaggedRight}

%--Todonotes mit tikz/externalize kompatibel machen
\LetLtxMacro{\oldtodo}{\todo}
\renewcommand{\todo}[2][]{\tikzexternaldisable\oldtodo[#1]{#2}\tikzexternalenable}
\LetLtxMacro{\oldmissingfigure}{\missingfigure}
\renewcommand{\missingfigure}[2][]{\tikzexternaldisable\oldmissingfigure[{#1}]{#2}\tikzexternalenable}

%--Konfiguration von Hyperref pdfstartview=FitH, 
\usepackage[hidelinks, pdfpagelabels,  bookmarksopen=true, bookmarksnumbered=true, linkcolor=black, urlcolor=SkyBlue2, plainpages=false,pagebackref, citecolor=black, hypertexnames=true, pdfauthor={Jannes Bantje}, pdfborderstyle={/S/U}, linkbordercolor=SkyBlue2, colorlinks=false]{hyperref}


\newcommand{\appendLink}[1]{#1\,\faExternalLink}
\newcommand{\hrefsym}[2]{\href{#1}{\texttt{\appendLink{#2}}}}
\newcommand{\hrefsymmail}[2]{\href{#1}{\texttt{\faEnvelopeO\,#2}}}
\renewcommand{\url}[1]{\hrefsym{#1}{\nolinkurl{#1}}}

\providecommand{\appendLink}{ }
\providecommand{\hrefsym}{\href}
\providecommand{\hrefsymmail}{\hrefsym}


% -- QR-Codes
\usepackage{qrcode} %-- hinter hyperref laden!

%--Römische Zahlen
\newcommand{\RM}[1]{\MakeUppercase{\romannumeral #1{}}}

%%--Abkürzungen etc.
\newcommand{\light}{\text{\Large $\lightning$}}

%-- Definitionen von weiteren Mathe-Befehlen
\DeclareMathOperator{\re}{Re} %Realteil
\DeclareMathOperator{\im}{Im} %Imaginaerteil
\DeclareMathOperator{\id}{id} %identische Abbildung
\DeclareMathOperator{\Sp}{Sp} %Spur
\DeclareMathOperator{\sgn}{sgn} %Signum
\DeclareMathOperator{\alt}{Alt} %Alternierende n-linearFormen
\DeclareMathOperator{\End}{End} %Endomorphismen
\DeclareMathOperator{\Vol}{Vol} %Volumen
\DeclareMathOperator{\dom}{dom} %Domain
\DeclareMathOperator{\Hom}{Hom} %Homomorphismen
\DeclareMathOperator{\bild}{Bild} %Bild
\DeclareMathOperator{\Kern}{Kern}
\DeclareMathOperator{\ev}{ev} %Auswertungsabbildung
\DeclareMathOperator{\rg}{rg} %Rang
\DeclareMathOperator{\Rg}{Rg} %Rang
\DeclareMathOperator{\diam}{diam} %Durchmesser
\DeclareMathOperator{\dist}{dist} %Distanz
\DeclareMathOperator{\grad}{grad} %Gradient
\DeclareMathOperator{\dive}{div} %Gradient
\DeclareMathOperator{\rot}{rot} %Rotation
\DeclareMathOperator{\hess}{Hess} %Hesse-Matrix
\DeclareMathOperator{\koker}{Koker} %Kokern
\DeclareMathOperator{\coker}{coker} %Kokern
\DeclareMathOperator{\aut}{Aut} %Automorphismen
\DeclareMathOperator{\ord}{ord} %Ordnung
\DeclareMathOperator{\ggT}{ggT} %ggT
\DeclareMathOperator{\kgV}{kgV} %kgV
\DeclareMathOperator{\Gr}{Gr} %Gerade
\DeclareMathOperator{\Kr}{Kr} %Kreis
\DeclareMathOperator{\Char}{char} %Charakteristik
\DeclareMathOperator{\Aut}{Aut} %Automorphismen
% \DeclareMathOperator{\D}{D} %formale Ableitung
\newcommand{\D}{\ensuremath{\mathrm{D}\mkern-1.5mu}}
\newcommand{\mathd}{\ensuremath{\mathrm{d}\mkern-1.0mu}}
\newcommand{\Tmap}{\ensuremath{\mathrm{T}\mkern-1.0mu}}
\newcommand{\opL}{\ensuremath{\mathrm{L}\mkern-0.6mu}}
\DeclareMathOperator{\cov}{cov} %Kovarianz
\DeclareMathOperator{\Gal}{Gal} %Galois
\DeclareMathOperator{\supp}{supp} %Träger
\DeclareMathOperator{\Sym}{Sym} %Symmetrische Gruppe
\DeclareMathOperator{\Zyl}{Zyl} %Zylinder
\DeclareMathOperator{\Mod}{Mod} %Moduln
\DeclareMathOperator{\EPK}{EPK} %Einpunktkompaktifizierung
\DeclareMathOperator{\conj}{conj} %Konjugation
\DeclareMathOperator{\ann}{ann} %Annulator
\DeclareMathOperator{\tor}{tor} %Torsionsmodul
\DeclareMathOperator{\Int}{Int} %Interior/Inneres
\DeclareMathOperator{\Br}{Br} %Brauerergruppe
\DeclareMathOperator{\GL}{GL} %allgemeine lineare Gruppe
\DeclareMathOperator{\SL}{SL} %spezielle lineare Gruppe
\DeclareMathOperator{\SU}{SU} %spezielle unitäre Gruppe
\DeclareMathOperator{\Tr}{Tr} %Spur/Trace
\DeclareMathOperator{\conv}{conv} %Konvexe Teilmenge
\DeclareMathOperator{\Span}{span} %span
\DeclareMathOperator{\Diff}{Diff} %Menge der Diffeomorphismen
\DeclareMathOperator{\Isom}{Isom} %Menge der Isometrien
\DeclareMathOperator{\spec}{spec} %Spektrum
\DeclareMathOperator{\res}{res} %Resolventenmenge
\DeclareMathOperator{\resolv}{R} %Resolventenmenge
\DeclareMathOperator{\idx}{ind} %Fredholm Index
\DeclareMathOperator{\Fred}{Fred} %menge der Fredholm-Operatoren
\DeclareMathOperator{\Landau}{\mathcal{O}} %großes Landau-Symbol
\DeclareMathOperator{\landau}{\scriptstyle \mathcal{O}} %kleines Landau-Symbol

\newcommand{\tr}{\mathrm{tr}}
\newcommand{\w}{\mathrm{w}}
\newcommand{\sa}{\mathrm{s.a.}}
\newcommand{\so}{\mathrm{s.o.}}
\newcommand{\weakT}[1]{\ensuremath{\mathcal{T}_{#1}^{\mkern+1.0mu\text{\raisebox{0.4ex}{$\mathrm{w}$}}}}}
\newcommand{\weakTstar}[1]{\ensuremath{\mathcal{T}_{#1}^{\mkern+1.0mu\text{\raisebox{0.4ex}{$\mathrm{w}$}}^*}}}
\newcommand{\finSub}{\subset\mkern-0.7mu \subset}
\newcommand{\rand}[1]{\ensuremath{\partial^{\scriptscriptstyle #1}}}


% -- Kategorien
\DeclareMathOperator{\Ob}{Ob} %Objekte
\DeclareMathOperator{\obj}{obj} %Objekte
\DeclareMathOperator{\Mor}{Mor} %Morphismen

\newcommand{\TOP}{\textsc{Top}}
\newcommand{\HTOP}{\textsc{HTop}}
\newcommand{\VR}{\textsc{VR}}
\newcommand{\MOD}{\textsc{Mod}}
\newcommand{\MENGEN}{\textsc{Mengen}}
\newcommand{\MAN}{\textsc{Man}}
\newcommand{\GRUPPEN}{\textsc{Gruppen}}
\newcommand{\ABELGRUPPEN}{\textsc{Abel.Gruppen}}
\newcommand{\KAT}{\textsc{Kat}}
\newcommand{\FUN}{\textsc{Fun}}
\newcommand{\SIMP}{\textsc{Simp}}
\newcommand{\VEKT}{\textsc{Vekt}}


\newcommand{\op}{\mathrm{op}}
\newcommand{\ab}{\mathrm{ab}}
\newcommand{\pr}{\mathrm{pr}}
\newcommand{\pt}{\mathrm{pt}}
\newcommand{\const}{\mathrm{const}}
\newcommand{\CW}{CW}
\DeclareMathOperator{\cell}{cell}
\DeclarePairedDelimiter{\homologieklasse}{\llbracket}{\rrbracket}

\newcommand{\Ker}{\ensuremath{\Kern \,}}
\newcommand{\zirkel}{\rotatebox[origin = c]{-90}{$\varangle$}}
\newcommand{\wirkung}{\!\curvearrowright\!}

\newcommand{\oE}{\OE}


%--Beweisende
\newcommand{\bewende}{\ifmmode \tag*{$\square$} \else \hfill $\square$ \fi}

%--Volumen
\newcommand{\vol}[1]{\ensuremath{\Vol_{#1}}}

%--Skalarprodukt
\DeclarePairedDelimiterX\sprod[2]{\langle}{\rangle}{#1\,\delimsize\vert\,#2}
\DeclarePairedDelimiterX\skal[2]{\langle}{\rangle}{#1\,,\,#2}

%--Betrag, Gaußklammer
\DeclarePairedDelimiter{\abs}{\lvert}{\rvert}
\DeclarePairedDelimiter{\floor}{\lfloor}{\rfloor}
\DeclarePairedDelimiter{\ceil}{\lceil}{\rceil}

%--Norm
\DeclarePairedDelimiter\doppelstrich{\Vert}{\Vert}
\newcommand{\norm}[2][\relax]{
\ifx#1\relax \ensuremath{\doppelstrich*{#2}}
\else \ensuremath{\doppelstrich*{#2}_{#1}}
\fi}


%--Umklammern
\DeclarePairedDelimiter\enbrace{(}{)}
\DeclarePairedDelimiter\benbrace{[}{]}
\DeclarePairedDelimiter\lenbrace{<}{>}
\newcommand{\ssbrace}[1]{{\scriptscriptstyle\enbrace{#1}}}

%--Mengen
\DeclarePairedDelimiterX\mengenA[1]{\lbrace}{\rbrace}{#1}
\DeclarePairedDelimiterX\mengenB[2]{\lbrace}{\rbrace}{#1\, \delimsize\vert \, #2}

\makeatletter
\newcommand{\set}[2][\relax]{
\ifx#1\relax \ensuremath{
\mengenA*{#2}}
\else \ensuremath{%
  \mengenB*{#1}{#2}}
\fi}
\makeatother

\DeclareRobustCommand{\minwidthbox}[2]{%
  \ifmmode
    \expandafter\mathmakebox
  \else
    \expandafter\makebox
  \fi
  [\ifdim#2<\width\width\else#2\fi]{#1}%
}


%--offen und abgeschlossen
\newcommand{\off}{\! \stackrel[\text{offen}]{}{\subset} \!}
\newcommand{\abg}{\! \stackrel[\text{abg.}]{}{\subset} \!}

%--Differential
\newcommand{\diff}[2]{\ensuremath{\frac{{\partial #1}}{{\partial #2}} }}
\newcommand{\diffd}[2]{\ensuremath{\frac{\mathd #1}{\mathd #2} }}

%--Diagonale Linien für Matrizen
\newcommand\tikzmark[1]{%
  \tikzexternaldisable\tikz[overlay,remember picture,baseline] \node [anchor=base] (#1) {};\tikzexternalenable}

\newcommand\MyLine[3][]{%
  \tikzexternaldisable\begin{tikzpicture}[overlay,remember picture]
    \draw[#1] (#2.south) -- (#3.north);
  \end{tikzpicture}\tikzexternalenable}
  
%--Jordankasten
\newcommand{\Jord}{ \tikzexternaldisable \begin{BMAT}{|ccc|}{|ccc|} %
				0 & & \\ %
				1\tikzmark{z} &  \diagdown & \\ %
				& \tikzmark{u}1 \MyLine[thick]{z}{u} &  0 %
			\end{BMAT} \tikzexternalenable}
			
%--Jordankasten mit Lambda
\newcommand{\JordLa}[1][\empty]{%
	\tikzexternaldisable\ifthenelse{\equal{#1}{\empty}}
		{\begin{BMAT}{|ccc|}{|ccc|} %
				\lambda & & \\ %
				1\tikzmark{z} &  \diagdown & \\ %
				& \tikzmark{u}1 \MyLine[thick]{z}{u} &  \lambda  %
			\end{BMAT}}
		{\begin{BMAT}{|ccc|}{|ccc|} %
				\lambda_{#1}  & & \\ %
				1\tikzmark{z} &  \diagdown & \\ %
				& \tikzmark{u}1 \MyLine[thick]{z}{u} &  \lambda_{#1}  %
			\end{BMAT}}\tikzexternalenable}

%--Inhaltsverzeichnis
\usepackage[tocindentauto]{tocstyle}
\usetocstyle{KOMAlike}	
\shorthandon{"}
		
%--Punkte (als letztes laden)
\usepackage{ellipsis}

\newcommand{\fach}{Höhere Algebra \RM{1}.}
%\newcommand{\shortFach}{Analysis, Topologie, Geometrie}
\newcommand{\semester}{SoSe 2014}
\newcommand{\homepage}{http://wwwmath.uni-muenster.de/u/franziska.jahnke/ha/}

\newcommand{\prof}{Prof.\ Dr.\ Dr.\ Katrin Tent}
\providecommand{\verfasser}{Jannes Bantje}

%--Konfiguration von scrheadings
\setheadsepline{1pt}[\color{light_gray}]
\pagestyle{scrheadings}
%\fancyhf[H, F]{}
\clearscrheadfoot

\providecommand{\shortFach}{\fach}
\lehead{\includegraphics[height=0.6 cm,keepaspectratio]{../!config/Bilder/Logo_WWU_Muenster_light_gray.pdf}}
\rehead{\rule{0cm}{0.6cm}\footnotesize\sffamily\color{light_gray}\verfasser{} -- Mitschrift \shortFach}
\lohead{\rule{0cm}{0.6cm}\footnotesize\sffamily\color{light_gray}Stand: \today \; \thistime[:]}
\rohead{\includegraphics[height=0.6 cm,keepaspectratio]{../!config/Bilder/fb10logo_gray.pdf}}


\ofoot[{ \color{dark_gray} \LARGE \sffamily \thepage}]{{ \color{dark_gray} \LARGE \sffamily \thepage}} %hier wir auch der plain Stil bearbeitet!
\automark{section}
\ifoot{ \color{dark_gray} \small \leftmark}

%--Metadaten
\providecommand{\mail}{j.bantje@wwu.de}
\author{\verfasser}
\titlehead{\includegraphics[height=1.5cm, keepaspectratio]{../!config/Bilder/Logo_WWU_Muenster.pdf}%
\hfill \includegraphics[height=1.3cm, keepaspectratio]{../!config/Bilder/fb10logo.pdf}}
\title{Skript \fach}
\subtitle{Mitschrift der Vorlesung  \enquote{\fach} von \prof}

\KOMAoptions{DIV=last}


\begin{document}
\pagenumbering{Roman}
\maketitle
\begin{abstract}
\section*{Aktuelle Version verfügbar bei}
\newcommand{\dieBreite}{11cm}
\begin{minipage}{4cm}
	\qrcode[height=3.3cm, version=6]{https://gitlab.com/JaMeZ-B/LaTeX-WWU}
\end{minipage}
\hfill
\begin{minipage}{\dieBreite}
	% \includegraphics[height=0.6cm, keepaspectratio]{../!config/Bilder/wm_no_bg.pdf}
	\includegraphics[height=0.8cm, keepaspectratio]{../!config/Bilder/wm_no_bg.pdf}\\
	\url{https://gitlab.com/JaMeZ-B/LaTeX-WWU} \smallskip\\
	Das zentrale Repository des \enquote{\LaTeX-WWU}-Projekts befindet sich auf der Plattform GitLab.com.
	Neben der Koordination aller Beteiligten werden über diesen Dienst auch die PDFs gebaut, die in der Readme verlinkt sind.
\end{minipage}\\[1cm]
\begin{minipage}{4cm}
	\qrcode[height=3.3cm, version=6]{https://github.com/JaMeZ-B/latex-wwu}
\end{minipage}
\hfill
\begin{minipage}{\dieBreite}
	\includegraphics[height=0.6cm, keepaspectratio]{../!config/Bilder/github_octo.pdf}
	\includegraphics[height=0.6cm, keepaspectratio]{../!config/Bilder/GitHub_Logo.pdf}\\
	\url{https://github.com/JaMeZ-B/latex-wwu} \smallskip\\
	Die Entwicklung des \enquote{\LaTeX-WWU}-Projekts hat ursprünglich auf GitHub stattgefunden, ist mittlerweile aber zu GitLab gewechselt.
	Das GitHub-Repository wird stündlich automatisch aktualisiert, Merge-Requests werden aber nicht mehr entgegengenommen.
\end{minipage}\\[1cm]
% \begin{minipage}{4cm}
% 	\qrcode[height=3.3cm, version=6]{https://uni-muenster.sciebo.de/public.php?service=files&t=965ae79080a473eb5b6d927d7d8b0462}
% \end{minipage}
% \hfill
% \begin{minipage}{\dieBreite}
% 	\raisebox{-2pt}{\includegraphics[height=0.6cm, keepaspectratio]{../!config/Bilder/sciebo_logo.pdf}}
% 	\resizebox{!}{0.5cm}{\large \sffamily\textbf{sciebo}} {\sffamily\large die Campuscloud} \\
% 	\resizebox{\dieBreite}{!}{\footnotesize\url{https://uni-muenster.sciebo.de/public.php?service=files&t=965ae79080a473eb5b6d927d7d8b0462}}\smallskip\\
% 	Sciebo ist ein Dropbox-Ersatz der Hochschulen in NRW, der von der Uni Münster in leitender Position auf Basis der OpenSource-Software Owncloud aufgebaut wurde.
% \end{minipage}\\[1cm]
\hrule \mbox{ }\\[0.7cm]
\begin{minipage}{4cm}
	\qrcode[height=3.3cm, version=6]{\homepage}
\end{minipage}
\hfill
\begin{minipage}{\dieBreite}
	\resizebox{!}{0.5cm}{\large\sffamily\textbf{Vorlesungshomepage}}\\
	\resizebox{\dieBreite}{!}{\footnotesize\url{\homepage}}\smallskip\\
	Hier ist ein Link zur offiziellen Vorlesungshomepage.
\end{minipage}
\newpage
\section*{Vorwort --- Mitarbeit am Skript}
Dieses Dokument ist eine Mitschrift aus der Vorlesung \enquote{\fach, \semester}, gelesen von \prof. 
Der Inhalt entspricht weitestgehend dem Tafelanschrieb. 
Für die Korrektheit des Inhalts übernehme ich keinerlei Garantie! 
Für Bemerkungen und Korrekturen -- und seien es nur Rechtschreibfehler -- bin ich sehr dankbar. 
Korrekturen lassen sich prinzipiell auf drei Wegen einreichen: 
\begin{itemize}
	\item Persönliches Ansprechen in der Uni, Mails an \hrefsymmail{mailto:\mail}{\mail} (gerne auch mit annotieren PDFs) oder Kommentare auf \url{https://gitlab.com/JaMeZ-B/LaTeX-WWU}.
	\item \emph{Direktes} Mitarbeiten am Skript: Den Quellcode poste ich auf GitLab (siehe oben), also stehen vielfältige Möglichkeiten der Zusammenarbeit zur Verfügung:
	Zum Beispiel durch Kommentare am Code über die Website und die Kombination Fork und Merge-Request. 
	Wer sich verdient macht oder ein Skript zu einer Vorlesung, die ich nicht besuche, beisteuern will, dem gewähre ich gerne auch Schreibzugriff.
	
	Beachten sollte man dabei, dass dazu ein Account bei \url{gitlab.com} notwendig ist, der allerdings ohne Angabe von persönlichen Daten angelegt werden kann. 
	Wer bei GitLab (bzw. dem zugrunde liegenden Open-Source-Programm \enquote{\texttt{git}}) -- verständlicherweise -- Hilfe beim Einstieg braucht, dem helfe ich gerne weiter. 
	Es gibt aber auch zahlreiche empfehlenswerte Tutorials im Internet.\footnote{zB. \url{https://try.github.io/levels/1/challenges/1}, ist auf Englisch, aber dafür interaktiv}
	\item \emph{Indirektes} Mitarbeiten: \TeX-Dateien per Mail verschicken. 
	
	Dies ist nur dann sinnvoll, wenn man einen ganzen Abschnitt ändern möchte (zB. einen alternativen Beweis geben), da ich die Änderungen dann per Hand einbauen muss! Ich freue mich aber auch über solche Beiträge!
\end{itemize}
\section*{Literatur}
\begin{itemize}
	\item P.M. Cohn: Basic Algebra, (Further Algebra) Springer
	\item N. Jacobsen: Basic Algebra \RM{1} + \RM{2}
	\item S. Lang : Algebra, Wiley
	\item F. Lorenz: Algebra \RM{3}, Springer
\end{itemize}
\end{abstract}
\tableofcontents
\cleardoubleoddemptypage
\newpage
\pagenumbering{arabic}
\setcounter{page}{1}
\setcounter{footnote}{0}

\section{Gruppentheorie: Wiederholung, Sylow-Sätze, Kompositionsreihen} % (fold)
\label{sec:1}

\subsection{Definition: Gruppenwirkung} % (fold)
\label{sub:11}
\begin{itemize}
	\item Sei $G$ eine Gruppe, $X \not= \emptyset$ Menge. Eine \Index{Gruppenwirkung} von $G$ auf $X$ ist (gegeben durch) einen Gruppenhomomorphismus 
	$\varphi : G \to \Sym(X)$. $\ker \varphi$ heißt \Index{Kern der Wirkung}. Wir schreiben auch kurz $g(x)$ für $\varphi(g)(x)$.
	\item Für $x \in X$ heißt $G_x = \set[g \in G]{g(x) = x} \le G $ der \Index{Stabilisator} von $x$.
	\item Die \Index{Bahn} von $x \in X$ unter $G$ ist $G(x) = \set[g(x)]{g \in G} \subseteq X$. 
	\item Eine Gruppenwirkung heißt \bet{transitiv}\index{Gruppenwirkung!transitive}, wenn $G(x)=X$ für ein $x \in X$.
	\item Eine Gruppenwirkung heißt \bet{treu}\index{Gruppenwirkung!treue}, falls $\ker \varphi = \set{1_G}$. 
\end{itemize}
Für eine alternative Definition siehe Anhang \ref{sub:anh_1} auf Seite \pageref{sub:anh_1}.
% subsection 11 (end)

\subsection[Bemerkung über eine Abbildung $\nicefrac{G}{G_x} \to G(x)$]{Bemerkung} % (fold)
\label{sub:12}
Für jedes $x \in X$ ist die Abbildung $\nicefrac{G}{G_x} \to G(x)$, $g G_x \mapsto g(x)$ eine Bijektion.
\minisec{Beweis}
Es ist $g(x) = h(x) \iff (h ^{-1} g)(x)=(h^{-1}h)(x)=x \iff h ^{-1} g \in G_x \iff g G_x = h G_x$. Daher ist die Abbildung wohldefiniert und injektiv. Surjektiv ist klar. \bewende
% subsection 12 (end)

\subsection*{Wiederholung Isomorphiesätze\footnote{siehe auch \url{http://de.wikipedia.org/wiki/Isomorphiesatz}}} % (fold)
\label{sub:wiederholung_iso}
\begin{description}
	\item[1. Isomorphiesatz:] Ist $\varphi : G \to H$ surjektiv, dann ist $H \simeq \nicefrac{G}{\ker \varphi}$. Allgemein ist für jeden Homomorphismus $\varphi : G \to H$
	dann $\im \varphi \simeq \nicefrac{G}{\ker \varphi}$.\hfill {\color{lightgray}(Homomorphiesatz)}
	\item[2. Isomorphiesatz:] Ist $H \le G, N \unlhd G$, dann ist $\nicefrac{H}{(H \cap N)} \simeq \nicefrac{H N}{N}$.\marginnote{{\footnotesize $H N = \set[h n]{h \in H, n \in N}$} ist
	Untergruppe, da $H N = N H$ für $N$ normal}\hfill {\color{lightgray} ("{}Erweitern mit $N$")}
	\item[3. Isomorphiesatz:] Sind $N,K \unlhd G, N \le K$, dann ist $\nicefrac{\nicefrac{G}{N}}{\nicefrac{K}{N}} \simeq \nicefrac{G}{K}$. \hfill
	{\color{lightgray} ("{}Kürzen mit $N$")}
\end{description}
Die letzten beiden Sätze lassen sich mit dem ersten beweisen!
% subsection wiederholung_isomorphiesätze (end)

\subsection[Beispiele für Gruppenwirkungen]{Beispiel} % (fold)
\label{sub:13}
\begin{enumerate}[(i)]
	\item\begin{enumerate}[(a)]
		\item $G$ wirkt durch Rechtsmultiplikation auf sich selbst ($X = G$). Dann ist $G_x = \set{1}$ für alle $x \in X$, d.h. die Wirkung ist treu und transitiv. Solche 
		Wirkungen heißen \bet{regulär}\index{Gruppenwirkung!reguläre}.
		\[
			\varphi : G \to \Sym(G) , \quad g \mapsto\rho_g \enspace \text{ mit } \enspace\rho_g(x) = x \cdot g
		\]
		Gruppenhomomorphismus\footnote{Beachte: Die Addition in $\Sym(X)$ ist die \emph{linksseitige} 
		Komposition von Abbildungen! Wir wollen also, dass $g \cdot h$ folgendermaßen wirkt: Zuerst $g$ wirken lassen und dann $h$.}: 
		$\rho_{g h} (x) = x \cdot g \cdot h = \rho_h \circ \rho_g (x)$.
		\item $G$ wirkt durch Linksmultiplikation auf sich selbst (regulär) $\lambda_g (x) = g ^{-1} \cdot x$.
	\end{enumerate}
	\item $G$ operiert durch Konjugation auf sich selbst, d.h. $\kappa : G  \to \Aut(G) \le \Sym(G)$, $g \mapsto \kappa_g$, wobei $\kappa_g(x) = g ^{-1} \cdot x \cdot  g$
	\[
		g \cdot  h \mapsto \kappa_{g \cdot h} \quad \kappa_{g \cdot h} (x) = h ^{-1} \cdot g ^{-1}\cdot x \cdot g \cdot h = \kappa_h \circ \kappa_g (x)
	\]
	Dann ist $G_x = \set[g \in G]{g ^{-1} \cdot x \cdot g = x} = Z_G(x) $ der \Index{Zentralisator} von $x$ in $G$. Der Kern der Wirkung ist das \Index{Zentrum}  von $G$ 
	$Z(G) = \set[g \in G]{x \cdot g = g \cdot x \text{ für alle } x \in G}$.
	
	\bet{Bemerkung:} $\ker \varphi = \bigcap_{x \in X} G_x$ gilt für alle Gruppenwirkungen $\varphi : G \to \Sym(X)$.
	\item $G = \mathrm{Gl}_n(K)$, $K$ Körper, operiert auf $K^n$ durch lineare Abbildungen.
\end{enumerate}
% subsection 13 (end)

\subsection{Bahnengleichung} % (fold)
\label{sub:14}
Es gilt: $X = \dot{\bigcup} \set[G(x)]{x \in X} $, denn für jedes $x\in X$ ist $x\in G(x)$. Falls $X$ endlich ist gilt also
\[
	\abs*{X} = \sum \abs*{G(x)} = \sum \abs*{\nicefrac{G}{G_x}} = \sum [G : G_x]. \tag{nach \ref{sub:12}}  
\]
Insbesondere ist $\abs*{G(x)} = \abs*{\nicefrac{G}{G_x}} = [G : G_x] = \frac{\abs*{G} }{\abs*{G_x} }  $ falls $G$ endlich ist. (Bijektion aus \ref{sub:12}) \\
\bet{Spezialfall:} Wirkung von $G$ durch Konjugation auf sich selbst. $\kappa_g(x) = g ^{-1}  \cdot  x \cdot g$.
% subsection 14 (end)

\subsection{Klassengleichung} % (fold)
\label{sub:15}
Sei $K_G = \set[G(x)]{x \in G}$ die Menge der Konjugationsklassen. Sei $K_G^* = \set[G(x)]{x \in G\backslash Z(G)} = \set[G(x)]{ \abs*{G(x)} \ge 2 }  $. Für jede endliche 
Gruppe $G$ gilt dann nach \ref{sub:14}
\[
	\abs*{G} = \sum_{K_G} [G : Z_G(x)] = \abs*{Z(G)} + \sum_{K_G^*} [G : Z_G(x)]. \bewende  
\]
% subsection 15 (end)

\subsection[Korollar: $p$-Gruppen haben eine nichttriviales Zentrum]{Korollar} % (fold)
\label{sub:16}
$G$ endlich, $\abs*{G} = p^m $, $p$ prim, $m \ge 1 \Rightarrow Z(G) \not= 1$. Also haben $p$-Gruppen ein nicht-triviales Zentrum.
\minisec{Beweis}
Nach Lagrange\footnote{Ordnung einer Untergruppe teilt die Gruppenordnung. Die Umkehrung gilt nicht!} gilt für jedes $x \in G$, dass $\abs*{Z_G(x)} = p^k $ für ein 
$k \le m$, also ist $[G : Z_G(x)] = p^{m-k}$. Für $x \notin Z(G)$ gilt dabei insbesondere $k < m$, denn aus $k=m$ würde $x \in Z(G)$ folgen.  Wegen $p \mid \abs*{G} $ und $\abs*{Z(G)} \ge 1 $ folgt $p \mid \abs*{Z(G)} $. \bewende
% subsection 16 (end)

\subsection[Definition: $p$-Sylowgruppe]{Definition} % (fold)
\label{sub:17}
Sei $G$ eine endliche Gruppe, $\abs*{G} = p^a \cdot m $ mit $(m,p) = 1$ und $p$ prim. Dann heißt eine Untergruppe $H \le G$ mit $\abs*{H} = p^a$ eine \bet{$p$-Sylowgruppe}\index{p-Sylowgruppe@$p$-Sylowgruppe} von $G$.
% subsection 17 (end)

\subsection{Satz (Sylow)} % (fold)
\label{sub:18}
Sei $G$ eine endliche Gruppe, $p$ prim, $\abs*{G} = p^a \cdot m $ mit $(p,m) = 1$. Dann gilt
\begin{enumerate}[(i)]
	\item Jede $p$-Untergruppe von $G$ ist in einer $p$-Sylowgruppe enthalten. Insbesondere existieren $p$-Sylowgruppen immer.
	\item Ist $n_p = \#$ $p$-Sylowgruppen von $G$, dann gilt: $n_p \mid m$ und $n_p \equiv 1 \mod p$.
	\item Alle $p$-Sylowgruppen sind konjugiert.
\end{enumerate}
\minisec{Beweis}
Sei $S := \set[X \subset G]{ \abs*{X} = p^a } $. $G$ operiert auf $S$ durch Rechtsmultiplikation. Es ist
\[
	\abs*{S} = \binom{p^a \cdot m}{p^a} = \frac{\cancel{p^a} \cdot m \cdot  (p^a \cdot m -1) \cdot \ldots \cdot \big(p^a \cdot m - (p^a -1)\big)}{1 \cdot 2 \cdot \ldots 
	\cdot (p^a -1) \cdot \cancel{p^a}}.  
\]
Behauptung: $p \nmid \abs*{S} $. Betrachte dazu $k_i := \frac{p^a \cdot m -i}{i}$, für $1 \le i < p^a$. Wenn $p^j \mid p^a \cdot m -i$, dann ist $j < a$ und $p^j \mid i$.
Daher sind $p^a \cdot m -i$ und $i$ durch dieselbe Potenz von $p$ teilbar, d.h. nach vollständigem Kürzen von $k_i$ ist dessen Zähler nicht durch $p$ teilbar. Damit ist $p \nmid m \cdot k_1 \cdots k_{p^a-1} = \abs*{S}$.\\
Daher existiert eine $G$-Bahn $S_1 \subseteq S$ mit $p \nmid \abs*{S_1} $. Wähle $X \in S_1$, d.h. $\abs*{X}= p^a$. Setze $P := G_X$. Dann ist \marginnote{siehe Bahnengleichung \ref{sub:14}}
\[
	\abs*{S_1} = [G: G_X] = [G: P] 
\]
Daher gilt $p \nmid \abs*{\nicefrac{G}{P}}$, also $p^a \mid \abs*{P}$. Andererseits ist $\abs*{P} \le p^a$, denn für $x \in X, g \in P$ ist $x \cdot g \in X$ und die 
$x \cdot g$ für $g \in P$ sind paarweise verschieden. Daher ist $\abs*{P} = p^a $ und $P$ eine $p$-Sylowgruppe. 

Sei nun $T \subseteq S$ die Menge aller Konjugierten von $P$ unter der Konjugationswirkung. Dann operiert auch $P$ durch Konjugation auf $T$. Nach der Bahnengleichung
(\ref{sub:14}) hat 
jede Bahn die Länge $p^i$ für ein $i \le a$. Offensichtlich ist $P$ ein Fixpunkt dieser Wirkung. Ist $P_1 \in T$ ein weiterer Fixpunkt, dann ist 
$P \subseteq N_G(P_1)$, daher ist $P \cdot  P_1 \le G$. Wegen\footnote{$P \cdot P_1$ ist Untergruppe, da $P$ $P_1$ normalisiert} 
$\abs*{P \cdot P_1} = \frac{\abs*{P} \cdot \abs*{P_1}}{\abs*{P \cap P_1} }$ ist $P \cdot P_1$ eine $p$-Untergruppe von $G$. Wegen $P \le P \cdot P_1$ und $p \nmid m$ folgt
$P= P \cdot P_1 = P_1$. Daher ist $\abs*{T} = 1 \mod p $. 

Noch zu zeigen: $T$ enthält alle Sylowgruppen und jede $p$-Gruppe ist in einer $p$-Sylowgruppe enthalten. Sei $P_2 \le G$ eine $p$-Sylowgruppe mit $P_2 \not\in T$.
Dann operiert auch $P_2$ durch Konjugation auf $T$. Wenn $P_2$ auf $T$ einen Fixpunkt $P' \in T$ hat, dann ist wie eben $P_2 \cdot P'$ eine $p$-Untergruppe und dann 
$P_2 = P_2 \cdot P' = P' \in T$ \light. Daher hat $P_2$  auf $T$ keinen Fixpunkt. Dann folgt aber $p \mid \abs*{T} $ \light
Damit sind alle $p$-Sylowgruppen in $T$ enthalten, d.h. $\abs*{T} = n_p = 1 \mod p $. 

Ist $H \le G$ eine $p$-Untergruppe, dann operiert auch $H$ durch Konjugation auf $T$. Wegen $p \nmid \abs*{T} $ muss $H$ einen Fixpunkt $P' \in T$ besitzen, dann folgt
$H \cdot P' = P'$, d.h. $H \le P'$.

Weil $G$ durch Konjugation transitiv auf $T$ operiert, folgt 
\[
	n_p = \abs*{T} = [G : N_G(P)] \mid [G:P] = m. \bewende 
\]
\minisec{Bemerkung}
Wenn $G$ nur eine $p$-Sylowgruppe $P \le G$ besitzt, dann ist $P \unlhd G$. 
% subsection 18 (end)

\subsection{Satz (Frattini-Argument)} % (fold)
\label{sub:19}
Ist $G$ eine beliebige Gruppe, $H \unlhd G$ endlich und $P \le H$ eine $p$-Sylowgruppe von $H$. Dann ist $G =  N_G(P) \cdot H = H \cdot N_G(P)$, wobei
$N_G(P) = \set[g \in G]{P^g = g ^{-1} \cdot P \cdot g = P} $.
\minisec{Beweis}
Sei $g \in G$. Dann ist $P^g \le H^g = H$ eine $p$-Sylowgruppe von $H$. Daher existiert ein $h \in H$ mit $P^g = P^h$. Dann ist $P = P^{g \cdot h^{-1}}$, d.h. 
$g \cdot h^{-1} \in N_G(P)$. Damit ist 
\[
	g = \underbrace{g \cdot h ^{-1}}_{\in N_G(P)} \cdot \underbrace{h}_{\in H} \bewende
\]
\minisec{Bemerkung}
Sind $H_1, H_2 \le G$, $H_2 \le N_G(H_1)$, dann ist $H_1 \cdot H_2 = H_2 \cdot  H_1 \le G$.

Beweis:
Für alle $h_1\in H_1$ und $h_2\in H_2$ gilt:
\[
	h_1 \cdot h_2 =h_2 \cdot \underbrace{h_2^{-1} \cdot h_1 \cdot h_2}_{\in H_1} \quad\text{und}\quad h_2 \cdot h_1=\underbrace{h_2 \cdot h_1 \cdot h_2^{-1}}_{\in H_1}\cdot h_2. \bewende
\]
% subsection 19 (end)

\subsection[Bemerkung zu $p$-Sylowgruppen in Normalteilern und Faktorgruppen]{Bemerkung} % (fold)
\label{sub:110}
Offensichtlich gilt für eine endliche Gruppe $G$, $P \le G$ $p$-Sylowgruppe, $N \unlhd G$
\begin{multicols}{2}
	\begin{enumerate}[(i)]
		\item $P \cap N$ ist $p$-Sylowgruppe von $N$
		\item $\nicefrac{P \cdot N}{N}$ ist $p$-Sylowgruppe von $\nicefrac{G}{N}$.
	\end{enumerate}
\end{multicols}
\minisec{Beweis}
Es ist $[N : P \cap N] \stackrel{\text{2. Iso}}{=} [PN : P] $ teilerfremd zu $p$, und $P \cap N$ ist $p$-Untergruppe von $N$. Es ist 
$\nicefrac{\enbrace*{\nicefrac{G}{N}}}{\enbrace*{\nicefrac{P N}{N}}} \stackrel{\text{3. Iso}}{\simeq} \nicefrac{G}{PN}$, und $[G : PN] \mid [G :P] $ teilerfremd zu $p$. 
Wegen $[PN : N] = [P : (P \cap N)]  $ ist $\nicefrac{PN}{N}$ eine $p$-Gruppe. \bewende
% subsection 110 (end)

\subsection[Definition: Normalreihe und Kompositionsreihe]{Definition} % (fold)
\label{sub:111}
Eine Folge von Untergruppen $(H_i)_{0 \le i \le n}$ mit $H_0 = G$, $H_n = \set{1_G} $, $H_{i+1} \unlhd H_i$ heißt \Index{Normalreihe} in $G$. Ist $\nicefrac{H_i}{H_{i+1}}$
einfach\footnote{ Eine Gruppe $G$ heißt einfach, wenn sie $G$ und $\set{1_G}$ als einzige Normalteiler hat.} für alle $i < n$, 
dann heißt die Folge \Index{Kompositionsreihe}. Zwei Normalreihen $(H_i)_{i \le n}$, $(K_j)_{j \le m}$ heißen \Index{äquivalent}, falls $n=m$
und die auftretenden Quotienten $\enbrace*{\nicefrac{H_i}{H_{i+1}}}_{i \le n-1} $  nach geeigneter Permutation isomorph sind zu den Quotienten 
$\enbrace*{\nicefrac{K_j}{K_{j+1}}}_{j \le n-1} $.
% subsection 111 (end)

\subsection[Beispiel zu Normalreihen]{Beispiel} % (fold)
\label{sub:112}
\begin{align*}
	\mathds{Z}_6 \simeq \mathds{Z}_3 \times \mathds{Z}_2 &\vartriangleright \mathds{Z}_3 \vartriangleright \set{1} \\
	&\vartriangleright \mathds{Z}_2 \vartriangleright \set{1}  
\end{align*}
\minisec{Bemerkung}
\begin{enumerate}[(i)]
	\item Nicht jede Gruppe besitzt eine Kompositionsreihe, zB. hat $\mathds{Z}$ \emph{keine} Kompositionsreihe.
	\item Eine Normalreihe ist genau dann eine Kompositionsreihe, wenn es keine echte Verfeinerung gibt. Insbesondere hat also jede endliche Gruppe eine Kompositionsreihe.
\end{enumerate}
% subsection 112 (end)

\subsection{Ziel: Satz von Jordan-Hölder} % (fold)
\label{sub:113}
Sei $G$ eine Gruppe mit Kompositionsreihen $(H_i)_{i \le n}$ und $(K_j)_{j \le m}$. Dann sind die Reihen äquivalent.
\bigskip\\
Für den Beweis brauchen wir
% subsection 113 (end)

\subsection{Schmetterlings-Lemma (Zassenhaus)} % (fold)
\label{sub:114}
Sei $G$ eine Gruppe, $H,K \le G$ und $H' \unlhd H, K' \unlhd K$. Dann ist 
\[
	H'( H \cap K') \unlhd H'(H \cap K) \enspace\text{ und } \enspace K'( K \cap H') \unlhd K'(K \cap H)
\]
und die Quotienten sind isomorph.
\tikzexternaldisable
\[	
	\begin{tikzpicture}[baseline= (a).base]
		\node[scale=0.7] (a) at (0,0){
		\begin{tikzcd}
			& H \dar[dash]& & K \dar[dash]& \\
			& H'(H \cap K)\arrow[dd, dash] \drar[dash]& & K'(H \cap K)\arrow[dd, dash] \dlar[dash] \\
			& & H \cap K  \arrow[dd, dash]& & \\
			& H'(H \cap K') \drar[dash] \dlar[dash]& & K'(H' \cap K)\dlar[dash]  \drar[dash]\\
			H' \drar[dash]& & (H \cap K') ( H' \cap K) \drar[dash] \dlar[dash] & & K' \dlar[dash]\\
			& H' \cap K & & H \cap K' &
		\end{tikzcd}
		};
	\end{tikzpicture}
\]
\tikzexternalenable
\minisec{Beweis}
Setze $N := H \cap K$ und $M := H'(H \cap K')$. Dann gilt $N \le N_G(M)$ wegen $H'  \unlhd H, K' \unlhd K$ und daher $M \unlhd N \cdot M = H'(H \cap K)$.\\
\uline{Behauptung:} Es ist $N \cap M = (H \cap K) \cap \big( H' (H \cap K')\big) = (H' \cap K) ( H \cap K')$.
\begin{description}
	\item["$\subseteq $":] Sei $h'\cdot k \in H \cap K$ mit $h' \in H'$ und $k \in H \cap K'$ $\Rightarrow h'\cdot  k \in (H' \cap K) (H \cap K')$
	\item["$\supseteq $":] $h' \cdot k$ mit $h' \in H' \cap K$ und $k \in H \cap K'$, dann ist $h'\cdot  k \in H \cap K$.
\end{description}
Daher ist 
\[
	\nicefrac{NM}{M} = \nicefrac{H' (H  \cap K') (H \cap K)}{H'(H \cap K)} \cong \nicefrac{N}{N \cap M} \cong \nicefrac{(H \cap K)}{(H' \cap K) (H \cap K')}
\]
(so steht es in den Notizen). Besser finde ich:
\[
	\nicefrac{N M}{M} \stackrel{\text{2. Iso}}{\simeq} \nicefrac{N}{N \cap M} = \nicefrac{(H \cap K)}{(H' \cap K) (H \cap K')}
\]
Die rechte Seite ist symmetrisch in $H$ und $K$. Daher sind beide Quotienten im Lemma isomorph zu $\nicefrac{N}{N \cap M}$ und das Lemma ist bewiesen. 
\marginnote{Damit zeigen wir nun folgenden Satz:} \bewende
 
% subsection schmetterlings_lemma_zassenhaus_ (end)

\subsection{Satz von Schreier} % (fold)
\label{sub:115}
Sind $(H_i)_{i \le n}$, $(K_j)_{j \le m}$ Normalreihen in $G$, dann existieren äquivalente Verfeinerungen.
\minisec{Beweis}
Für $j= 1, \ldots , m-1$, $i=0, \ldots , n-1$ setze 
\[
	H'_{i m + j} := H_{i+1}( H_i \cap K_j)
\]
und  für
$i=0, \ldots ,n$ sei
\[
	H'_{i m} := H_i \qquad {\color{light_gray} =H_{i+1} (H_i \cap K_0) = H_i (H_{i-1} \cap K_m).}
\]
Für $i=1, \ldots , n-1$, $j=0, \ldots , m-1$ setze dementsprechend $K'_{j n +i} := K_{j+1}(K_j \cap H_i)$ und $K'_{j n} := K_j (= K_{j+1}(K_j \cap H_0) = K_j(K_{j-1} \cap H_n))$ für $j=0, \ldots , m$.

Nach dem Zassenhaus-Lemma (\ref{sub:114}) sind dann
\[
	\nicefrac{H'_{im +j}}{H'_{im+j+1}} \simeq \nicefrac{K'_{jn+i}}{K'_{jn + i +1}}
\]
und damit sind diese Verfeinerungen äquivalent. Damit folgt der Satz von Jordan-Hölder: Kompositionsreihen haben keine echten Verfeinerungen, müssen also bereits äquivalent 
sein! \marginnote{reviewed 22.4.14} \bewende
% subsection 115 (end)

\subsection[Definition: Auflösbare und nilpotente Gruppen]{Definition} % (fold)
\label{sub:116}
Eine Gruppe heißt \Index{auflösbar}, wenn sie eine abelsche Normalreihe besitzt, d.h. eine Normalreihe mit abelschen Quotienten. Eine Gruppe heißt \Index{nilpotent}, wenn 
es eine Normalreihe $(H_i)_{i \le n}$ gibt mit $H_i \unlhd G$ und $\nicefrac{H_i}{H_{i+1}} \le Z \enbrace*{\nicefrac{G}{H_{i+1}}}$.
% subsection 116 (end)

\subsection[Bemerkung: Nilpotente Gruppen sind auflösbar, Umkehrung gilt nicht]{Bemerkung} % (fold)
\label{sub:117}
Jede nilpotente Gruppe ist auflösbar, aber nicht umgekehrt: $S_3$ ist auflösbar $1 \unlhd \langle (1 2 3)\rangle \unlhd S_3$, aber $Z(S_3) = 1$, d.h. $S_3$ ist \emph{nicht}
nilpotent.
% subsection 117 (end)

\subsection[Satz: Auflösbarkeit von Untergruppen, Quotienten und Produkten auflösbarer Gruppen]{Satz} % (fold)
\label{sub:118}
Untergruppen und Quotienten auflösbarer Gruppen sind auflösbar, direkte Produkte auflösbarer Gruppen sind ebenfalls auflösbar.
\minisec{Beweis}
Ist $1 = G_0 \unlhd G_1 \unlhd \ldots \unlhd G_n  = G$ abelsche Normalreihe, $H \le G$, dann ist $1 = G_0 \cap H \unlhd G_1 \cap H \unlhd \ldots  \unlhd G_n \cap H = H$
abelsche Normalreihe in $H$, denn 
\[
	\nicefrac{(G_{i+1} \cap H)}{G_i \cap H} \simeq \nicefrac{G_i (G_{i+1} \cap H)}{G_i} \le \nicefrac{G_{i+1}}{G_i} \text{ ist abelsch.}
\]
Ist $N \unlhd G$, dann ist $(\nicefrac{G_i N}{N})$ abelsche Normalreihe für $\nicefrac{G}{N}$, denn es ist 
\[
	\nicefrac{(\nicefrac{G_{i+1}N}{N})}{(\nicefrac{G_i N}{N})} \simeq \nicefrac{G_{i+1}N}{G_i N} \simeq \nicefrac{G_{i+1}}{G_{i+1} \cap (G_i N)}
\]
ein Quotient von $\nicefrac{G_{i+1}}{G_i}$ und daher abelsch. (Da $G_i \le G_{i+1} \cap (G_i N)$, ist $\nicefrac{G_{i+1}}{G_i} \to \nicefrac{G_i+1}{ G_{i+1} \cap G_i N}$
ein Epimorphismus und daher ist die rechte Seite abelsch.) \bewende
% subsection 118 (end)

\subsection[Korollar: Auflösbarkeit ist äquivalent zur Auflösbarkeit von Normalteilern und Quotienten]{Korollar} % (fold)
\label{sub:119}
Sei $N  \unlhd G$. Dann ist $G$ auflösbar genau dann, wenn $N$ und $\nicefrac{G}{N}$ auflösbar sind. 
\minisec{Beweis}
\begin{description}
	\item["$\Rightarrow $":] \ref{sub:118}
	\item["$\Leftarrow$":] klar: Wir können die abelschen Normalreihen für $N$ und $\nicefrac{G}{N}$ zusammensetzen:
	\[
		1 = H_0 \unlhd H_1 \unlhd \ldots \unlhd H_k = N, \qquad \nicefrac{K_0}{N} = N \unlhd \nicefrac{K_1}{N} \unlhd \ldots \unlhd \nicefrac{K_m}{N} = \nicefrac{G}{N}
	\]
	Setze $1= H_0 \unlhd \ldots \unlhd H_k = K_0 \unlhd K_1 \unlhd \ldots \unlhd K_m = G$. Wegen $\nicefrac{(\nicefrac{K_{i+1}}{N})}{(\nicefrac{K_i}{N})} \simeq \nicefrac{K_{i+1}}{K_i}$. \bewende
\end{description}
% subsection 119 (end)

\subsection[Korollar: Das Produkt auflösbarer Normalteiler ist auflösbar]{Korollar} % (fold)
\label{sub:120}
Sind $M,N \unlhd G$ auflösbar, dann auch $M N$ auflösbar.
\minisec{Beweis}
$\nicefrac{M N}{N} \simeq \nicefrac{M}{M \cap N}$ ist auflösbar. Nach \ref{sub:119} ist $M N$ auflösbar. \bewende
% subsection 120 (end)

\subsection*{Einschub: Direktes Produkt} % (fold)
\label{sub:einschub_direktes_produkt}
Sind $G, H$ Gruppen, dann ist das direkte Produkt $G \times H$ die Gruppe mit Multiplikation
\[
	(g,h) \cdot (g', h') = (g \cdot g, h \cdot h')
\] 
% subsection einschub_direktes_produkt (end)

\subsection[Satz: Untergruppen und Quotienten nilpotenter Gruppen]{Satz} % (fold)
\label{sub:121}
Untergruppen und Quotienten nilpotenter Gruppen sind wieder nilpotent, die Produkte nilpotenter Gruppen sind nilpotent.
\minisec{Beweis}
Wie Satz \ref{sub:118}: 

Ist $1 = G_0 \unlhd G_1 \unlhd \ldots \unlhd G_n  = G$ Zentralreihe, $H \le G$, dann ist $1 = G_0 \cap H \unlhd G_1 \cap H \unlhd \ldots  \unlhd G_n \cap H = H$
Zentralreihe in $H$, denn 
\[
	\nicefrac{(G_{i+1} \cap H)}{G_i \cap H} \stackrel{\text{2. Iso}}{\simeq} \nicefrac{G_i (G_{i+1} \cap H)}{G_i} \le \nicefrac{G_{i+1}}{G_i} \text{ ist abelsch.}
\]
Ist $N \unlhd G$, dann ist $(\nicefrac{G_i N}{N})$ Zentralreihe für $\nicefrac{G}{N}$, denn es ist 
\[
	\nicefrac{(\nicefrac{G_{i+1}N}{N})}{(\nicefrac{G_i N}{N})} \simeq \nicefrac{G_{i+1 }N}{G_i N} \simeq \nicefrac{G_{i+1}}{G_{i+1} \cap (G_i N)}
\]
ein Quotient von $\nicefrac{G_{i+1}}{G_i}$ und daher abelsch. (Da $G_i \le G_{i+1} \cap (G_i N)$, ist $\nicefrac{G_{i+1}}{G_i} \to \nicefrac{G_i+1}{ G_{i+1} \cap G_i N}$
ein Epimorphismus und daher ist die rechte Seite abelsch.) \bewende
% subsection 121 (end)

\subsection[Satz: Endliche $p$-Gruppen sind nilpotent]{Satz} % (fold)
\label{sub:122}
Endliche $p$-Gruppen sind nilpotent.
\minisec{Beweis}
Nach Satz \ref{sub:16} ist $H_1 := Z(G) \not= 1$. Da $\nicefrac{G}{Z(G)}$ wieder $p$-Gruppe ist, ist $Z(\nicefrac{G}{Z(G)}) \not= 1$. Setze 
\[
	H_2 := \pi_{Z(G)} ^{-1} \enbrace[\big]{Z (\nicefrac{G}{Z(G)})}  \enspace \text{ usw.}
\]
Nach endlich vielen Schritten ist $H_k=G$. Es gilt dann 
\[
	\nicefrac{H_{i+1}}{H_i} = Z(\nicefrac{G}{H_i})
\]
d.h. die $H_i$ bilden die \Index{obere Zentralreihe}. 
% subsection 122 (end)

\subsection[Definition: Kommutator]{Definition} % (fold)
\label{sub:123}
Für $a,b \in G$ heißt $[a,b] = a ^{-1} \cdot b ^{-1} \cdot a \cdot b$ der \Index{Kommutator} von $a$ und $b$.
\begin{enumerate}[(i)]
	\item\label{123:1} Es ist $a \cdot b = b \cdot a \cdot [a,b]$ und $[a,b]= 1$, genau dann wenn $a \cdot b=b \cdot a$.
	\item\label{123:2} Ist $\varphi : G \to H$, dann ist $\varphi([a,b])= [\varphi(a), \varphi(b)]$.
	\item\label{123:3} Produkte von Kommutatoren sind nicht unbedingt selber wieder ein Kommutator!
\end{enumerate} 
Für Untergruppen $H, K \le G$ setze $[K,H] := \big\langle [k,h] \vert k \in K, h \in H \big\rangle$. Ist $K \le N_G(H)$, dann ist $[K,H] \le H$, denn 
$k ^{-1} \cdot h ^{-1} \cdot k \cdot h  = \enbrace*{h ^{-1}}^k \cdot h\in H$. Die Gruppe $G' = [G,G] = \big\langle [g,h] \vert g,h \in G \big\rangle$
heißt \Index{Kommutatorgruppe} von $G$.
% subsection 123 (end)

\subsection[Satz: Eigenschaften der Kommutatorgruppe]{Satz} % (fold)
\label{sub:124}
\begin{enumerate}[(i)]
	\item $G' \unlhd G$
	\item $\nicefrac{G}{G'}$ ist abelsch.
	\item Ist $\varphi : G \to A$ ein Gruppenhomomorphismus und $A$ abelsch, dann ist $G' \le \ker \varphi$.
\end{enumerate}
\minisec{Beweis}
\begin{enumerate}[(i)]
	\item Es ist $g ^{-1} [a,b] g = [a^g, b^g]$ nach \ref{sub:123} (\ref{123:2}).
	\item Klar nach \ref{sub:123} (\ref{123:1}).
	\item Es ist $\varphi([a,b]) = [\varphi(a), \varphi(b)] = 1$, d.h. $G' \le \ker \varphi$. \bewende
\end{enumerate}
\minisec{Bemerkung}
Mit anderen Worten: $G'$ ist der kleinste Normalteiler von $G$ mit $\nicefrac{G}{G'}$ abelsch, denn ist $\nicefrac{G}{N}$ abelsch, dann ist nach (iii) mit 
$\varphi : G \to \nicefrac{G}{N}$, $G' \le \ker \varphi = N$
% subsection 124 (end)

\subsection[Definition: Konstruktion weiterer Kommutatorgruppen]{Definition} % (fold)
\label{sub:125}
Wir setzen $G^{(0)} = G$, $G^{(i+1)} = \left[G^{(i)}, G^{(i)}\right]$. Dann ist $G^{(1)}= G'$ und $G^{(i+1)} \unlhd G^{(i)}$, $\nicefrac{G^{(i)}}{G^{(i+1)}}$ abelsch.
% subsection 125 (end)

\subsection[Satz: Auflösbarkeit einer Gruppe $G$]{Satz} % (fold)
\label{sub:126}
$G$ ist auflösbar genau dann, wenn $G^{(k)}= 1_G$ für ein $k \ge 0$.
\minisec{Beweis}
\begin{description}
	\item["$\Leftarrow$":] Die $G^{(i)}$ bilden eine abelsche Normalreihe.
	\item["$\Rightarrow $":] Ist $(N_i)_{i \le n}$, $N_0 = G$, $N_n = \set{1}$, dann ist mit Induktion $G^{(i)} \le N_i$ nach voriger Bemerkung, also 
	$G^{(n)} \le \set{1} = N_n$ 
\end{description}
\minisec{Bemerkung}
Damit ist $(G^{(i)})_{i \le k}$ die am schnellsten absteigende untere Normalreihe für $G$. $k$ heißt \Index{auflösbare Länge} von $G$.
% subsection 126 (end)

\subsection[Definition: Untere Zentralreihe]{Definition} % (fold)
\label{sub:127}
Die \Index{untere Zentralreihe} einer nilpotenten Gruppe $G$ ist definiert durch $G^{[0]} = G, \, G^{[i]} = \left[ G^{[i-1]}, G\right]$.
Es ist $\nicefrac{G^{[i]}}{G^{[i+1]}} \le Z(\nicefrac{G}{G^{[i+1]}})$ nach Definition.
% subsection 127 (end)

\subsection[Satz: Charakterisierung von Nilpotenz über die untere Zentralreihe]{Satz} % (fold)
\label{sub:128}
Eine Gruppe $G$ ist nilpotent genau dann, wenn $G^{[k]} = 1_G$ für ein $k \ge 0$.
\minisec{Beweis}
\begin{description}
	\item["$\Leftarrow $":] Klar nach voriger Bemerkung: $(G^{[i]})_{i \le k}$ bilden Zentralreihe. 
	\item["$\Rightarrow$":] Ist $1 = N_0 \le N_1 \ldots  \le N_n  = G$ eine Zentralreihe, dann ist  $G^{[1]} \le N_{n-1}$, denn $\nicefrac{G}{N_{n-1}}$ ist abelsch. Zeige
	$G^{[i]} \le N_{n-i}$ für $i=1, \ldots ,n$, denn dann folgt $G^{[n]}=1$. Weil $\nicefrac{N_{n-i}}{N_{n-(i+1)}} \le Z \enbrace*{\nicefrac{G}{N_{n-(i+1)}}}$ folgt 
	$[N_{n-i},G] \le N_{n-(i+1)}$. Nach Induktion ist wegen $G^{[i]} \le N_{n-i}$ dann 
	\[
		G^{[i+1]} = [G^{[i]}, G] \le [N_{n-i}, G] \le N_{n-(i+1)}
	\]
\end{description}
% subsection 128 (end)

\subsection{Einschub über direkte und semidirekte Produkte} % (fold)
\label{sub:129}
\begin{enumerate}[a)]
	\item Sei $G$ eine Gruppe, $H \le G$, $N \unlhd G$ ein Normalteiler mit $H \cap N = \set{1}$ und $N \cdot H = G$. Dann ist die Abbildung $\varphi : N \times H \to G$,
	$(n,h) \mapsto n \cdot h$ bijektiv, d.h. für jedes $g \in G$ existiert ein eindeutig bestimmtes $n \in N, h \in H$ mit $n \cdot h = g$.
	Denn ist 
	\[
		n_1 \cdot h_1 = n_2 \cdot h_2 \iff \underbrace{n_2 ^{-1} n_1}_{\in N} = \underbrace{h_2 \cdot h_1 ^{-1}}_{\in H} \in N \cap H = 1
	\]
	Aber: Im Allgemeinen ist $\varphi$ \emph{kein} Gruppenhomomorphismus, denn es ist 
	\[
		(n_1, h_1) (n_2 h_2) = \enbrace{n_1 \underbrace{h_1 n_2 h_1 ^{-1}}_{\in N}} \underbrace{(h_1 h_2)}_{\in H} = (n_1 n_2) (n_2 ^{-1} h_1 n_2 h_2)
	\]
	Daher ist $\varphi$ ein Gruppenhomomorphismus genau dann, wenn $H$ die Elemente aus $N$ zentralisiert ($n \cdot h = h \cdot n$),
	 d.h. wenn $H\unlhd G$. In dem Fall ist dann $G\simeq N \times H$. 
	
	Ist $H \unlhd G$, dann gilt $\varphi(n_1, h_1) \cdot (n_2, h_2) = (n_1 \cdot n_2, h_1 \cdot h_2)$
	\item Sind $H,N, G$, $\varphi : H \to \Aut(N)$ ein Homomorphismus, dann definiere eine Verknüpfung auf der Menge $G = N \times H$ durch
	\[
		(n_1, h_1) \cdot (h_2, h_2) = (n_1 \cdot \varphi(h_1) (n_2), h_1 \cdot h_2)
	\]
	Mit dieser Verknüpfung wird $G$ zu einer Gruppe mit Untergruppen $\set{1} \times H \simeq H $, $N \times \set{1} \simeq N $. Man schreibt 
	$G = N \rtimes H = N \rtimes_\varphi H$ für das \bet{semidirekte Produkt}\index{semidirektes Produkt}. Mit $\set{1} \times H $ und $N \times \set{1} $ können wir $G$ 
	wie in a) beschreiben. Dabei ist $\varphi : H \to \Aut(N)$, $h \mapsto \kappa_h$ (Konjugation mit $h$).
\end{enumerate}
% subsection 129 (end)
% section 1 (end)
\newpage

\section{Moduln: Halbeinfache Moduln, freie Moduln} % (fold)
\label{sec:2}

\minisec{Erinnerung}
In Algebra \RM{1} wurden hauptsächlich kommutative Ringe betrachtet:
\begin{enumerate}[a)]
	\item Körper, Polynomringe, $\mathds{Z}$
	\item nicht kommutative Ringe: $R$ Ring, Matrizenring $M_n(R)= R^{n \times n}$. 
	
	Sei $A$ eine abelsche Gruppe, dann ist $\End(A) = \Hom(A,A)$ ein Ring 
	\begin{align*}
		(\varphi + \psi)(x) &= \varphi(x) + \psi(x) \\
		(\varphi \cdot \psi)(x) &= \varphi(\psi(x))
	\end{align*}
	$\End(A)$ heißt der \Index{Endomorphismenring} von $A$. Dies ist das allgemeinste Beispiel, denn es gilt:
\end{enumerate}

\subsection[Satz: Jeder Ring ist isomorph zu einem Endomorphismenring]{Satz} % (fold)
\label{sub:21}
Jeder Ring $R$ ist isomorph zu einem Ring von Endomorphismen einer abelschen Gruppe.
\minisec{Beweis}
Ist $(R,+, \cdot )$ gegeben, dann ist $A=(R,+)$ eine abelsche Gruppe. Die Abbildung $R \to \End(A), a \mapsto \lambda_a $ mit 
$\lambda_a : A \to A, x \mapsto a \cdot x$ ist ein injektiver Ringhomomorphismus, eingeschränkt auf das Bild also ein Isomorphismus. \bewende
% subsection 21 (end)

\subsection[Definition: Modul]{Definition} % (fold)
\label{sub:22}
Sei $R$ ein Ring, $(M,+)$ eine abelsche Gruppe. Eine $R$-(Links-)\Index{Modulstruktur} auf $M$ ist eine Verknüpfung $R \times M \to M, (r,m) \mapsto r \cdot m$ mit
\begin{enumerate}[(i)]
	\item $r(x+y)= r \cdot x + r \cdot y$
	\item $(r+s) \cdot x = r \cdot x + s \cdot x$
	\item $(r \cdot s) x = r \cdot (s \cdot x)$
	\item $1_R \cdot x = x$
\end{enumerate}
für alle $r,s \in R, \,x,y \in M$. Ist $R$ ein Körper, dann sind die $R$-Moduln genau die $R$-Vektorräume. Mit anderen Worten: Eine $R$-Modulstruktur auf $M$ ist 
(gegeben durch) einen Ringhomomorphismus $\varphi : R \to \End(M,+)$ mit $r \cdot x = \varphi(r) (x)$.
% subsection 22 (end)

\subsection[Beispiele für Moduln]{Beispiele} % (fold)
\label{sub:23}
\begin{enumerate}[(i)]
	\item Ist $R$ ein Körper, dann ist ein $R$-Modul ein $R$-Vektorraum.
	\item $(R,+)$ ist $R$-Modul durch Produktwirkung, d.h. $\varphi : R \to \End(R),\, r \mapsto \lambda_r$.
	\item Ist $I \unlhd R$ ein \Index{Ideal} (d.h. für alle $i,j \in I$, $r \in R$ ist $i + j, i \cdot r, r \cdot i \in I$), dann ist auch $(I,+)$ ein
	$R$-Modul, ein $R$-Untermodul von $(R,+)$.
	\item Jede abelsche Gruppe ist ein $\mathds{Z}$-Modul. \hfill {\color{light_gray}(mittels $\varphi(z)(x) = x^z$, für $z \in \mathds{Z}$)}
\end{enumerate}
% subsection 23 (end)

\subsection[Definition: Untermodul, einfache Moduln und Ringe]{Definition} % (fold)
\label{sub:24}
Ist $M$ ein $R$-Modul, $N \le M$ Untergruppe mit $r \cdot x \in N$ für alle $x \in N$, $r \in R$, dann heißt $N$ ein $R$-\Index{Untermodul} von $M$.
\marginnote{$\set{0}$, $M$ sind immer Untermoduln} \smallskip \\
Ein Modul $M \not= \set{0}$ heißt \bet{einfach}\index{Modul!einfach} (oder \bet{irreduzibel}\index{Modul!irreduzibel}), wenn $\set{0},M$ die einzigen Untermoduln sind. 
Ein Ring $R$ heißt (links-)\bet{einfach}\index{Ring!einfach}, wenn er als (Links-)$R$-Modul einfach ist.  
\minisec{Bemerkung}
Einfache kommutative Ringe sind genau die Körper. Jedes Ideal in $R$ ist Untermodul, aber nicht jedes Untermodul ist ein Ideal.
% subsection 24 (end)

\subsection[Definition: erzeugte Untermoduln]{Definition} % (fold)
\label{sub:25}
\begin{itemize}
	\item Ist $\set{N_\alpha}_{\alpha \in I} $ Menge von Untermoduln von $M$, dann ist $\bigcap_{\alpha \in J} N_\alpha$ ein Untermodul.
	\item Ist $\emptyset \not= S \subseteq M$, dann ist $\langle S \rangle = \bigcap_{N \supseteq S} N$ der von $S$ \Index{erzeugte Untermodul}. Der von einer Summe 
	erzeugte Modul ist gegeben durch
	\[
		\sum_{\alpha \in I} N_\alpha = \langle n_{\alpha_1} + \ldots + n_{\alpha_k} : \alpha_i \in I, n_{\alpha_i} \in N_{\alpha_i}\rangle
	\] 
	Ist $S$ endlich, dann heißt $\langle S \rangle$ \Index{endlich erzeugt}. Ist $\abs*{S}=1 $, dann heißt $\langle S \rangle = M$ \Index{zyklisch}.
	\item Ein einfacher Modul ist zyklisch, aber nicht umgekehrt (zB. $\mathds{Z}$).   
\end{itemize}
% subsection 25 (end)

\subsection[Bemerkung zu Modulstrukturen auf Quotienten]{Bemerkung} % (fold)
\label{sub:26}
\begin{enumerate}[(i)]
	\item Ist $M$ ein zyklischer $R$-Modul, dann ist $M \simeq \nicefrac{R}{I}$ für ein Ideal $I \unlhd R$. \hfill(siehe Blatt 3 oder \ref{sub:anh_2})
	\item Ist $N \le M$ ein $R$-Untermodul, dann ist auch $\nicefrac{M}{N}$ ein $R$-Modul durch 
	\[
		r(m + N) = r \cdot m + N
	\]
\end{enumerate}
% subsection 26 (end)

\subsection[Definition: Klasse der $R$-Moduln, $R$-Modul-Homomorphismen]{Definition} % (fold)
\label{sub:27}
Die Klasse aller $R$-Links-Moduln bezeichnen wir mit $_R \Mod$. Sind $M,N \in {_R \Mod}$ und $\varphi : (M,+)\to (N,+)$ ein Homomorphismus (der additiven Gruppen), 
dann ist\marginnote{$\varphi$ ist $R$-linear}
$\varphi$ ein \Index{$R$-Modul-Homomorphismus}, falls
\[
	\varphi(r \cdot m) = r \cdot \varphi(m)\qquad  {\color{light_gray} \varphi \enbrace[\big]{\lambda_r(m)}= \lambda_r \enbrace[\big]{\varphi(m)}  .}
\]
% subsection 27 (end)

\subsection[Bemerkung zu $\Hom_R(M,N)$]{Bemerkung} % (fold)
\label{sub:28}
Kerne und Bilder von $R$-Modul-Homomorphismen sind $R$-Untermoduln. Die Menge $\Hom_R(M,N) := \set[\varphi: M \to N]{\varphi \text{ ist $R$-Modul-Homomorphismus}} $ ist 
eine abelsche Gruppe mit 
\[
	(\psi+ \varphi)(m) = \psi(m)+ \varphi(m)
\]
und $\End_R(M) := \Hom_R(M,M)$ ist mit $(\varphi \cdot \psi)(m) = \varphi\big(\psi(m)\big)$ der Endomorphismenring von $M$.
Die Homomorphie- und Isomorphiesätze für Gruppen gelten auch für Moduln:
% subsection 28 (end)

\subsection[Satz: Homomorphie- und Isomorphiesätze für Moduln]{Satz (Isomorphiesätze)} % (fold)
\label{sub:29}
\begin{enumerate}[(i)]
	\item Ist $f : M \to N$ ein $R$-Modul-Homomorphismus, $M' \subseteq M$ Untermodul mit $M'  \subseteq \ker f$, dann existiert ein eindeutiger $R$-Modul-Homomorphismus 
	$f'  : \nicefrac{M}{M'} \to N$ mit 
	\[
		\begin{tikzcd}
			M \rar{f} \dar{\pi }& N \\
			\nicefrac{M}{M'} \urar[dashed, swap]{f'}& 
		\end{tikzcd}
	\]
	und $f'$ ist injektiv genau dann, wenn $M' = \ker f$.
	\item Sind $A,B \subseteq M$ Untermoduln, dann gilt \marginnote{Vgl. Kapitel 1 \hyperref[sub:wiederholung_iso]{Seite \pageref{sub:wiederholung_iso}}}
	\[
		\nicefrac{(A+B)}{B} \simeq \nicefrac{A}{A \cap B}
	\]
	\item Ist $M' \subseteq M$ ein Untermodul, dann existiert ein \Index{Verbandsisomorphismus} zwischen den Untermoduln von $M$, die $M'$ enthalten, und den Untermoduln
	von $\nicefrac{M}{M'}$, nämlich $N \mapsto \nicefrac{N}{M'}$ und es gilt\hfill(vgl. \ref{sub:12})
	\[
		\nicefrac{\enbrace*{\nicefrac{M}{M'}} }{\enbrace*{\nicefrac{N}{M'}} } \simeq \nicefrac{M}{N}
	\] 
\end{enumerate}
\minisec{Beweis für (i)}
Es ist nur nachzurechnen, dass der (einzige mögliche) Gruppenhomomorphismus $f' : \nicefrac{M}{M'} \to N$, $m + M' \mapsto f(m)$ $R$-linear ist. Das folgt sofort:
\begin{align*}
	f' \enbrace[\big]{r(m + M')} = f' (r \cdot m + M')= f(r \cdot m) = r \cdot f(m) = r \cdot f'(m+ M')
\end{align*} 
Die Wohldefiniertheit sowie (ii) und (iii) sind einfache Übungsaufgaben. \bewende
% subsection 29 (end)

\subsection[Definition: Exakte Sequenz]{Definition (Sprechweise)} % (fold)
\label{sub:210}
Eine Folge von $R$-Moduln $(M_i)_i$ und Homomorphismen $f_i : M_{i-1} \to M_{i}$ heißt \bet{exakt in $M_i$}\index{exakte Folge}, falls $\im f_{i} = \ker{f_{i+1}}$.
Eine \Index{exakte Sequenz} ist eine Folge, die überall exakt ist. Eine exakte Sequenz von der Form 
\[
	0 \to M_1 \xrightarrow{f_1} M_2  \xrightarrow{f_2} M_3 \to 0  
\]
heißt \Index{kurze exakte Sequenz}. Dieses bedeutet $f_1$ injektiv, $f_2$ surjektiv und daher ist dann $M_3 \simeq \nicefrac{M_2}{M_1}$.
% subsection 210 (end)

\subsection[Definition: noethersch und artinsch]{Definition} % (fold)
\label{sub:211}
Ein (Links)$R$-Modul $M$ heißt (links-)\Index{noethersch}\footnote{nach Emmy Noether, 1882-1935, siehe \url{http://de.wikipedia.org/wiki/Emmy_Noether}}, 
wenn es keine unendliche echt aufsteigende Kette von Untermoduln gibt.
$M$ heißt (links-)\Index{artinsch}\footnote{nach Emil Artin, 1898-1962, siehe \url{http://de.wikipedia.org/wiki/Emil_Artin}}, 
wenn es keine unendliche echt absteigende Kette von Untermoduln gibt.\\
Ein Ring $R$ heißt \bet{noethersch}\index{Ring!noethersch} (bzw. \bet{artinsch}\index{Ring!artinsch}), wenn er als $R$-Modul noethersch (bzw. artinsch) ist.
\minisec{Beispiel}
$\mathds{Z}$ ist noethersch aber nicht artinsch. Allgemein gilt: HIR sind noethersch. Körper sind artinsch und noethersch.
% subsection 211 (end)

\subsection[Proposition: noethersch $\iff$ alle Untermoduln endlich erzeugt]{Proposition} % (fold)
\label{sub:212}
Ein $R$-Modul $M$ ist noethersch genau dann, wenn alle Untermoduln endlich erzeugt sind.
\minisec{Beweis}
\begin{description}
	\item["$\Rightarrow $":] Sei $N \subseteq M$, wähle induktiv $x_1, x_2, \ldots \in N$ mit $x_i \not\in \langle x_1, \ldots , x_{i-1} \rangle =: N_{i-1}$. Dann ist
	$(N_i)$ eine echt aufsteigende Kette und muss --- da $M$ noethersch ist --- nach endlich vielen Schritten mit $\langle x_1, \ldots , x_k \rangle =N$ enden. 
	Das heißt $N$ ist endlich erzeugt.
	\item["$\Leftarrow$":] Sei $N_0 \subseteq N_1 \subseteq \ldots $ eine echt aufsteigende Kette von Untermoduln in $M$ und $N :=\sum N_i$. 
	Da $N$ endlich erzeugt ist, existieren $x_1,\ldots ,x_r\in N$ mit $N=\langle x_1, \ldots , x_r\rangle$. Dann existiert ein $k$ mit $x_1,\ldots , x_r\in N_k$; 
	d.h. $N=N_k$ und die Kette ist endlich.\bewende
\end{description}
\minisec{Bemerkung}
Offensichtlich gilt: Ist ein $R$-Modul $M$ noethersch (bzw. artinsch) und $N \subseteq M$ ein Untermodul. Dann sind auch $N$ und $\nicefrac{M}{N}$ noethersch 
(bzw. artinsch). Dies gilt nach den Isomorphiesätzen. 
% subsection 212 (end)

\subsection[Satz: noethersch (artinsch) innerhalb einer kurzen exakten Sequenz]{Satz} % (fold)
\label{sub:213}
Ist $0 \to M' \xrightarrow{f} M \xrightarrow{g} M'' \to 0  $ eine kurze exakte Sequenz von $R$-Moduln, dann gilt: $M$ noethersch (bzw. artinsch) genau dann, wenn
$M'$ und $M''$ noethersch (bzw. artinsch).
\minisec{Beweis}
\begin{description}
	\item["$\Rightarrow$":] klar nach voriger Bemerkung. \hfill {\color{light_gray}($f$ injektiv, $M'' \simeq \nicefrac{M}{M'}$)}
	\item["$\Leftarrow$":] (artinsch) Sei $P_0 \supseteq P_1 \supseteq P_2 \supseteq \ldots $ eine echt absteigende Kette in $M$. Dann betrachte 
	\[
		P_0 \cap M' \supseteq P_1 \cap M' \supseteq \ldots 
	\]
	in $M'$ und 
	\[
		\nicefrac{\enbrace*{P_0 +  M'} }{M'} \supseteq \nicefrac{\enbrace*{P_1 + M'} }{M'} \supseteq \ldots 
	\]
	in $\nicefrac{M}{M'}$. Nach Vorraussetzung existiert $k \in \mathds{N}$ mit $P_k \cap M' = P_l \cap M'$ für $l \ge k$ und 
	$\nicefrac{\enbrace*{P_k +M'} }{M'} = \nicefrac{\enbrace*{P_l + M'} }{M'}$ für $l \ge k$. Wegen
	\[
		\nicefrac{\enbrace*{P_l + M'} }{M'} \simeq \nicefrac{P_l}{\enbrace*{P_l \cap M'} }\marginnote{Isomorphiesätze}
	\]
	folgt $P_l = P_k$ für $l \ge k$. Entsprechend für aufsteigende Ketten. \bewende
\end{description}
% subsection 213 (end)

\subsection[Korollar: Endliche Summen noetherscher Moduln sind noethersch (artinsch)]{Korollar} % (fold)
\label{sub:214}
Endliche Summen von noetherschen (bzw. artinschen) Moduln sind wieder noethersch (bzw. artinsch).
\minisec{Beweis}
Ist $M = N + P$ und $N,P$ noethersch, dann betrachte $0 \to N \hookrightarrow M \twoheadrightarrow \nicefrac{M}{N} \to 0$. Wegen 
$\nicefrac{M}{N} = \nicefrac{(N+P)}{N} \simeq \nicefrac{P}{(N \cap P)}$ ist $\nicefrac{M}{N}$ noethersch, also ist nach Satz \ref{sub:213} auch $M$ noethersch. 
Entsprechend für artinsch. \bewende
% subsection 214 (end)

\subsection[Korollar: Moduln über einen noetherschen Ring sind noethersch]{Korollar} % (fold)
\label{sub:215}
Ist $R$ ein noetherscher (bzw. artinschen) Ring, dann ist jeder endlich erzeugte $R$-Modul noethersch (bzw. artinsch).
\minisec{Beweis}
Durch Induktion über die Anzahl der Erzeuger. Ist $M$ zyklisch, dann ist $M \simeq \nicefrac{R}{J}$ (siehe \ref{sub:26})
und $\nicefrac{R}{J}$ noethersch (bzw. artinsch) nach Satz \ref{sub:213}. Sei nun 
\[
	M = \langle x_1, \ldots , x_n \rangle \enspace, \enspace M' =  \langle x_1, \ldots , x_{n-1}\rangle
\]
Nach Induktionsvorraussetzung ist $M'$ noethersch (artinsch) und $\nicefrac{M}{M'}$ ist zyklisch, daher auch noethersch (artinsch), nach \ref{sub:213} ist $M$ noethersch 
(artinsch). \bewende
% subsection 215 (end)

\subsection[Korollar: Endlich erzeugte Moduln über einem Hauptidealring sind noethersch]{Korollar} % (fold)
\label{sub:216}
Ist $R$ ein Hauptidealring, dann ist jeder endlich erzeugte $R$-Modul noethersch.
% subsection 216 (end)

\subsection[Definition: Halbeinfacher Modul]{Definition} % (fold)
\label{sub:217}
Ein $R$-Modul $M$ heißt \bet{halbeinfach}\index{Modul!halbeinfach} (oder vollständig zerlegbar), wenn jeder Untermodul $N$ ein Komplement hat, d.h. wenn $N' \subseteq M$ existiert mit $M = N \oplus N'$, d.h. $N \cap N' = \set{0}, N+ N' = M $.
% subsection 217 (end)

\subsection[Beispiele zu halbeinfachen Moduln]{Beispiele} % (fold)
\label{sub:218}
\begin{enumerate}[(i)]
	\item $\mathds{Z}$ (als $\mathds{Z}$-Modul) ist \emph{nicht} halbeinfach. $m \mathds{Z} \subseteq \mathds{Z}, m \not= 0$ hat kein Komplement, denn für 
	$k \cdot \mathds{Z}$ gilt immer $k \cdot \mathds{Z} \cap m \cdot \mathds{Z} \ni k \cdot m \not= 0$.
	\item Ist $R$ ein Körper, dann sind alle $R$-Vektorräume halbeinfach nach dem Basisergänzungssatz.
	\item Untermoduln und Quotienten halbeinfacher Moduln sind halbeinfach.
	\item Einfache Moduln sind halbeinfach.
	\item $\mathds{Z}_6 \simeq \mathds{Z}_2 \times \mathds{Z}_3$ ist halbeinfach, weil $\mathds{Z}_2, \mathds{Z}_3$ die einzigen nicht-trivialen Untermoduln sind.
\end{enumerate}
Ein Ring $R$ heißt (links-)halbeinfach, wenn er als $R$-Modul halbeinfach ist.
% subsection 218 (end)

\subsection[Satz: Äquivalenzen zu halbeinfach durch Summen aus einfachen Untermoduln]{Satz} % (fold)
\label{sub:219}
Für einen $R$-Modul $M$ sind äquivalent:
\begin{enumerate}[(i)]
	\item $M$ ist halbeinfach.
	\item $M$ ist Summe von einfachen Moduln, d.h. es existiert Familie $(N_\alpha)_{\alpha \in I}$ von einfachen Untermoduln, die $M$ erzeugen, also 
	$M = \sum_{\alpha \in I} N_\alpha$.
	\item $M$ ist direkte Summe von einfachen Untermoduln, d.h. es existiert eine Familie $(N'_\alpha)_{\alpha \in I'}$ von einfachen Untermoduln mit 
	\[
		M = \bigoplus_{\alpha \in I'} N'_\alpha
	\]
	d.h. $N'_\alpha \cap \sum_{\beta \not= \alpha} N'_\beta = 0$
\end{enumerate}
\minisec{Beweis (mit Lemma \ref{sub:220})}
\begin{description}
	\item["(i)$\Rightarrow$(ii)":] Sei $\set{N_\alpha}_{\alpha\in I}$ die Menge \emph{aller} einfachen Untermoduln von $M$. Diese ist nach Lemma \ref{sub:220} nicht leer. 
	Setze $M_1 := \sum N_\alpha$.
	
	\uline{Behauptung:} $M_1 = M$. Sonst existiert ein Komplement $P \subseteq M$ mit $P \not= 0$, $P \oplus M_1 = M$. Dann ist $P$ halbeinfach und enthält daher einen 
	einfachen Untermodul $N \subseteq P$ \light zu $N \in \set{N_\alpha}_{\alpha \in I}$.
	\item["(iii)$\Rightarrow $(ii)":] Klar.
	\item["(ii)$\Rightarrow $(i)":] Sei $P \subseteq M$ Untermodul. Betrachte die Menge aller $J \subseteq I$ mit 
	\begin{enumerate}[(a)]
		\item $N_i \cap \sum_{j \not= i, j \in J} N_j = 0$ für alle $i \in J$
		\item $P \cap \sum_j N_j = 0$
	\end{enumerate}
	Weil $J = \emptyset$ die Bedingungen (a) und (b) erfüllt, können wir Zorns Lemma anwenden und finden eine maximale Teilmenge $J \subseteq I$ mit (a) und (b).
	
	\uline{Behauptung:} 
	\[
		M_1 := P \oplus \sum_{j \in J} N_j = P \oplus \bigoplus_{j \in J} N_j = M
	\]
	Für $\alpha \in I$ ist $N_\alpha \cap M_1 \in\set{0, N_\alpha}$, da $N_\alpha$ einfach ist. Ist $N_\alpha \cap M_1 = 0$, dann erfüllt $J \cup \set{\alpha}$ die Bedingungen
	(a) und (b). \light Maximalität von $J$. Daher ist $N_\alpha \subseteq M_1$, also $M_1 = M$.
	\item["(ii)$\Rightarrow $(iii)":] folgt aus dem Beweis "(ii) $\Rightarrow$ (i)"{} mit $P=0$. \bewende
\end{description}
% subsection 219 (end)

\subsection[Lemma: Jeder halbeinfache Modul hat einen einfachen Untermodul]{Lemma} % (fold)
\label{sub:220}
Ist $M \not= 0$ halbeinfach, dann hat $M$ einen einfachen Untermodul.
\minisec{Beweis}
Sei $m \in M, m \not= 0$. Betrachte $N := \langle m \rangle \subseteq M$. Nach Zorns Lemma existiert ein maximaler Untermodul $P \le N$ mit $m \not\in P$ (denn $0$ 
ist ein solcher Untermodul). Sei $Q$ ein Komplement von $P$ in $N$, also $P \oplus Q = N$, $Q \not= 0$, da $m \not\in P$. \\
Behauptung: $Q$ ist einfach.  Beweis: $Q \subseteq N \subseteq M$. Ist $0 \not= Q' \subseteq Q$ ein Untermodul, dann ist $Q' \oplus P  \supsetneq P$, also gilt wegen der
Maximalität von $P$: $m \in Q' \oplus P$,
also $Q' \oplus P = N$ und daher $Q=Q'$. \bewende
% subsection 220 (end)

\subsection[Satz: Alle $R$-Moduln halbeinfach $\Leftrightarrow$ $(R,+)$ halbeinfach]{Satz} % (fold)
\label{sub:221}
Für einen Ring $R$ sind äquivalent:
\begin{enumerate}[(i)]
	\item Alle $R$-Moduln sind halbeinfach.
	\item Alle endlich erzeugten $R$-Moduln sind halbeinfach.
	\item Alle zyklischen $R$-Moduln sind halbeinfach.
	\item $(R,+)$ ist als $R$-Modul halbeinfach.
\end{enumerate}
\minisec{Beweis}\leavevmode
\begin{description}
	\item[\enquote{(i)$\Rightarrow$(ii)$\Rightarrow$(iii)$\Rightarrow$(iv)}:] Klar.
	\item[\enquote{(iv)$\Rightarrow $(iii)}:] Jeder zyklische $R$-Modul ist von der Form $\nicefrac{R}{I}$ und Quotienten halbeinfacher Moduln sind halbeinfach. \hfill 
	{\color{light_gray} siehe \ref{sub:26} bzw. Blatt 3}
	\item[\enquote{(iii)$\Rightarrow$(i)}:] Sei $M \in {_R \Mod}$, dann ist $M = \sum_{m \in M} R m$ Summe zyklischer Moduln. Da jeder zyklische $R$-Modul halbeinfach ist
	und Summen halbeinfacher Moduln wieder halbeinfach sind, folgt die Behauptung. \bewende
\end{description}
% subsection 221 (end)

\subsection[Korollar: $M$ direkte Summe einfacher Untermoduln $\Rightarrow P \le M$ isomorph zu Teilsumme]{Korollar} % (fold)
\label{sub:222}
Sei $M = \bigoplus_{\alpha \in I} N_\alpha$ mit $N_\alpha$ einfach. Ist $P \subseteq M$ Untermodul, dann existiert $J \subseteq I$ mit 
\[
	P \simeq \bigoplus_{\alpha \in J} N_\alpha.
\]
Ist $P$ einfach, dann ist $P \simeq N_\alpha$ für ein $\alpha \in I$, aber nicht unbedingt gleich $N_\alpha$.
\minisec{Beweis}
Nach Satz \ref{sub:219} existiert $J' \subseteq J$ mit $P \oplus \bigoplus_{j \in J'} N_j = M$ (Beweis "(ii)$\Rightarrow$(i)"). Daher ist 
\[
	P \simeq \nicefrac{M}{\bigoplus_{j \in J'} N_j} \simeq \bigoplus_{j \in J \setminus J'} N_j \bewende
\]
% subsection 222 (end)

\subsection{Korollar (Krull-Remak-Schmidt)} % (fold)
\label{sub:223}
Ist 
\[
	M= \bigoplus_{i \in I} N_i = \bigoplus_{k \in K} L_k,
\]
mit $N_i, L_k$ einfach und $I$ endlich. Dann ist $\abs*{I} = \abs*{K}$ und es existiert ein $\pi  \in \Sym(K)$ mit 
$L_k \simeq N_{\pi(k)}$.
\minisec{Beweis}
Durch Induktion über $n = \abs*{I}$. Für $n=1$ folgt $k = 1$, weil $M$ einfach. Im Allgemeinen existiert ein $j \in I$ mit $L_i \simeq N_j$ (nach Korollar \ref{sub:222}). Dann ist
\[
	\bigoplus_{j \not= i} L_j \simeq \nicefrac{M}{L_i} \simeq \nicefrac{M}{N_j} \simeq \bigoplus_{i \not= j} N_i
\]
Nach Induktionsvorraussetzung folgt die Behauptung. \bewende
% subsection 223 (end)

\subsection[Bemerkung: $M$ endlich erzeugt $\Leftrightarrow M$ endliche direkte Summe einfacher Untermoduln]{Bemerkung} % (fold)
\label{sub:224}
Sei $M = \bigoplus_{i \in I} S_i$ mit $S_i$ einfach. Dann ist $M$ endlich erzeugt genau dann, wenn $\abs*{I}$ endlich ist.
\minisec{Beweis}
\begin{description}
	\item["$\Leftarrow$":] klar, weil $S_i$ zyklisch.
	\item["$\Rightarrow $":] Ist $M = \langle x_1, \ldots , x_r \rangle$, dann existiert für jedes $j=1, \ldots ,r $ endlich viele $S_{j1}, \ldots , S_{jk}$ mit 
	$x_j \in \bigoplus_{i=1}^k S_{j i}$, also ist $M$ die Summe von endlich vielen $S_i$, d.h. $\abs*{I}$ endlich. \bewende
\end{description}
Das heißt wenn $R$ halbeinfach als $R$-Modul ist, dann ist $R$ noethersch und artinsch (nach \ref{sub:214}).
% subsection 224 (end)

\subsection[Satz über die Ideale eines Ringes $R$ und die Ideale in $M_k(R)$]{Satz} % (fold)
\label{sub:225}
Sei $R$ ein Ring. Die Ideale in $M_k(R)$ sind genau von der Form $M_k(I)$ für ein Ideal $I \unlhd R$. Insbesondere ist $M_k(R)$ einfach genau dann, wenn $R$ einfach ist.
\minisec{Beweis}
Klar ist: Wenn $I \unlhd R$, dann $M_k(I) \unlhd M_k(R)$. Sei nun $I \unlhd M_k(R)$ ein Ideal und $\overline{I} := \set[x_{11}]{(x_{ij}) \in I}$ die Menge aller 
$(1,1)$-Koeffizienten in $X \in I$. Man rechnet leicht nach, dass $\overline{I} \unlhd R $ ein beidseitiges Ideal ist:
\[
	\underbrace{\begin{pmatrix}
		j & * \\
		* & *
	\end{pmatrix}}_{\in I} \cdot \underbrace{\begin{pmatrix}
		r & \\
		 & 0
	\end{pmatrix}}_{\in M_k(R)} = \underbrace{\begin{pmatrix}
		j \cdot r & \\
		& 0
	\end{pmatrix}}_{\in I}
\]
Sei $E(s,t) \in M_k(R)$ mit 
\[
	E(s,t)_{\mu, \nu} = \begin{cases}
		1, &\text{ falls }s= \mu, t= \nu \\
		0, &\text{ sonst }
	\end{cases}
\]
eine Elementarmatrix. Dann ist $E(s,t) X E(u,v) = x_{t,u} E(s,v)$. Für $X \in I$ folgt aufgrund der Tatsache  $E(1,s) X E(t,1) = x_{s,t} E(1,1)$, also $x_{s,t} \in \overline{I}$. Daher ist
$I \subseteq M_k(\overline{I})$.

Noch zu zeigen: $M_k(\overline{I}) \subseteq I$. Sei also $Y \in M_k(\overline{I})$. Dann existiert für $s,t$ ein $X \in I$ mit $y_{s,t} = x_{1,1}$. Das heißt 
$y_{s,t} E(s,t) = x_{1,1} E(s,t) = E(s,1) X E(1,t) \in I$. Daher ist 
\[
	Y = \sum y_{s,t} E(s,t) \in I
\]
d.h. $I= M_k(\overline{I} )$. \bewende \smallskip \\
\fbox{$\Longrightarrow$ Matrizenringe über Körpern und Schiefkörpern sind einfache Ringe.}
% subsection 225 (end)

\subsection{Satz (Schurs Lemma)} % (fold)
\label{sub:226}
Ist $M \in {_R \Mod}$ einfach, dann ist $\End_R(M)$ ein Schiefkörper.
\minisec{Beweis}
Wegen $M \not= 0$ ist $\End_R(M) \not= 0$. Ist $\varphi \in \End_R(M) \setminus \set{0}$, dann ist $\varphi(M) \not= 0$, also $\varphi(M)=M$, da $M$ einfach. Ebenso
$\ker \varphi \not= M$, daher $\ker \varphi = 0$. Daher ist $\varphi$ ein Isomorphismus und hat ein Inverses in $\End_R(M)$. \bewende 
% subsection 226 (end)

\subsection[Lemma: Darstellung von $\varphi \in \End_R(M)$ als Matrix]{Lemma} % (fold)
\label{sub:227}
Ist $M = \bigoplus_{i \le k} M_i$, $\varphi \in \End_R(M)$, dann existieren $\varphi_{i,j} \in \Hom(M_i, M_j)$ für $1 \le i,j \le k$, so dass für alle 
$x = (x_1, \ldots , x_k) \in M_1 \oplus \ldots \oplus M_k$ gilt
\[
	\varphi(x) = \begin{pmatrix}
		\varphi_{1,1} & \cdots & \varphi_{i,k} \\
		\vdots & &\vdots \\
		\varphi_{k,1} & \cdots & \varphi_{k,k}
	\end{pmatrix}\cdot 
	\begin{pmatrix}
		x_1 \\ \vdots \\ x_k
	\end{pmatrix}
\]
Insbesondere ist 
\[
	\End_R(M^k) \simeq M_k \enbrace[\big]{\End_R(M)}
\]
\minisec{Beweis}
Seien $e_j : M_j \to M$ die Einbettungen, $\pi_i : M \to M_i$ die Projektionen. Dann ist für jeden Endomorphismus $\varphi \in \End_R(M)$
\[
	\varphi = \sum e_i \circ \varphi_{i,j} \circ \pi_j
\]
wobei $\varphi_{i,j} := \pi_i \circ \varphi \circ e_j \in \Hom(M_j, M_i)$.
Es ist leicht nachzurechnen, dass die Verknüpfung von Endomorphismen der Multiplikation von Matrizen entspricht, d.h. $\varphi \psi = \sum \varphi_{i,j} \psi_{j,k}$. Damit
folgt die Behauptung. \bewende
% subsection 227 (end)

\subsection[Definition: Entgegengesetzter Ring $R^{\op} \simeq \End_R(R,+)$]{Definition + Lemma} % (fold)
\label{sub:228}
Ist $(R,+,\cdot ,1)$ ein Ring, dann ist $R^{\op} := R(+, *, 1)$ ein Ring mit $r * s = s \cdot r$, der \Index{entgegengesetzte Ring}. \\
\uline{Klar:} Ist $R$ kommutativ, dann ist $R^{\op} \simeq R$. Jeder Links-$R$-Modul lässt sich als Rechts-$R^\op$-Modul auffassen.
Es gilt
\[
	\End_R(R,+) \simeq R^\op
\]
\minisec{Beweis}
Setze $\rho : R^\op \to \End_R(R,+)$, $\rho(r)= \rho_r : x \mapsto x \cdot r$. Homomorphie: $\rho_{r * s} : x \mapsto x \cdot s \cdot r$. Dann ist $\rho$ ein
Ringhomomorphismus mit $\ker \rho = 0$, also injektiv. Ist $\varphi \in \End_R(R,+)$, so betrachte $\varphi(1)= r \in R$. Dann ist 
\[
	\rho(r)(x) = x \cdot r = x \cdot \varphi(1) = \varphi(x \cdot 1) = \varphi(x)
\]
d.h. $\varphi = \rho_r$. Daher ist $\rho$ ein Isomorphismus. \bewende
% subsection 228 (end)

\subsection{Satz (Wedderburn, 1. Struktursatz)} % (fold)
\label{sub:229}
Sei $R \not= 0$ ein Ring. Dann sind äquivalent:
\begin{enumerate}[(i)]
	\item $R$ ist einfach als Ring und links-artinsch als $R$-Modul.
	\item Alle einfachen $R$-Moduln sind isomorph und $R$ ist halbeinfach als Modul. 
	\item $R \simeq M_k(D)$ für einen Schiefkörper $D$ und ein $k\ge 1$. 
\end{enumerate}
Darüber hinaus sind $k$ und $D$ eindeutig bestimmt und der einfache $R$-Modul ist isomorph zu $D^k$.
\minisec{Beweis}
\begin{description}
	\item["(i)$\Rightarrow $(ii)":] Sei $0  \not= I \subseteq R$ ein minimales Linksideal (existiert, da $R$ links-artinsch als Modul). Das heißt $I = R c \not= 0$ für ein
	$c \in R$ (zyklisch!). Dann ist 
	\[
		\sum_{r \in R} R c \cdot r = J
	\] 
	ein beidseitiges Ideal im Ring $R$, also $R= \sum_{r\in R} Rc\cdot r$, da $R$ einfach. Ist $Rc r \not= 0$, dann ist $R c \to R c r$, $s\cdot c\mapsto s \cdot c\cdot r$
	ein Epimorphismus, also $R c r = R c$ wegen der Minimalität von $R c$. Daher ist $(R,+)$ Summe von einfachen Untermoduln, also halbeinfach nach Satz \ref{sub:219}.
	
	Noch zu zeigen: Alle einfachen $R$-Moduln sind isomorph. \\
	Ist $M$ ein einfacher $R$-Modul, dann ist $M = R \cdot m$ zyklisch und also $M \simeq \nicefrac{R}{I}$ für ein \marginnote{siehe \ref{sub:26}}maximales Linksideal
	$I$ in $(R,+)$. Weil $R$ halbeinfach
	ist, ist $M$ isomorph zu einem minimalen Linksideal in $(R,+)$ und wegen $R = \bigoplus_{i=1}^k R c_i$ nach Satz \ref{sub:219} ist dann $M \simeq R c_i \simeq Rc$. Die 
	direkte Summe ist endlich, da $R$ artinsch ist nach Bemerkung \ref{sub:224}.
	\item["(ii)$\Rightarrow $(iii)":]  $(R,+)$ ist als Links-$R$-Modul endlich erzeugt (sogar zyklisch, $R= R \cdot 1$). Daher ist nach Vorraussetzung 
	\[
		(R,+) \simeq \bigoplus_{i=1}^k L_i
	\]
	mit $L_i$ minimale Linksideale in $R$, $L_1 \simeq \ldots  \simeq L_k$ und nach Schurs Lemma (\ref{sub:226}) ist $\End_R(L_i) \simeq D'$ ein Schiefkörper. Daher ist
	\[
		R^\op \simeq \End_R(R,+) \simeq M_k(D')
	\]
	nach \ref{sub:227} und damit $R \simeq M_k(D)$ mit $D= (D')^\op$, da $M_k(R)^\op = M_k(R^\op)$.
	\item["(iii)$\Rightarrow$(i)":] $M_k(D)$ ist ein $k^2$-dimensionaler $D$-Vektorraum. Linksideale sind $D$-Untervektorräume, also ist $M_k(D)$ noethersch und artinsch
	\marginnote{artinsch, noethersch, da Dimension endlich}% 
	und nach Satz \ref{sub:225} einfach, da $D$ ein Schiefkörper ist. 
	Die Eindeutigkeit von $k$ und $D$ folgt aus Korollar \ref{sub:223} (Krull-Remak-Schmidt) und Schurs Lemma (\ref{sub:226}). 
	\bewende
	\minisec{Bemerkung}
	$(R,+)$ ist also auch noethersch (rechts und links). Achtung: Es gibt Ringe, die links-noethersch aber nicht rechts-noethersch sind!
\end{description}
% subsection 229 (end)

\subsection[Bemerkung: Homomorphismen zwischen einfachen, nicht-isomorphen Moduln sind trivial]{Bemerkung} % (fold)
\label{sub:230}
Sind $M,N$ einfache, nicht-isomorphe $R$-Moduln, dann ist $\Hom_R(M,N) = 0$. (Klar!)
% subsection 230 (end)

\subsection{Satz (2. Struktursatz von Wedderburn)} % (fold)
\label{sub:231}
Sei $R \not= 0$ ein Ring, halbeinfach als $R$-Modul. Dann existieren Schiefkörper $D_1, \ldots , D_l$, die paarweise nicht isomorph sind und 
$k_1, \ldots , k_l \in \mathds{N}$, so dass 
\[
	R \simeq M_{k_1}(D_1) \oplus \ldots \oplus M_{k_l}(D_l).
\]
als Ringisomorphismus.
\minisec{Beweis}
$(R,+) \simeq L_1^{k_1} \oplus \ldots \oplus L_l^{k_l}$ mit $L_i \subseteq R$ minimale Linksideale, $L_i$ paarweise nicht isomorph. Dann ist
\[
	R^\op \simeq \End_R(R,+) \stackrel{\text{\ref{sub:230}}}{\simeq} \End_R(L_1^{k_1}) \oplus \ldots \oplus \End_R(L_l^{k_l})  
	\simeq M_{k_1}(\tilde D_1) \oplus \ldots \oplus M_{k_l}(\tilde D_l)
\]
mit $\tilde D_i = \End_R(L_i)$. Setze $D_i := \tilde D_i^\op$. Damit ist dann 
\[
	R \simeq M_{k_1} \oplus \ldots \oplus M_{k_l}(D_l) \bewende
\]
% subsection 231 (end)
\addcontentsline{toc}{section}{Freie Moduln}
\section*{Freie Moduln}
\subsection[Definition: linear unabhängig und Basen]{Definition} % (fold)
\label{sub:232}
Sei $M$ ein Links-$R$-Modul. Elemente $x_1, \ldots , x_n \in M$ heißen \bet{$R$-linear unabhängig}\index{linear unabhängig in Moduln}, wenn gilt: 
\[
	\alpha_1, \ldots , \alpha_n \in R \enspace : \enspace\sum \alpha_i \cdot x_i = 0 \quad \Longrightarrow \quad \forall i : \alpha_i = 0
\]
\marginnote{äquivalent dazu: minimale erzeugende Menge}%
Sonst heißen $x_1, \ldots , x_n$ linear abhängig. Eine Menge $X \subseteq M$ heißt linear unabhängig, wenn jede endliche Teilmenge linear abhängig ist. Eine 
linear unabhängige Erzeuger-Menge heißt \bet{Basis}\index{Basis (Modul)} für $M$.
% subsection 232 (end)

\subsection[Definition: Freier Modul]{Definition} % (fold)
\label{sub:233}
Sei $R$ ein Ring. Ein $R$-Modul heißt \Index{frei}, wenn er eine Basis hat. \\
\uline{Achtung:} Die Mächtigkeit einer Basis ist \emph{nicht} notwendig eindeutig bestimmt!
\minisec{Beispiele}
\begin{enumerate}[(i)]
	\item Wenn $R=K$ für einen Körper $K$, dann sind alle $R$-Moduln frei.
	\item $R$ als $R$-Modul ist frei mit der Basis $\set{1_R}$.
	\item $R^n$ für beliebiges $n \in \mathds{N}$ ist frei.
\end{enumerate}
% subsection 233 (end)

\subsection[Definition: Freier $R$-Modul mit Basis der Mächtigkeit $\abs*{I}$]{Definition} % (fold)
\label{sub:234}
Für eine beliebige Menge $I$ heißt der (Links-)$R$-Modul $\mathcal{F}_I = \bigoplus_{I} R$ \emph{der} freie $R$-Modul mit Basis der Mächtigkeit $\abs*{I}$. \smallskip \\
Elemente von $\mathcal{F}_I$ sind von der Form $(a_i)_{i \in I}$ mit $a_i=0$ für fast alle $i \in I$, d.h. $a_i = 0$ für alle bis auf endlich viele $i \in I$. Jeder andere
freie $R$-Modul mit Basis der Mächtigkeit $\abs*{I}$ ist isomorph zu $\mathcal{F}_I$ via Bijektion der Basen: $(u_i)_{i \in I}$ und $(v_i)_{i \in I}$ induzieren einen 
Isomorphismus durch
\[
	\sum_{i \in I} \alpha_i \cdot u_i \mapsto \sum_{i \in I} \alpha_i \cdot v_i
\]
\uline{Achtung:} Die Umkehrung gilt nicht! (siehe \ref{sub:233}) \smallskip\\
Freie Moduln können durch ihre \bet{universelle Eigenschaft}\index{universelle Eigenschaft!freier Modul} charakterisiert werden:
% subsection 234 (end)

\subsection[Satz: Universelle Eigenschaft der freien $R$-Moduln]{Satz} % (fold)
\label{sub:235}
Sei $R$ ein Ring. Dann existiert für jede Menge $I$ ein $R$-Modul $\mathcal{F}_I $ und eine Abbildung $\varphi : I \to \mathcal{F}_I $, die universell ist für 
$R$-Moduln. D.h. für jede Abbildung $f : I \to M$ in einen $R$-Modul $M$ existiert ein eindeutiger Homomorphismus $f' : \mathcal{F}_i  \to M $ mit $f = f' \circ  \varphi$.
Also
\[
	\begin{tikzcd}
		I \rar{\varphi} \drar[swap]{f}& \mathcal{F}_I \dar[dashed]{f'} \\
		& M 
	\end{tikzcd}
\]
\minisec{Beweis}
Sei $\mathcal{F}_I = \bigoplus_I R$ mit Basis $(u_i)_{i \in I}$ und $\varphi(i)=u_i$. Ist $f : I \to M$ eine Abbildung, dann gilt für $f' : \mathcal{F}_I \to M$ mit 
$f=  f' \circ \varphi$ offensichtlich $f'(u_i)= f(i)$. Also muss gelten 
\[
	f'\enbrace*{\sum \alpha_i u_i} = \sum \alpha_i f(i) \tag*{($\star$)}
\]
Daher folgt die Existenz von $f'$ und die Eindeutigkeit ebenfalls. \bewende

\minisec{Bemerkung}
Aus der universellen Eigenschaft folgt, dass der freie $R$-Modul mit Basis der Mächtigkeit $\abs{I}$ bis auf Isomorphie eindeutig bestimmt ist.
% subsection 235 (end)

\subsection[Korollar: Freie Moduln sind projektiv]{Korollar} % (fold)
\label{sub:236}
Ist 
\[
	\begin{tikzcd}
		0 \rar& M' \rar["\alpha"]& M \rar["\beta"]& M'' \rar \lar["\sigma", bend left, dashed]& 0
	\end{tikzcd}
\]
eine kurze exakte Sequenz und $M''$ ein freier $R$-Modul, dann \Index{spaltet} die Sequenz, d.h. es 
existiert ein $R$-Modulhomomorphismus $\sigma : M''  \to M$ mit $\beta \circ \sigma = \id_{M''}$. Dann ist $M \cong M' \oplus M''$. $\sigma$ heißt dann ein \Index{Schnitt}. 
\minisec{Bemerkung}
$M''$ heißt \Index{projektiv} genau dann, wenn jede kurze Sequenz mit $M''$ an dritter Stelle spaltet. Also bedeutet dieses Korollar: Freie $R$-Moduln sind projektiv. 
\minisec{Beweis}
Sei $I$ Basis für $M''$. Da $\beta$ surjektiv ist, existiert für jedes $i \in I$ ein $u_i \in M$ mit $\beta(u_i)= i$. Die Abbildung $f : I \to M$, $i \mapsto u_i$ lässt 
sich fortsetzen zu $\sigma : M'' \to M$ (nach \ref{sub:235}) und dann gilt $\beta \circ \sigma = \id_{M''}$. Damit ist 
$M = \im \alpha \oplus \im \sigma \simeq M' \oplus M''$. \bewende
% subsection 236 (end)

\subsection[Satz: Jeder $R$-Modul ist Quotient eines freien $R$-Moduls]{Satz} % (fold)
\label{sub:237}
Jeder $R$-Modul ist Quotient eines freien $R$-Moduls, also $M \simeq \nicefrac{\mathcal{F}}{G}$ für einen Untermodul $G \le \mathcal{F}$.
\minisec{Beweis}
Ist $M \in {_R \Mod}$, setze $\mathcal{F}_M = \bigoplus_{m \in M} R$. Dann lässt sich die Abbildung
$f : M \to M$ fortsetzen zu einem Epimorphismus $f' : \mathcal{F}_M \to M$, d.h. $M \simeq \nicefrac{\mathcal{F}_M}{\ker f'}$. \bewende\smallskip\\
Daher sind Quotienten von freien Moduln im Allgemeinen \emph{nicht} frei! \\
Nur über Körpern sind alle Moduln frei!
% subsection 237 (end)

\subsection[Satz: Alle $R$-Moduln sind frei $\iff$ $R$ ist Schiefkörper]{Satz} % (fold)
\label{sub:238}
Sei $R$ ein nicht-trivialer Ring. Dann sind äquivalent:
\begin{enumerate}[(i)]
	\item Jeder (Links-)$R$-Modul ist frei.
	\item Jeder zyklischer (Links-)$R$-Modul ist frei.
	\item $R$ ist einfach als Links-$R$-Modul.
	\item Jedes $x \in R \setminus \set{0}$ hat ein Linksinverses.
	\item $R$ ist ein Schiefkörper.
\end{enumerate}
(i) - (v) sind auch äquivalent zu $\text{(i)}_R$ - $\text{(iv)}_R$ für Rechts-$R$-Moduln.
\minisec{Beweis}
\begin{description}
	\item["(i)$\Rightarrow$(ii)":] Klar. 
	\item["(ii)$\Rightarrow $(iii)":] Sei $I$ ein maximales Links-Ideal in $(R,+)$. Dann ist 
	\[
		0 \longrightarrow I  \longrightarrow R \longrightarrow \mathcal{F} \longrightarrow 0
	\]
	mit $\mathcal{F} \simeq \nicefrac{R}{I}$ eine exakte Sequenz mit $\mathcal{F}$ zyklisch und daher frei nach Vorraussetzung. Da $I$ ein maximales Ideal war, ist 
	$\mathcal{F}$ einfach. $R$ ist auch zyklisch und daher frei, d.h. $R \simeq \mathcal{F}$ und daher auch $R$ einfach als $R$-Modul.
	\item["(iii)$\Rightarrow $(iv)":] Sei $c \in R \setminus \set{0} $. Dann ist $Rc=R$, also existiert $b \in R$ mit $b c = 1$, d.h. $c$ hat ein Linksinverses.
	\item["(iv)$\Rightarrow$(v)":] Sei $c \in R \setminus \set{0}$ mit Linksinversem $b$ und $a$ Linksinverses von $b$, also
	\[
		b c = 1 = a b.
	\]
	Dann ist also $a = a (b c) = (a b) c = c$. Dann folgt $a=c$ und $b c= c b = 1$. Damit ist $R$ ein Schiefkörper.
	\item["(v)$\Rightarrow$(i)":] Lineare Algebra I. (Vektorräume haben eine Basis)
\end{description}
Die Äquivalenz von $\text{(i)}_R$ - $\text{(v)}_R$ folgt entsprechend. Wegen (v)=$\text{(v)}_R$ (weil (v) symmetrisch ist), sind also (i)- (v) auch äquivalent zu
$\text{(i)}_R$ - $\text{(iv)}_R$. \bewende

\minisec{Achtung}
Auch freie Moduln haben nicht notwendig eine Dimension! \\
\uline{Beispiel:} Sei $V$ ein unendlich dimensionaler $K$-Vektorraum, dann ist $V \simeq V \oplus V$, also ist 
\[
	R = \End_K(V) \simeq \End_K(V^2) \simeq M_2(R) \simeq R^4
\]
Das heißt $R$ als freier $R$-Modul hat eine Basis der Mächtigkeit 1, aber auch Basen jeder anderen Mächtigkeit. Genaueres in Aufgabe 1.b) von Blatt 5.
% subsection 238 (end)

\subsection[Definition: Invariante Basislänge (IBL)]{Definition} % (fold)
\label{sub:239}
Ein Ring hat \Index{invariante Basislänge} (IBL), wenn aus $R^m \simeq R^n$ (als Links-Moduln) schon $m=n$ folgt.
\minisec{Bemerkung}
\begin{itemize}
	\item Körper haben IBL, $\End_R(V)$ nicht immer, $0$ nie!
	\item Man kann zeigen: $R$ hat IBL genau dann, wenn jeder freie $R$-Modul eindeutige Basislänge hat.
\end{itemize}
% subsection 239 (end)

\subsection[Lemma: Charakterisierung von IBL in Matrizenformulierung]{Lemma} % (fold)
\label{sub:240}
Ein Ring $R$ hat in IBL, wenn für alle $m,n \in \mathds{N}$ gilt: Ist $A \in R^{m \times n}$, $B \in R^{n \times m}$ mit 
\[
	AB = \ind_{m \times m}\quad , 	\quad  BA= \ind_{n \times n} 
\]
dann ist $n=m$.
\minisec{Beweis}
Dies ist die Matrizenformulierung mit $R^m \simeq R^n$ mit Lemma \ref{sub:227}. \bewende
\minisec{Folgerung}
$R$ hat IBL für Links-$R$-Moduln genau dann, wenn $R$ IBL für Rechts-$R$-Moduln hat.
% subsection 240 (end)

\subsection[Satz: Kommutative, noethersche Ringe und Urbilder von Ringhomom. haben IBL]{Satz} % (fold)
\label{sub:241}
Sei $R \not= 0$ ein Ring.
\begin{enumerate}[(i)]
	\item Ist $R$ kommutativ, dann hat $R$ IBL.
	\item Ist $R$ noethersch (bzw. artinsch), dann hat $R$ IBL. Insbesondere haben halbeinfache Ringe IBL.
	\item Ist $\varphi : S \to R$ ein Ringhomomorphismus ($\varphi(1_S)=1_R$) und $R$ hat IBL, dann auch $S$. Insbesondere vererbt sich IBL auf Unterringe mittels Inklusion.
\end{enumerate}
\minisec{Beweis}
\begin{enumerate}[(i)]
	\item Angenommen $A \in R^{m \times n}$, $B \in R^{n \times m}$ mit 
	\(
		AB = \ind_{m \times m}, BA= \ind_{n \times n} 
	\)
	und o.B.d.A. $m < n$. Dann ist
	\[
		\underset{\tilde{B}}{\begin{BMAT}(e,2pt,3pt){|ccc|c|}{|ccc|}
			& & & 0 \\
			& B & &   \\
			& & & 0
		\end{BMAT}} \cdot 
		\underset{\tilde{A}}{\begin{BMAT}(e,2pt,4pt){|ccc|}{|ccc|c|}
			& &  \\
			& A&  \\
			&  &  \\
			0 & & 0
		\end{BMAT}} = \ind
	\]
	Es gilt $\det( \tilde{B}\tilde{A})= \det \ind = 1 $ und $\det(\tilde{A}\tilde{B}) = \det(\tilde{B}\tilde{A})$, aber 
	\[
		\tilde{A} \cdot \tilde{B} = 
		\begin{BMAT}(e,2pt,4pt){|ccc|}{|ccc|c|}
			& &  \\
			& A&  \\
			&  &  \\
			0 & & 0
		\end{BMAT} \cdot 
		\begin{BMAT}(e,2pt,3pt){|ccc|c|}{|ccc|}
			& & & 0 \\
			& B & &   \\
			& & & 0
		\end{BMAT} =
		\begin{BMAT}(e){|cccc|}{|cccc|}
			& & & 0\\
			& 1 & & \\
			& & & \\
			0 & & & 0
			\addpath{(0,1,|)rrruuu}
		\end{BMAT}
	\]
	Also $\det(\tilde{A} \cdot \tilde{B}) = 0$ \light
	\item Ist $R^m \simeq R^{m+k}$ mit $k \ge 1$, dann ist $R^m \simeq R^{m+k} \simeq R^{m+2k}$ und man erhält eine unendliche, echt aufsteigende (absteigende) Kette in
	$R^m$. Mit $R$ ist auch $R^m$ noethersch (bzw. artinsch). \light
	
	Halbeinfacher Ring $\Rightarrow $ noethersch und artinsch.
	\item Seien $A \in S^{m \times n}$, $B \in S^{n \times m}$ mit $AB = \ind$, $BA = \ind$. Dann gilt 
	\[
		\varphi(A) \varphi(B)= 1 \quad , \quad \varphi(B) \varphi(A) = 1
	\]
	das heißt es gilt $n=m$. \bewende
\end{enumerate}
% subsection 241 (end)

\subsection*{Erinnerung} % (fold)
\label{sub:erinnerung_hir}
$R$ ist ein \Index{Hauptidealring} (HIR)\footnote{englisch: PID, prime ideal domain} , falls $R$ kommutativ und nullteilerfrei ist und jedes Ideal ein Hauptideal ist, das 
heißt von einem Element erzeugt ist. In Hauptidealringen hat man eine eindeutige Primfaktorzerlegung (bis auf Einheiten).
% subsection erinnerung (end)

\subsection[Satz: Für $R$ HIR sind Untermoduln freier $R$-Moduln frei]{Satz} % (fold)
\label{sub:242}
Sei $D$ ein Hauptidealring, $L \subseteq D^m$ ein Untermodul, dann ist $L \simeq D^n$ für ein $n \le m$, d.h. Untermoduln freier Moduln sind frei.
\minisec{Beweis}
Durch Induktion über $m \ge 0$. Für $m=0$ ist die Aussage klar. Für $m=1$ ist $L \subseteq D$ ein Linksideal und da $D$ kommutativ ist, ist $L \unlhd D$. Weil $D$
Hauptidealring ist, folgt $L=D a$ für ein $a \in D$, d.h. $L$ ist frei. Weiter ist $D \simeq D a$ via $\varphi : D \to D a$, $s \mapsto s \cdot a$ (Isomorphismus, da D
nullteilerfrei). \\
\uline{Induktionsschritt:} Sei $L \subseteq D^{m+1} = D \oplus D^m$, $\pi  : D \oplus D^m \to D$ die Projektion auf die erste Komponente. Betrachte die exakte Sequenz
\[
	0 \longrightarrow \ker \pi\big|_L \longrightarrow L \xrightarrow{\enspace \pi \enspace} \pi(L) \longrightarrow 0 
\]
\begin{description}
	\item[Fall 1:] $\pi (L)=0$, dann ist $L \subseteq \ker \pi = 0 \oplus D^m \simeq D^m$ und nach Induktionsvorraussetzung ist $L$ frei, $L \simeq D^n$, $n \le m$.
	\item[Fall 2:] $\pi (L) \not= 0$, dann ist $\pi(L) \subseteq D$ ein freier $D$-Modul. Nach Korollar \ref{sub:236} ist 
	\[
		L \simeq \ker \pi\big|_L \oplus \pi(L)
	\]
	Nach Induktionsvorraussetzung ist $\ker \pi\big|_L \subseteq D^m$ ein freier Modul, also $\ker \pi\big|_L \simeq D^n$, $n \le m$, und $\pi (L) \simeq D$, daher ist 
	$L \simeq D \oplus D^n$. \bewende
\end{description}
	Ist $D$ ein Hauptidealring, $M$ endlich erzeugt über $D$, also $M = D m_1 + \ldots + D m_s$. Betrachte $\varphi : D^s \to M$ mit $(d_1, \ldots , d_s) \mapsto \sum d_i m_i$.
	Nach Satz \ref{sub:242} ist $\ker \varphi \simeq D^t$ für $t \le s$ und $M \simeq \nicefrac{D^s}{\ker \varphi}$. Um $M$ zu beschreiben, müssen wir untersuchen, 
	\emph{wie} $\ker \varphi$ in $D^s$ liegt.
% subsection 242 (end)

\subsection[Satz: (Smith-Normalform) Invariante Faktoren einer Matrix über einem Hauptidealring]{Satz} % (fold)
\label{sub:243}
Sei $A= (a_{ij}) \in D^{m \times n}$, $D$ ein Hauptidealring. Dann gibt es invertierbare Matrizen $P  \in D^{m \times m}, Q \in D^{n \times n}$ mit 
\[
	P \cdot A \cdot Q = \begin{pmatrix}
		d_1 & & & 0 & & \\
		& \ddots & & & & \\
		& & d_k & & & \\
		0 & & & 0 & & \\
		& & & & \ddots & \\
		& & & & & 0
	\end{pmatrix}
	\qquad d_1, \ldots , d_k \not= 0, d_i \mid d_{i+1}
\]
Falls $D$ ein Körper ist, folgt dies aus dem Gauß-Verfahren und es gilt $d_i =1$ für alle $i$.
\minisec{Beweis}
Für $a \in D \setminus \set{0}$ sei $l(a)=s$, falls $a= p_1 \cdot \ldots \cdot p_s$ mit $p_i$ prim (eindeutig, da $D$ HIR). Ist $a \in D^\times$ Einheit, setze $l(a)=0$
und $l(0)= \infty$. Für $A=0$ ist nichts zu zeigen. \\
\uline{Erinnerung:} Durch Links- und Rechtsmultiplikation mit geeigneten invertierbaren Matrizen können wir:
\begin{itemize}
	\item Zeilen von $A$ vertauschen
	\item Spalten von $A$ vertauschen
	\item Zeilen/Spalten mit Einheiten multiplizieren
	\item zu einer Zeile/Spalte beliebige Vielfache einer anderen Zeile/Spalte aufaddieren \marginnote{siehe LA I.}
\end{itemize}
Wähle $a_{ij} $ in $A$ mit $l(a_{ij})$ minimal. Durch Zeilen- und Spaltenvertauschungen können wir erreichen, dass $a_{ij}$ links oben steht (vertausche $i$-te und 1.
Zeile, dann $j$-te und 1. Spalte). 

Wenn $a_{11} \nmid a_{1k}$ für $k \ge 2$, vertausche 2-te und $k$-te Spalte. Dann gilt $a_{11} \nmid a_{12}$. Sei $d \in \ggT(a_{11}, a_{12})$,
dann $l(d) < l(a_{11})$. Schreibe $d = a_{11} \cdot x + a_{12} \cdot y$\marginnote{$x,y \in D$} und $d \cdot e = a_{12}$ sowie $ d \cdot f = - a_{11}$. Dann ist 
\[
	\begin{pmatrix}
		-f & e \\
		y & -x
	\end{pmatrix} \cdot \begin{pmatrix}
		x  & e \\
		y & f
	\end{pmatrix} = \begin{pmatrix}
		- f \cdot x + y \cdot e & 0 \\
		0 & e \cdot y - f \cdot x
	\end{pmatrix}
	= \begin{pmatrix}
		1 & 0 \\
		0 & 1
	\end{pmatrix}
\]
denn $d \cdot e \cdot y - d \cdot f \cdot x = a_{12} \cdot y + a_{11} \cdot x = d$, daher $e \cdot y- f x \cdot = 1$. Ebenso 
\[
	\begin{pmatrix}
		x & e \\
		y & f
	\end{pmatrix} \cdot \begin{pmatrix}
		- f & e \\
		y & -x
	\end{pmatrix} = \begin{pmatrix}
		1 & 0\\
		0 & 1
	\end{pmatrix}
\]
das heißt $\begin{psmallmatrix}
	x & e \\
	y & f
\end{psmallmatrix}$ ist invertierbar. Es gilt \marginnote{$a_{11} \cdot e +a_{12}\cdot f  = - d \cdot f \cdot e + d \cdot e \cdot f = 0$}
\[
	\begin{pmatrix}
		\begin{BMAT}(c){:cc:}{:cc:}
			a_{11} & a_{12} \\
			a_{21} & a_{22}
		\end{BMAT} & * \\
		* & *
	\end{pmatrix} \cdot \begin{pmatrix}
		\begin{BMAT}(c){:cc:}{:cc:}
			x & e \\
			y & f
		\end{BMAT} & 0 \\
		0 & \begin{matrix}
			1  & & \\
			& \fixedddots & \\
			& & 1
		\end{matrix}
	\end{pmatrix} = 
	\begin{pmatrix}
		d & a_{11}e +a_{12}f & a_{13} \cdots  \\
		* & * & * \\
		* & * & * 
	\end{pmatrix} =
	\begin{pmatrix}
		d & 0 & a_{13} \cdots  \\
		* & * & * \\
		* & * & * 
	\end{pmatrix}
\]
Damit transformieren wir nach und nach die Matrix $A$ auf die Gestalt
\[
	\tilde{A} = \begin{pmatrix}
		\tilde{a_{11}} & \tilde{a_{12}} & \ldots \\
		* & * & *
	\end{pmatrix}
\]
mit $\tilde{a}_{11} \mid \tilde{a}_{ik}$ und $\tilde{a}_{11} \mid \tilde{a}_{k1}$. Durch Addieren von geeigneten Vielfachen der $1$-ten Zeile/Spalte erhalten wir eine 
Matrix der Form
\[
	\left( \begin{BMAT}(@,1pt,1.7em){cc}{cc}
		\tilde{a}_{11} & \, 0  \enspace\cdots  \enspace 0 \, \\
		\begin{matrix} % (c){c}{ttt}
			0\\ \vdots \\ 0
		\end{matrix} & B 
		\addpath{(1,0,:)ur}
	\end{BMAT} \right)
\]
und $l(\tilde{a}_{11}) \le l (b_{ij})$ für $B= (b_{ij})$. Induktiv erhalten wir eine Matrix der Form
\[
	\begin{pmatrix}
		d_1 &  & \\
		 & \ddots  & &\\
		  & & d_n & \\
		  & & & 0
	\end{pmatrix}
\]
mit $l(d_i) \le l(d_{i+1})$. Ist $d_i \nmid d_{i+1}$. Dann transformieren wir weiter:
\[
	\begin{pmatrix}
		d_1 &  \\
		& d_{i+1}
	\end{pmatrix} \leadsto \begin{pmatrix}
		d_i & d_{i+1} \\
		& d_{i+1}
	\end{pmatrix} \leadsto \begin{pmatrix}
		d_i & d_{i+1} \\
		& d_{i+1}
	\end{pmatrix} \cdot \begin{pmatrix}
		x & e \\
		y & f
	\end{pmatrix} = \begin{pmatrix}
		d_i' & 0 \\
		d_{i+1}y & d_{i+1} f
	\end{pmatrix}
\]
mit $d_i x + d_{i+1} y = d_i' \in \ggT(d_i, d_{i+1})$. Wie eben erhalten wir $\begin{psmallmatrix}
	d_i' & 0 \\
	0 & d_{i+1} f
\end{psmallmatrix} \leadsto \ldots $ fertig. \bewende \medskip\\
Die $d_i, \ldots , d_k$ heißen die \Index{invarianten Faktoren} von $A$. Diese sind eindeutig (bis auf Einheiten). 
% subsection 243 (end)

\subsection[Definition: $i$-Minoren und Rang einer Matrix]{Definition} % (fold)
\label{sub:244}
Ist $A \in D^{m \times n}$ und $i \le m,n$, dann ist die Determinante einer Matrix $A' \in D^{i \times i}$ ein $i$-\Index{Minor} von $A$, wenn $A'$ aus $A$ hervorgeht
durch Streichen von $m-i$ Zeilen und $n-i$ Spalten. Der \Index{Rang} von $A$ ist das größte $i$, für das $A$ einen $i$-Minor $\not= 0$ besitzt.
% subsection 244 (end)

\subsection{Satz (Elementarteilersatz)} % (fold)
\label{sub:245}
Sei $D$ ein Hauptidealring, $A \in D^{m\times n}$ und $k$ der Rang von $A$. Für $j \ne k$ sei $\delta_j = \delta_j(A)$ ein $\ggT$ aller $j$-Minoren von $A$. Dann sind die
invarianten Faktoren eindeutig bestimmt (bis auf Einheiten) und von der Form $d_1 = \delta_1$, $d_i = \prod_{j \le i} \delta_j$.
\minisec{Beweis}
Ist $k$ der Rang von $A$, dann existiert für alle $i \le k$ ein $i$-Minor $\not= 0$ in $A$ (Laplace-Entwicklung der Determinante). Insbesondere ist 
$\delta_1, \ldots , \delta _k \not= 0$, also $d_i \not= 0$ für $i=1, \ldots ,k$. 

Für $P \in D^{m \times m}$ sind die Zeilen von $P \cdot A$ Linearkombinationen der 
Zeilen von $A$, also sind die $j$-Minoren von $P \cdot A$ Linearkombinationen der $j$-Minoren von $A$. Entsprechendes gilt für die Spalten von $Q \in D^{n \times n}$ und 
$A \cdot Q$. Wenn also $P \in D^{m \times m}, Q \in D^{n \times n}$ invertierbare Matrizen sind, dann folgt $\delta_j (P \cdot A \cdot Q) \mid \delta_j (A)$ und 
$\delta_j(A) \mid \delta_j (P \cdot A \cdot Q)$ ist. Das heißt $\delta_j (P \cdot A \cdot Q) = u \cdot \delta_j(A)$ für eine Einheit $u \in D^\times$. Ist nach Satz 
\ref{sub:243}
\[
	P \cdot A \cdot Q = \begin{pmatrix}
		d_1 & & & 0 & \cdots & 0\\
		& \ddots & & \vdots  & & \vdots \\
		& & d_k & \vdots& & \vdots \\
		0 & \cdots & \cdots & 0 & \cdots & 0 \\
		\vdots & & & \vdots & \ddots & \vdots \\
		0 & \cdots & \cdots & 0 & \cdots & 0
	\end{pmatrix}
\] 
mit $d_i \mid d_{i+1}$, dann ist $\delta_j(P \cdot A \cdot Q) = d_i \cdot \ldots \cdot d_j$. Dies zeigt, dass die invarianten Faktoren bis auf Einheiten eindeutig sind.
\bewende\\
Die Matrix heißt die \Index{Smith-Normalform} von A. 
% subsection 245 (end)

\subsection[Beispiel: Anwendung des Elementarteilersatzes]{Beispiel} % (fold)
\label{sub:246}
Sei $D= \mathds{Z}$ und $A = \begin{psmallmatrix}
	3 & 1 \\
	0 & 4
\end{psmallmatrix}$. Die $1$-Minoren sind $1,3,4$, also ist $\delta_1 = 1 = d_1$. Der $2$-Minor ist $12$, also gilt $\delta_2=12= d_2$. Damit ist 
$A \sim_\mathds{Z} \begin{psmallmatrix}
	1 & 0 \\
	0 & 12
\end{psmallmatrix}$.
% subsection 246 (end)

\subsection[Definition: Annulatorideal und Torsionselement]{Definition} % (fold)
\label{sub:247}
Ist $M \in {_R\Mod}$, $m \in M$, dann heißt 
\[
	\ann(m) = \set[r \in R]{r \cdot m = 0} 
\]
das \Index{Annulatorideal} von $m$ in $R$. $m \in M$ heißt \Index{Torsionselement}, falls $\ann(m) \not= 0$  
% subsection 247 (end)

\subsection{Satz: Struktursatz für endlich erzeugte Moduln über HIR} % (fold)
\label{sub:248}
Sei $D$ ein Hauptidealring, $M$ ein endlich erzeugter $D$-Modul. Dann ist $M$ die Summe von zyklischen Moduln $M= D {m_1} \oplus \ldots \oplus D {m_k}$ mit 
\[
	D \supsetneq \ann(m_1) \supsetneq \ann(m_2) \supsetneq \ldots \supsetneq \ann(m_k)
\]
\minisec{Beweis}
Schreibe $M \simeq \nicefrac{D^s}{L}$ mit $L \simeq D^t$, $t \le s$ (siehe \ref{sub:242}). Dann ist $L= Dl_1 \oplus \ldots \oplus D l_t$ mit 
$l_j = (a_{1j}, \ldots , a_{sj}) \in D^s$. Dann erhalten wir eine Matrix $A= (a_{ij}) \in D^{s \times t}$. Die Smith-Normalform
\[
	P \cdot A \cdot Q = \begin{pmatrix}
		a_1 & & & & & 0 \\
		& \ddots & & & & \\
		& & a_k & & & \\
		& & & 0 & & \\
		& & && \ddots & \\
		0 & & & & & 0
	\end{pmatrix}
\]
liefert ein neues Erzeugendensystem für $D^s$, $L$ nämlich $e_1, \ldots , e_s, f_1, \ldots , f_t$ mit 
\[
	D^s = \bigoplus_{i=1}^s D{e_i} \quad , \quad L= \bigoplus_{j=1}^t D f_j
\]
sowie $f_i = d_i \cdot e_i$ und 
\[
	\nicefrac{D^s}{L} \simeq \nicefrac{D}{(d_1)} \oplus \ldots \oplus \nicefrac{D}{(d_k)} \oplus D \oplus \ldots \oplus D
\]
Für $d_i \in D^\times$ Einheit ist $\nicefrac{D}{(d_i)} = 0$, also einfach weglassen.\bewende
% subsection 248 (end)

\subsection{Korollar: Struktursatz für endlich erzeugte abelsche Gruppen} % (fold)
\label{sub:249}
Ist $A$ eine endlich erzeugte abelsche Gruppe, dann ist $A$ die direkte Summe von endlich vielen zyklischen Gruppen.
\minisec{Beweis}
Dies folgt unmittelbar aus der Tatsache, dass jede abelsche Gruppe ein $\mathds{Z}$-Modul ist, vergleiche Bemerkung \ref{sub:23} (iv). \bewende
% subsection 249 (end)

\subsection[Definition und Satz: Torsionsmodul]{Definition und Satz} % (fold)
\label{sub:250}
Setze $\tor_D(M) := \set[m \in M]{\ann(m)\not= 0}$. Ist $D$ kommutativ und nullteilerfrei, dann ist $\tor_D(M)$ ein Untermodul, der \Index{Torsionsmodul} von $M$.
\minisec{Beweis}
Ist $a \cdot m = 0$, dann ist $d \cdot a \cdot m= a \cdot d \cdot m$ also $d \cdot m \in \tor(M)$ für alle $d \in D$. Sind $m_1, m_2 \in \tor(M)$, also
$a m_1 = b m_2 = 0$, dann ist $a b( m_1 + m_2)=0$. \bewende
% subsection 250 (end)

\subsection[Satz: Zerlegung eines endlich erzeugten Moduls in Torsionsmodul und freien Modul]{Satz} % (fold)
\label{sub:251}
Ist $M$ ein endlich erzeugter Modul über einem Hauptidealring $D$, dann ist 
\[
	M \simeq \tor_D(M) \oplus D^k
\]
und $\tor_D(M)$, $k$ sind eindeutig bestimmt.
\minisec{Beweis}
Nach Satz \ref{sub:248} ist 
\[
	M \simeq D m_1 \oplus \ldots \oplus D m_s \oplus D m_{s+1} \oplus \ldots \oplus D m_{s+k}
\]
mit 
\begin{align*}
	\ann(m_i) &\not= 0 \text{ für }i=1, \ldots ,s \\
	\ann(m_i) &= 0 \text{ für } i=s+1, \ldots , s+k
\end{align*}
Ein Element $d_1 m_1 + \ldots  + d_{s+k} m_{s+k} \in M$ ist ein Torsionselement
genau dann, wenn $d_{s+1} = \ldots = d_{s+k}=0$, also wenn es in $Dm_1\oplus \ldots \oplus D m_s$ liegt. Damit ist $\tor(M)= D m_1 \oplus \ldots \oplus D m_s$.
Es ist $\nicefrac{M}{\tor(M)} \simeq D^k$ und $k$ ist nach Satz \ref{sub:241} eindeutig bestimmt ($M$ ist halbeinfach).\bewende
% subsection 251 (end)

\subsection[Beispiele für Zerlegungen mit dem Torsionsmodul]{Beispiel} % (fold)
\label{sub:252}
\begin{enumerate}[(i)]
	\item Sei $D=Z$ und $M= \nicefrac{\mathds{R}}{\mathds{Z}}$. Dann ist 
	\(
		\tor_\mathds{Z}(\nicefrac{\mathds{R}}{\mathds{Z}}) = \nicefrac{\mathds{Q}}{\mathds{Z}}
	\), da $d \cdot r \in \mathds{Z} \Rightarrow r \in \mathds{Q}$. Also gilt
	\[
		\nicefrac{\mathds{R}}{\mathds{Z}} \simeq \nicefrac{\mathds{Q}}{\mathds{Z}} \oplus \bigoplus \mathds{Q}.
	\]
	Anmerkung\footnote{$\nicefrac{\mathds{R}}{\mathds{Z}}\simeq S^1 = \set[z \in \mathds{C}]{\abs*{z}=1}$ indem man $S^1$ in Polardarstellung, also nur als Menge der Winkel 
	betrachtet.}.
	\item Sei $D=\mathds{Z}$ und $M=\mathds{Q}$. Dann ist $\tor_\mathds{Z}(\mathds{Q})= 0$, aber $(\mathds{Q},+)$ ist \emph{kein} freier $\mathds{Z}$-Modul: Seien 
	$\frac{r}{s}, \frac{p}{q} \in \mathds{Q}$. Dann ist 
	\[
		r \cdot p \in \mathds{Z} \frac{r}{s} \cap \mathds{Z} \frac{p}{q} \not= 0 
	\]
	$\mathds{Q}$ ist nicht endlich erzeugt als $\mathds{Z}$-Modul.
\end{enumerate}
Ist $A$ eine endlich erzeugte abelsche Gruppe, dann existiert $k \in \mathds{Z}$ und eine endliche Gruppe $A_0$ mit $A \simeq A_0 \oplus \mathds{Z}^k$. Dabei heißt $k$ der
Rang von $A$. Zusammengefasst ergibt sich:
% subsection 252 (end)

\subsection[Satz: Umformulierung des Struktursatzes für endl. erzeugte abelsche Gruppen]{Satz} % (fold)
\label{sub:253}
Eine endlich erzeugte abelsche Gruppe ist direkte Summe einer endlichen Gruppe, der Torsionsgruppe, und einer freien abelschen Gruppe isomorph zu $\mathds{Z}^k$. Dabei ist 
der Rang $k$ invariant. Jede endliche abelsche Gruppe ist direkte Summe von zyklischen Gruppen.

Jede endliche abelsche Gruppe ist die Summe von zyklischen Gruppen von Primzahlordnung. Diese Ordnung und ihre Vielfachheiten sind eindeutig bestimmt. Zwei endlich erzeugte
abelsche Gruppen sind genau dann isomorph, wenn der Rang und diese Ordnungen übereinstimmten.
\minisec{Bemerkung}
Die Zerlegung ist hier anders als in Satz \ref{sub:248}. Die Zerlegung von $\tor(A)$ folgt aus den Sylow-Sätzen.
% subsection 253 (end)
% section 2 (end)
\newpage

\section{Tensorprodukte und Algebren} % (fold)
\label{sec:3}
Zuerst betrachten wir Tensorprodukte für Moduln über kommutativen Ringen $R$. Sind $U,V, W \in {_R \Mod}$, $f : U \times V \to W$ bilinear, dann suchen wir ein 
universelles Objekt $T$ und eine bilineare Abbildung $\varphi$, sodass es für jedes bilineare $f$ und jedes $W$ einen eindeutigen $\mathds{R}$"=Modulhomomorphismus $f'$ 
gibt, sodass das folgende Diagramm kommutiert.\index{Tensorprodukt}
\[
	\begin{tikzcd}
		U \times V \rar{\varphi} \drar{f} & T \dar[dashed]{f'}\\
		& W
	\end{tikzcd}
\]
\subsection[Satz: Universelle Eigenschaft des Tensorproduktes]{Satz} % (fold)
\label{sub:31}
Sei $R$ ein kommutativer Ring und $U,V \in {_R \Mod}$. Dann existiert ein bis auf Isomorphie eindeutiger $R$-Modul $U \otimes V$ und eine bilineare Abbildung 
$\varphi : U \times V \to U \otimes V$, die universell für alle bilinearen Abbildungen $f : U \times V \to W$ ist.
\minisec{Beweis}
Eindeutigkeit ist klar: Seien $T_1, \varphi_1$ und $T_2, \varphi_2$ zwei universelle Objekte. Dann betrachte
\[
	\begin{tikzcd}
		U \times V \rar{\varphi_1} \drar{\varphi_2} \arrow{ddr}{\varphi_1} \arrow{dddr}{\varphi_2}& T_1  \dar[dashed]{\varphi_2'} \arrow[bend left = 60]{dd}{\id}	\\
		& T_2 \dar[dashed]{\varphi_1'}  \arrow[bend left = 60, crossing over]{dd}{\id}\\
		& T_1 \dar[dashed]{\varphi_2'} \\
		& T_2
	\end{tikzcd}
\]
also $\varphi_2' \circ \varphi_1' = \id_{T_2}$, $\varphi_1' \circ \varphi_2' = \id_{T_1}$ $\Rightarrow T_1 \simeq T_2$.

Für die Existenz sei $\mathcal{F}_{U \times V}$ der freie $R$-Modul mit Basis (indiziert durch) $U \times V$ und sei $B \subseteq \mathcal{F}_{U \times V}$ der Untermodul,
der von allen Elementen der folgenden Form erzeugt wird
\begin{align*}
	&(u+u', v) - (u,v) - (u', v) \\
	&(u,v+v') - (u,v) - (u,v') \\
	&(\alpha \cdot u, v) - \alpha(u,v) \\
	&(u, \alpha \cdot v) - \alpha(u,v)
\end{align*}
für alle $u,u' \in U$, $v,v' \in V, \alpha \in R$. Sei $\varphi$ die Hintereinanderausführung
\[
	0 \longrightarrow U \times V \xhookrightarrow{\quad} \mathcal{F}_{U \times V} \longrightarrow \nicefrac{\mathcal{F}_{U \times V}}{B} \longrightarrow 0
\]
Offensichtlich ist $\varphi$ bilinear: Es ist $\varphi((u,v+ v')) = (u, v+ v')+ B $. Wegen $(u, v+v')= (u,v)+ (u, v')$ usw. in $\nicefrac{\mathcal{F}_{U\times V}}{B}$ ist
$\varphi$ bilinear. Setze $U \otimes V := \nicefrac{\mathcal{F}_{U \times V}}{B}$. Noch zu zeigen: $U \otimes V$, $\varphi$ erfüllen die universelle Eigenschaft:
Sei also $f : U \times V \to W$ bilinear. Dann lässt sich wegen der universellen Eigenschaft des freien Moduls $\mathcal{F}_{U \times V}$ $f$ fortsetzen zu 
$f_1 : \mathcal{F}_{U \times V} \to W$ (Satz \ref{sub:235})
\[
	\begin{tikzcd}
		U \times V \rar[hook]{\varphi'} \drar{f}& \mathcal{F}_{U \times V} \dar[dashed]{f_1}\\
		& W
	\end{tikzcd}
\]
Nach Definition von $B$ folgt $B \subseteq \ker f_1$, denn es ist zB.
\[
	f_1 \benbrace[\big]{(u+ u', v) - (u,v)- (u',v)} = f(u+u',v) - f(u,v) - f(u',v) = 0 
\]
da $f$ bilinear ist. Ebenso ist
\[
	f_1 \benbrace[\big]{(\alpha u, v) - \alpha(u,v)} = f(\alpha u,v) - \alpha f(u,v)  = 0 
\]
Daher erhalten wir $f' : U \otimes V = \nicefrac{\mathcal{F}_{U \times V}}{B} \to W$ und dieses $f'$ ist eindeutig bestimmt, da die Bilder von $(u,v)$ den Modul 
$U \otimes V$ erzeugen. \bewende
% subsection 31 (end)

\subsection[Bemerkung: Elemente und Erzeuger des Tensorprodukts, Modulstruktur des Tensorprodukts]{Bemerkung} % (fold)
\label{sub:32}
Das Bild von $(u,v)$ in $U \otimes V$ unter $\varphi$ wird mit $u \otimes v$ bezeichnet. D.h. $U \otimes V$ ist $R$-Modul mit Erzeugermenge 
$\set[u \otimes v]{u \in U, v \in V} $ und definiert Relationen
\begin{align*}
	(u+ u') \otimes v &= u \otimes v + u' \otimes v \\
	u \otimes (v+v') &= u \otimes v + u \otimes v' \\
	\alpha u \otimes v &=  \alpha \enbrace*{u \otimes v} = u \otimes \alpha v 
\end{align*}
für alle $u,u' \in U, v,v' \in V, \alpha \in R$. 
\minisec{Achtung}
Nicht jedes Element in $U \otimes V$ ist von der Form $u \otimes v$, $u \in U, v \in V$! Ein allgemeines Element ist von der Form $\sum_{i=1}^{k} u_i \otimes v_i $ für
$u_i \in U$, $v_i \in V$, $k \in \mathds{N}$.
\todo[noline]{RevChap3}
% subsection 32 (end)

\subsection[Proposition: Das Tensorprodukt als Operation auf $_R\Mod$]{Proposition} % (fold)
\label{sub:33}
Sei $R$ ein kommutativer Ring, $U,V,W$ seien $R$-Moduln. Dann gilt 
\begin{enumerate}[(i)]
	\item $U \otimes V \simeq V \otimes U$
	\item $U \otimes ( V \otimes W) \simeq (U \otimes V) \otimes W$
	\item $U \otimes( V \oplus W) \simeq (U \otimes V) \oplus (U \otimes W)$
\end{enumerate}
\minisec{Beweis}
\begin{enumerate}[(i)]
	\item $f : U \times V \to V \times U$, $(u,v) \mapsto (v,u)$ ist bilinear. Daher existiert ein Homomorphismus $\alpha : U \otimes V \to V \otimes U$, 
	$u \otimes v \mapsto v \otimes u$, d.h. 
	\[
		\alpha \enbrace*{ \sum u_i \otimes v_i} = \sum v_i \otimes u_i
	\]
	Ebenso ist $\beta : V \otimes U \to U \otimes V$, $v \otimes u \mapsto u \otimes v$ ein $R$-Modulhomomorphismus. Dann ist $\beta \circ  \alpha = \id_{U \otimes V}$ und
	$\alpha \circ  \beta = \id_{V \otimes U}$, d.h. $\alpha$ und $ \beta$ sind Isomorphismen.
	\item Betrachte $f : U \times V \times W  \to U \otimes (V \otimes W)$, $(u,v,w) \mapsto u \otimes (v \otimes w)$. Für festes $w \in W$ ist $f$ bilinear in $u, v$
	und wir erhalten $f_w' : (U \otimes  V) \to U \otimes (V \otimes W)$, $(u \otimes v)  \mapsto u \otimes (v \otimes w)$. Dann ist $f_w'$ $R$-linear und wir 
	erhalten eine bilineare Abbildung
	\[
		\tilde{f} : (U \otimes V) \times W \to U \otimes (V \otimes W) \quad , \quad \enbrace[\big]{(u \otimes v), w} \mapsto f_w(u \otimes v) =  u \otimes ( v \otimes w) 
	\]
	Daraus erhalten wir $f' : (U \otimes V) \otimes W \to U \otimes (V \otimes W)$. Entsprechend erhalten wir Inverses 
	$g' : U \otimes (V \otimes W) \to (U \otimes V) \otimes W $.
	\item Sei $\varphi: U \times (V \oplus W) \to (U \otimes V) \oplus (U \otimes W)$, $(u,v,w) \mapsto (u \otimes v, u \otimes w)$. Ist $f : U \times (V \oplus  W) \to Z$
	eine bilineare Abbildung in einen $R$-Modul $Z$, dann ist $f(u,v,w) = f(u,v)+ f(u,w)$ und dies kann auch als Abbildung von $(U \otimes V) \oplus (U \otimes W)$
	aufgefasst werden. Daher erfüllt $(U \otimes V) \oplus (U \otimes W)$ die erforderliche universelle Eigenschaft, d.h.
	\[
		U \otimes (V \oplus W) \simeq (U \otimes V) \oplus (U \otimes W) \bewende
	\]
\end{enumerate}
% subsection 33 (end)

\subsection[Beispiele zu Tensorprodukten]{Beispiel} % (fold)
\label{sub:34}
\begin{enumerate}[(i)]
	\item $\mathds{Q} \otimes_\mathds{Z} \mathds{R} \simeq \mathds{R}$ via $q \otimes r \mapsto q \cdot r$.\marginnote{Ich bin mir noch unsicher, ob die Rechnung stimmt … bzw verstehe nicht ganz, was sie zeigen soll…}
	\begin{align*}
		\sum_{i=1}^{k} \frac{m_i}{n_i} \otimes r_i = \sum \frac{\prod_{i=1}^k m_i}{n_i \prod_{j \not= i} m_j} \otimes r_i  = \sum \frac{1}{\prod n_i} \otimes
		\enbrace*{r_i m_i \prod_{j\not= i} n_j} = \frac{1}{\prod n_i} \otimes \enspace(\sum \ldots )   
	\end{align*}
	\item Seien $r,s \in \mathds{N}$, $(r,s)=1$ und $m r + n s =1$. Dann ist 
	\[
		\nicefrac{\mathds{Z}}{(r)} \otimes_\mathds{Z} \nicefrac{\mathds{Z}}{(s)} = 0,
	\]
	denn für jeden Erzeuger $a \otimes b$ gilt 
	\begin{align*}
		a \otimes b = (m r + n s) (a \otimes b) = m r (a \otimes b) + n s (a \otimes b) &= m( r a \otimes b) + n(a \otimes  s b) \\&= m(0 \otimes b) + n( a \otimes 0) = 0
	\end{align*}
	\item Ist $A$ eine abelsche Torsionsgruppe, dann ist $A \otimes_\mathds{Z} \mathds{Q} = 0$. Denn ist $a \in A$ mit $m a = 0$, dann ist 
	\[
		a \otimes  \frac{p}{q} = a \otimes \frac{p \cdot m}{q \cdot m} = a \cdot m \otimes \frac{p}{q \cdot m} = 0  
	\]
\end{enumerate}
% subsection 34 (end)

\subsection[Satz: Tensorprodukt eines $R$-Moduls und $R^n$]{Satz} % (fold)
\label{sub:35}
Sei $R$ ein kommutativer Ring, $U \in {_R \Mod}$. Dann ist $U \otimes R^n \simeq U^n= U \oplus \ldots \oplus U$. 
\minisec{Beweis}
Durch Induktion nach $n$: Wegen Proposition \ref{sub:33} (iii) genügt es, die Behauptung für $n=1$ zu zeigen, also $U \otimes_R R \simeq U \simeq R \otimes_R U$.

Sei $\theta : R \times U \to U$, $(r,u) \mapsto r \cdot u$. Dann ist $\theta$ bilinear. Ist $W \in {_R  \Mod}$ und $f : R \times U  \to W$ bilinear, dann ist 
$f(r,u)= f(1, r \cdot u)$, also ist mit $f' : u \mapsto f(1,u)$ gerade $f = \theta \circ f'$. Offensichtlich ist $f'$ eindeutig bestimmt, d.h. $U, \theta$ erfüllen
die universelle Eigenschaft aus Satz \ref{sub:31}, also ist $R \otimes U \simeq U$ wegen der Eindeutigkeit. \bewende
% subsection 35 (end)

\subsection[Korollar: Tensorprodukt von freien Moduln]{Korollar} % (fold)
\label{sub:36}
Sind $U,V$ freie $R$-Moduln, $R$ kommutativ, $U \simeq R^m, V \simeq R^n$. Dann ist $U \otimes V \simeq R^{n \cdot m}$.
\minisec{Bemerkung}
Für Vektorräume über Körpern gilt daher $\dim_K(U \otimes V)= \dim_K U \cdot \dim_K V$. Daher gilt $\set{e_1, \ldots , e_m}$, $ \set{f_1, \ldots , f_n}$ Basen für $U$ bzw. 
V, dann ist $\set[e_i \otimes f_j]{i=1,\ldots ,m , j=1,\ldots ,n} $ eine Basis für $U \otimes V$.
% subsection 36 (end)

\subsection[Korollar: Darstellung der Elemente eines Tensorprodukts mit freiem Modul]{Korollar} % (fold)
\label{sub:37}
Ist $U \in {_R \Mod}$, $V$ ein freier $R$-Modul, $R$ kommutativ. Dann hat jedes Element von $U \otimes V$ eine eindeutige Darstellung der Form $\sum u_i \otimes e_i$ mit
$u_i \in U$, $\set{e_i}_{i \in I}$ Basis für $V$.
\minisec{Vorsicht}
Ist $\sum u_i \otimes v_i = 0$ in $U \otimes V$ und die $v_i$ sind in $V$ linear unabhängig über $R$. Dann folgt \emph{nicht}, dass $u_i =0$ gilt. 
\marginnote{Siehe Beispiel \ref{sub:34} (iii)}

Sei $V' = \langle v_1, \ldots , v_k \rangle \subseteq V$. Ist $\sum u_i \otimes v_i = 0$ in $U \otimes V'$, dann folgt $u_i=0$. Die Einbettung $V' \hookrightarrow V$
induziert einen Homomorphismus 
\[
	U \otimes V' \to U \otimes V \tag*{$\star$}
\]
\uline{Achtung}: Dieser Homomorphismus muss nicht injektiv sein!
\minisec{Beispiel}
$2 \mathds{Z} \hookrightarrow \mathds{Z}$ ist injektiv, aber bleibt \emph{nicht} injektiv in
\[
	2 \mathds{Z} \otimes \nicefrac{\mathds{Z}}{2} \to \mathds{Z} \otimes \nicefrac{\mathds{Z}}{(2)} \simeq \nicefrac{\mathds{Z}}{(2)}
\]
Seien $e, f, f'$ Erzeuger der zyklischen Moduln $\nicefrac{\mathds{Z}}{(2)}, \mathds{Z}, 2 \mathds{Z}$. $2 \mathds{Z} \hookrightarrow \mathds{Z}$, $f' \mapsto 2 f$ und dann
\[
	f' \otimes  e \mapsto 2 f \otimes e = f \otimes 2 e  = 0
\]
Bemerkung: Ist $V'$ direkter Summand in $V$, dann bleibt nach Proposition \ref{sub:33} (iii) die induzierte Abbildung ($\star$) aber doch eine Einbettung. Über Körpern ist 
das immer der Fall.
% subsection 37 (end)

\subsection[Definition: Algebra]{Definition} % (fold)
\label{sub:38}
Sei $R$ ein kommutativer Ring. Eine $R$-Algebra ist ein Ring $A$, der gleichzeitig ein $R$-Modul ist, sodass die Multiplikation in $A$ $R$-bilinear ist, d.h. so dass gilt:
\[
	\alpha (x \cdot y) = x \cdot \alpha y = (\alpha x) \cdot y \quad \forall x,y \in A, r \in R
\]
Eine \Index{Unteralgebra} einer $R$-Algebra $A$ ist ein Unterring, der gleichzeitig ein $R$-Untermodul ist.
% subsection 38 (end)

\subsection[Beispiele für Algebren, Zentrum einer Algebra]{Beispiel} % (fold)
\label{sub:39}
\begin{enumerate}[(i)]
	\item Ist $A$ eine $R$-Algebra, dann ist das \Index{Zentrum} $C(A) = \set[z \in A]{x \cdot z = z \cdot x \enspace \forall x \in A} $ eine Unteralgebra. 
	$C(A)$ ist $R$-Untermodul: $z \in C(A)$, dann ist für $\alpha \in R, x \in A$
		\[
			x \cdot (\alpha z) = (\alpha x) \cdot z = z \cdot (\alpha x) = (\alpha z) \cdot x
		\]
	\item Jeder Ring $R$ ist auch $R$-Algebra: 
	\item Jeder Ring wird zu einer $\mathds{Z}$-Algebra durch $n x = x + \ldots + x$, $(-n) x = - (n x)$
\end{enumerate}
Ein Homomorphismus von $R$-Algebren ist ein Ringhomomorphismus, der gleichzeitig $R$-linear ist.
\minisec{Bemerkung}
Für $y=1$ in Definition \ref{sub:38} ergibt sich $x \cdot  \alpha 1 = \alpha x$ und für $x=1$ ergibt sich $\alpha 1 \cdot y = \alpha \cdot y$. Also gilt 
$\alpha 1  \cdot x = x \cdot  \alpha 1 = \alpha x$ für alle $x \in A, \alpha \in \mathds{R}$. Für $\alpha = \beta \cdot 1$ erhalten wir 
\[
	\alpha 1 \cdot \beta \cdot 1 = \alpha (\beta 1) = \alpha \beta \cdot 1,
\]
d.h. die Abbildung $R \to A$, $\alpha \mapsto \alpha \cdot 1$ ist ein Algebrahomomorphismus in das Zentrum von $A$, d.h. $R \cdot 1 \subseteq C(A)$.

Umgekehrt gilt: Ist $R$ ein Ring, $f : R \to C(A)$ ein Ringhomomorphismus. Dann wir $A$ zu einer $R$-Algebra durch $\alpha \cdot x = f(x) \cdot x$ für $x\in A,\alpha\in R$.
Mit anderen Worten: Eine $R$-Algebra ist ein Ring $A$ zusammen mit einem Homomorphismus $f : R \to C(A)$. Manchmal werden Algebren allgemeiner definiert: Eine $R$-Algebra 
ist ein $R$-Modul $A$ mit einer bilinearen Multiplikation (nicht notwendig assoziativ oder mit Eins). Bei uns sind alle Algebren assoziativ und \Index{unital}, d.h. mit 1. 
% subsection 39 (end)

\subsection[Weitere Beispiele für Algebren]{Beispiel} % (fold)
\label{sub:310}
Sei $R$ ein kommutativer Ring.
\begin{enumerate}[(i)]
	\item $U \in {_R \Mod}$, $A = \End_R(U)$. Dann ist $A$ ein $R$-Modul: Für $f \in A, \alpha \in R$, $x \in U$ ist
	\[
		(\alpha f) (x) = f (\alpha x)
	\]
	Damit wir $A$ zu einer $R$-Algebra. Ist $V \simeq R^n$, dann ist $A \simeq M_n(R)$ und $C(A) \simeq R$, nämlich $\alpha \cdot \ind$, $\alpha \in R$.
	\item $B_n(R)=$ Menge der oberen Dreiecksmatrizen ist Unteralgebra von $M_n(R)$.
	\item Sei $M= \set[u_i]{i \in I}$ ein Monoid (= Halbgruppe mit 1) und sei $A= \mathcal{F}_M$ der freie Modul über $M$ mit Multiplikation $u_i \cdot u_j = u_k$ wie in 
	$M$. Durch lineares Fortsetzen erhalten wir eine Multiplikation auf $A$, die $A$ zu einer $R$-Algebra macht. Dann ist $A$ assoziativ und unital, weil $M$ diese 
	Eigenschaften hat. Diese Algebra wird mit $RM$ bezeichnet.
	
	Häufig: $M=G$ eine Gruppe, $R=K$ ein Körper: Dann heißt $K G$ die \Index{Gruppenalgebra} von $G$ über $K$. \uline{Spezialfälle:}
	\begin{enumerate}[(a)]
		\item $M= \set{1, X, X^2, X^3, \ldots }$, $RM= R[X]$.
		\item $G=Z = \langle X \rangle$. Dann ist $KG= K[X, X ^{-1}]$ der Ring der Laurent-Polynome in $X$.
	\end{enumerate}
	\item Endlich dimensionale Algebra über einem Körper $K$: Ist $A$ $K$-Algebra mit $\set{u_1, \ldots , u_n}$, dann ist die Struktur von $A$ eindeutig bestimmt durch die
	$n^3$ Strukturkonstanten $c_{ijk}$, $i,j,k = 1, \ldots ,n$ mit $u_i \cdot u_j = \sum_{k=1}^{n} c_{i,j,k} \cdot u_k$.
	
	Ist $u_1 = 1_A$, dann ist $c_{1ir} = c_{i1r} = \delta_{ir}$ und die Assoziativität von $A$ ist äquivalent zu
	\[
		\sum_{i} c_{i,j,k} = \sum_{i,l,m} c_{lik} \cdot c_{mji}
	\]
	Die Sätze über einfache Artinsche Ringe(Schur, Wedderburn, \ldots) gelten insbesondere für endlich dimensionale einfache $K$-Algebren.
	\item Ist $U \in {_R \Mod}$, dann sei $T(U) = \bigoplus_{n \in \mathds{N}} U^{\otimes n}$, mit $U^{\otimes n} = U \otimes  \ldots \otimes  U$, $U^{\otimes  0} = R$.
	Dann heißt $T(U)$ die \Index{Tensoralgebra} von $U$. 
	
	Ist $U=R$, dann ist $T(U) \simeq R[X]$.
\end{enumerate}
% subsection 310 (end)

\subsection[Definition: Darstellung einer $R$-Algebra]{Definition} % (fold)
\label{sub:311}
Eine \Index{Darstellung} der $R$-Algebra $A$ ist (gegeben durch) einen $R$"=Algebrenhomomorphismus $A \to \End_R(U)$ für einen $R$-Modul $U$.
\marginnote{vgl. $G \to \Sym(X)$ in Def. \ref{sub:11}} 
Für jede Algebra $A$ bezeichnet man mit $\rho$ die rechts-reguläre Darstellung $A$ auf sich selbst, d.h. 
\begin{align*}
	\rho : A \to \End_R(A), \quad a \mapsto \rho_a : x \mapsto x \cdot a 
\end{align*}
Diese Darstellung ist treu, d.h. $\ker \rho = 0$. Ist $A$ endlich dimensionale Algebra über einem Körper $K$ mit Basis $\set{u_1, \ldots , u_n}$, dann ist 
$\End_K(A) \simeq M_n(K)$. Für $a = \sum \alpha_i u_i \in A$ wird dann $\rho_a$ beschrieben durch die Matrix $(\rho_a)_{ij}$ mit 
\[
	(\rho_a)_{ij} = \sum_{k=1}^{n} \alpha_k \cdot c_{ijk}
\]
mit $c_{ijk}$ wie in \ref{sub:310} (iv).
% subsection 311 (end)

\subsection[Korollar: Jede $n$-dimensionale $K$-Algebra ist isomorph zu Unteralgebra von $M_n(K)$]{Korollar} % (fold)
\label{sub:312}
Jede $n$-dimensionale $K$-Algebra ist isomorph zu einer Unteralgebra von $M_n(K)$.
\minisec{Beweis}
$A \simeq \im \rho \subseteq \End_K(A) \simeq M_n(K)$. \bewende
% subsection 312 (end)

\subsection[Satz: Tensorprodukt von Algebren ist eine Algebra]{Satz} % (fold)
\label{sub:313}
Ist $R$ ein kommutativer Ring, dann ist für $R$-Algebren $A,B$ auch $A \otimes B$ eine $R$-Algebra mit $\mu (A \otimes B) \times (A \otimes B) \to A \otimes B$,
$(a \otimes b)(a' \otimes b') = a a' \otimes  b b'$ und linearer Fortsetzung. Die Tensoralgebra  $A \otimes B$ ist genau dann kommutativ, assoziativ bzw. unital, wenn $A$
und $B$ kommutativ, assoziativ und unital sind.
\minisec{Beweis}
Es genügt nachzurechnen, dass diese Multiplikation bilinear ist und die Eigenschaften  von $A,B$ sich übertragen. \bewende
% subsection 313 (end)

\subsection[Beispiel: Erweiterung der Skalare]{Beispiel} % (fold)
\label{sub:314}
\begin{enumerate}[(i)]
	\item Sei $K$ ein Körper, $E \supseteq K$ eine Körpererweiterung und eine $K$-Algebra $A$ gegeben. Dann ist $A \otimes E$ eine $E$-Algebra mit 
	$\dim_E(A \otimes E)= \dim_K(A)$ nach Beispiel \ref{sub:39}, denn $1 \otimes E \subseteq C(A \otimes_K E)$ und für eine $K$-Basis $\set{u_1, \ldots , u_n}$ von $A$ ist 
	$\set{u_1 \otimes 1, \ldots , u_n \otimes 1} $ eine $E$-Basis für $E =: A_E$. Man sagt: $A_E$ entsteht aus $A$ durch Erweiterung der Skalare.
	\item Sei $A= M_n(R)$. Dann hat $A$ eine $R$-Basis $\set[e_{ij}]{i,j \le n}$ mit Multiplikation $e_{ij} \cdot e_{kl}= \delta_{jk} e_{il}$ und $\sum e_{ii} = 1$.
	Ist $B$ eine $R$-Algebra, dann ist $A \otimes B$ ein freier $B$-(Rechts-)Modul mit derselben Basis als $B$-Modul. d.h. $M_n(R) \otimes B \simeq M_n(B)$. Für $B=M_k(R)$
	ist also $M_n(R) \otimes M_k(R) \simeq M_{n \cdot k}(R)$.
\end{enumerate}
Im Folgenden untersuchen wir endlich dimensionale Algebren über einem kommutativen Körper $K$.
% subsection 314 (end)

\subsection[Definition: Zentrale Algebra]{Definition} % (fold)
\label{sub:315}
Eine $K$-Algebra $A$ heißt \bet{zentral}\index{Algebra!zentral}, wenn $C(A)= K \cdot 1_A$. Ist $A$ eine endlich dimensionale einfache $K$-Algebra, dann ist nach Wedderburn (\ref{sub:229}) $A \simeq M_n(D)$ für einen Schiefkörper $D$, der endliche Dimension über $K$ hat. Weiter ist 
\[
	C(A) \simeq C(D) =: E \supseteq K.
\]
Dann ist $E$ ein kommutativer Körper und wir können $A$ als $E$-Algebra auffassen. Also solche ist $A$ dann eine zentral einfache endlich dimensionale $E$-Algebra.
\minisec{Bemerkung}
Sind $A,B$ $K$-Algebren, dann enthält $A \otimes B$ Unteralgebren $1 \otimes B \simeq B$ und $A \otimes 1 \simeq A$. Wir geben eine interne Charakterisierung von
Tensorprodukten an (vgl. \ref{sub:129}). 
% subsection 315 (end)

\subsection[Definition: Linear disjunkt]{Definition} % (fold)
\label{sub:316}
Sei $C$ eine $K$-Algebra, $U,V \subseteq C$ $K$-Untervektorräume. Dann heißen $U$ und $V$ \Index{linear disjunkt} über $K$, wenn für linear unabhängige $u_i \in U$, 
$v_j \in V$ die Elemente $u_i v_j \in C$ linear unabhängig sind. Mit anderen Worten ist die natürliche Abbildung $U  \otimes V \to C$, $u \otimes v \mapsto u v$ injektiv.
% subsection 316 (end)

\subsection[Proposition: Isomorphie einer $K$-Algebra zum Tensorprodukt von Unteralgebren]{Proposition} % (fold)
\label{sub:317}
Sei $C$ eine $K$-Algebra, $A,B \subseteq C$ Unteralgebren, dann ist $C \simeq A \otimes B$, falls gilt
\begin{enumerate}[(i)]
	\item $A$ und $B$ sind linear disjunkt über $K$
	\item $A B = C$
	\item $a b = b a$ für alle $a \in A, b \in B$. {\small \color{light_gray} $(a \otimes 1)(1 \otimes b)= (1 \otimes b)(a	\otimes 1)= a \otimes b$}
\end{enumerate}
\minisec{Beweis}
Die Abbildung $A \times B \to C$, $(x,y) \mapsto xy$ ist bilinear und induziert daher eine Abbildung $A \otimes B \to C$. Diese Abbildung ist injektiv nach (i), surjektiv
nach (ii) und ein Homomorphismus nach (iii). \bewende
% subsection 317 (end)

\subsection[Definition: Einhüllende Algebra]{Definition} % (fold)
\label{sub:318}
Für eine $K$-Algebra $A$ heißt $A^e = A \otimes_K A^\op$ die \Index{einhüllende Algebra} von $A$. Sei $A \subseteq B$ eine Unteralgebra. Dann wird $B$ in 
natürlicher Weise zu einem $A^e$-Modul, nämlich durch 
\[
	\enbrace*{\sum a_i \otimes b_i} y = \sum a_i \cdot y \cdot b_i
\]
für $a_i \in A, b_i \in A^\op, y \in B$. Aus der universellen Eigenschaft des Tensorprodukts folgt, dass dies eine wohldefinierte Modulwirkung beschreibt. 
Insbesondere wird $A$ selber zum $A^e$-Modul. Die $A^e$-Untermoduln von $A$ sind genau die beidseitigen Ideale von $A$. Nach Lemma \ref{sub:225} ist 
$\End_R(R,+) \simeq R^\op$, d.h. $R$-lineare Endomorphismen von $R \in \,{_R\Mod}$ sind genau die Rechtsmultiplikationen $\rho_r$ für $r \in R$. Daher ist 
$\End_{A^e}(A)$ die Menge der Abbildungen, die sowohl Rechts- als auch Links-Multiplikation sind, d.h. Abbildungen $x \mapsto c \cdot x = x \cdot d$. Für $x=1$ erhält
man $c=d$, d.h. $c \in C(A)$. Damit folgt $\End_{A^e}(A) \simeq C(A)$.
% subsection 318 (end)

\subsection{Satz (Dichtheitssatz für halbeinfache Moduln)} % (fold)
\label{sub:319}
Sei $R$ ein Ring, $M \in {_R\Mod}$ halbeinfach, $R' = \End_R (M) = C_{\End (M,+)} (\Lambda_R)$ mit $\Lambda_R = \set[\lambda_r]{r \in R} $, $R'' = \End_{R'} (M)$ \marginnote{vertauschen} 
(für $M$ aufgefasst als $R'$-Modul). Sei $\set{x_1, \ldots , x_n} \subseteq M$,
$a'' \in R''$, dann existiert $a \in R$ mit $a x_i = a'' x_i$ für $i=1, \ldots ,n$.
\minisec{Beweis}
Zuerst überlegen wir, dass jeder $R$-Untermodul $N \subseteq M$ auch ein $R''$-Untermodul ist. Schreibe $M= N \oplus P$ und $\pi : M = P \oplus N \to N$. Dann ist
$\pi \in R'$ und $N = \pi(M)$. Daher ist für $\varphi \in R''$ also 
\[
	\varphi(N)= \varphi \pi (M) = \pi  \varphi(M) \subseteq N.
\]
D.h. $N$ ist $R''$-Untermodul von $M$. Ist nun $n=1$, dann ist $N= R x_i$ ein $R$-Untermoduln und daher auch $R''$-Untermodul. Wegen $x_1 \in N$ ist für $a'' \in R''$
auch $a'' x_1 \in N= R x_1$. Also existiert $a \in R$ mit $a'' x_1 = a x_1$. Für $n\ge 1$ benutzen wir folgenden Trick:
% subsection 319 (end)

\subsection{Lemma} % (fold)
\label{sub:320}
Sei $M \in {_R\Mod}$, $R' = \End_R(M)$, $R'' = \End_{R'}(M)$. Dann ist für $\alpha \in R''$ die Abbildung $\varphi_\alpha : M^n  \to M^n$, 
$(u_1, \ldots ,u_n)  \mapsto \enbrace*{\alpha(u_1), \ldots , \alpha(n_n)} $ ein Endomorphismus von $M^n$, ist also $\End_R(M^n)$-Modul.
\minisec{Beweis}
Nach Lemma \ref{sub:227} ist $\End_R(M^n) \simeq M_n(R')$. Offensichtlich ist $M^n$ auch $\End_R(M^n)$-Modul. Für die Behauptung ist also nur zu zeigen, dass für jedes 
$\alpha \in R''$ die Abbildung $\varphi_\alpha$ mit jedem $R$-Endomorphismus von $M^n$ vertauscht. Dies ist der Fall, da $\alpha$ gerade mit jedem $R$-Endomorphismus von 
$M$ vertauscht. \bewende
\minisec{Beweis von \ref{sub:319}}
Mit $M$ ist auch $M^n$ halbeinfach. Nach Lemma \ref{sub:320} ist für $\alpha'' \in R''$ und $\varphi_{a''}$ wie eben auch $\varphi_{\alpha''} \in \End_{\End_R(M^n)} (M^n)$.
Nun wenden wir den Fall $n=1$ auf den $R$-Modul $M^n$ und $\varphi_{\alpha''}$. Damit finden wir für $(x_1, \ldots , x_n) \in M^n$ ein $\alpha \in R$ mit 
\[
	\varphi_{\alpha''}(x_1, \ldots , x_n)= \enbrace[\big]{\alpha''(x_1), \ldots , \alpha''(x_n)} = \alpha (x_1, \ldots , x_n) = (\alpha x_1, \ldots , \alpha x_n).
\]
% subsection 320 (end)

\subsection[Satz: Modulisomorphie der einhüllenden Algebra zum Matrizenmodul]{Satz} % (fold)
\label{sub:321}
Ist $A$ eine endlich dimensionale zentral einfache $K$-Algebra, dann ist $A^e = A \otimes A^\op \simeq M_n(K)$ mit $n= \dim A$.
\minisec{Beweis}
Betrachte $A$ als $A^e$-Modul. Dann ist $A$ einfach und $\End_{A^e} (A) \simeq K$. Da $A$ endlich dimensional über $K$ ist, folgt aus dem Dichtheitssatz (\ref{sub:319}),
dass $A^e \to \End_K(A)$ surjektiv ist. Da 
\[
	\dim_K A^e = \dim_K \End_K(A) = n^2
\]
folgt $A^e \simeq \End_K(A)$ und wegen $\End_K(A) \simeq M_n(K)$ folgt die Behauptung. \bewende
% subsection 321 (end)

\subsection{Satz} % (fold)
\label{sub:322}
Sei $B$ eine $K$-Algebra. $A$ eine endlich dimensionale zentral einfache Unteralgebra. Dann gilt $B \simeq A \otimes_K C$, wobei $C= C_B(A)$ der Zentralisator von $A$ in 
$B$ ist, d.h. $C= \set[x \in B]{x a = a x \enspace \forall a \in A}$. 

Die Abbildung $I \mapsto A \otimes I$ ist Verbandsisomorphismus zwischen den Idealen von $B$ und 
den Idealen von $C$ und es ist $C(B)= C(C)$
\minisec{Beweis}
Betrachte $B$ als $A^e$-Modul. Da $A^e$ nach Satz \ref{sub:321} einfach ist, $A^e = M_n(K)$, ist nach Wedderburn (\ref{sub:231}) 
\(
	B = \bigoplus A^e x_i
\)
direkte Summe von einfachen $A^e$-Moduln, die alle isomorph zu $A$ sind. Als $A^e$-Modul wird $A$ von $1_A$ erzeugt und es ist 
\[
	(a \otimes 1)1 = a \cdot 1 = 1 \cdot a = (1 \otimes a) 1
\]
und falls $(a \otimes 1)1 = 0$, dann ist $a=0$. Da alle $A^e x_i \simeq A$ als $A^e$-Moduln, existiert ein $c_i \in A^e x_i$ mit $(a \otimes 1)c_i = (1 \otimes a)c_i$
und $(a \otimes 1)c_i =0$ $\Rightarrow a=0$. Daher ist $B= \bigoplus Ac_i$ mit $a c_i = c_i a$ für alle $a \in A$ und $a c_i = 0 \Rightarrow a=0$. Daher sind alle 
$c_i \in C= C_B(A)$ und jedes Element in $B$ hat \emph{eindeutige} Darstellung der Form $\sum a_i c_i$ mit $a_i \in A$.

Ist $c \in C$, $c = \sum a_i c_i$, dann folgt wegen $a c = c a$ für alle $a \in A$ auch $a a_i = a_i a$ für alle $a \in A$, d.h. $a_i \in C(A) \simeq K$. Daher ist
$C = \sum K c_i$, d.h. $\set[c_i]{i \in I}$ ist Basis für $C$. Nach Proposition \ref{sub:317} ist $B = A \otimes_K C$. 

Für den zweiten Teil sei $I \unlhd C$ ein Ideal. Dann ist $A \otimes I$ ein Ideal in $B = A \otimes C$ und es ist $(A \otimes I) \cap (1 \otimes C) = I$. Ist 
$(x_1=1, x_2, \ldots , x_n)$ Basis für $A$, dann hat jedes $b \in B$ eine \emph{eindeutige} Darstellung der Form $b= \sum c_i \otimes x_i = \sum c_i x_i$ mit
$c_i \in C$. Die Elemente von $A \otimes I$ sind von der Form $\sum d_i x_i$ mit $d_i \in I$. Daher folgt aus $b= \sum d_i x_i \in C$ schon $b= c_1 x_1$ mit $c_1 \in C$,
d.h. $C \cap A I = \set{d_1 x_1 = d_1 \in I} = I$. Daher ist $I \mapsto A \otimes I$ injektiv.
Ist $I'$ ein Ideal in $B$, dann ist $I'$ ein $A^e$-Untermodul von $B$, d.h. für $I = I' \cap C$ ist dann $I' = \sum A d_j$ mit $d_j \in I$. Daher ist $I' = A \cdot I_i$, 
d.h. die Abbildung $I \mapsto A \otimes I$ ist surjektiv.

Noch zu zeigen: $C(B) = C(C)$. "$\subseteq$"{} ist klar. Ist $c \in C(C)$, dann vertauscht $c$ mit allen $a \in A$ und mit allen $d \in C$, also $c \in C(B)$. \bewende
% subsection 322 (end)

\subsection{Korollar} % (fold)
\label{sub:323}
Ist $A$ eine endlich dimensionale zentral einfache $K$-Algebra, $C$ eine beliebige $K$-Algebra, dann ist die Abbildung $I \mapsto A \otimes I$ ein Verbandsisomorphismus 
zwischen Idealen von $B = A \otimes C$ und den Idealen von $C$ und es ist $C(B) = C(C)$.
\minisec{Beweis}
Identifiziere $A,C$ mit $A \otimes I$ bzw. $1 \otimes C$ in $B= A \otimes C$. Behauptung: $C = C_B(A)$. Klar "$\subseteq $"{}. "$\supseteq$"{} wie eben. Daher folgt die
Behauptung aus \ref{sub:322}. \bewende
% subsection 323 (end)

\subsection{Korollar} % (fold)
\label{sub:324}
Ist $A$ eine endlich dimensionale zentral einfache $K$-Algebra, $C$ eine beliebige $K$-Algebra, dann ist $A \otimes C$ einfach genau dann, wenn $C$ einfach ist und 
$A \otimes C$ ist zentral genau dann, wenn $C$ zentral ist.
Insbesondere gilt: Ist $A$ endlich dimensionale zentral einfache $K$-Algebra, $E \supseteq K$ eine Körpererweiterung, dann ist $A_E = A \otimes E$ endlich dimensionale
zentral einfache $E$-Algebra.
% subsection 324 (end)

\subsection{Korollar} % (fold)
\label{sub:325}
Sind $A_1, \ldots , A_n$ endlich dimensionale zentral einfache $K$-Algebren, dann ist auch $A_1 \otimes \ldots \otimes A_n$ zentral einfach und endlich dimensional.
% subsection 325 (end)

\subsection[Definition: Zerfällungskörper über zerfallende zentral einfache $K$-Algebra]{Definition} % (fold)
\label{sub:326}
$M_n(K)$ heißt zerfallende zentral einfache $K$-Algebra. Ist $A$ endlich dimensionale zentral einfache $K$-Algebra, dann heißt eine Körpererweiterung $E \supseteq K$
ein \Index{Zerfällungskörper} für $A$, falls $A_E = A \otimes_K E \simeq M_n(E)$ für ein $n \in \mathds{N}$. Man sagt $A$ zerfällt über $E$. 
% subsection 326 (end)

\subsection{Bemerkung} % (fold)
\label{sub:327}
\begin{enumerate}[(i)]
	\item Ist $E$ Zerfällungskörper für $A$, dann auch jedes $E' \supseteq E $. Insbesondere ist jedes $E \supseteq K$ Zerfällungskörper für $M_n(K)$. 
	\item Sind $A,B$ zentral einfach und $E$ Zerfällungskörper für $A,B$, dann auch für $A \otimes B$: Es ist 
	\[
		(A \otimes B)_E = (A \otimes B) \otimes E \simeq (A \otimes_k B) \otimes_K E \simeq (A \otimes E) \otimes (B \otimes E).
	\]
	$M_{n k}(E) \simeq M_n(E) \otimes M_k(E)$.
	\item Ist $E$ Zerfällungskörper für $A$ ,dann auch für $A^\op$: 
	\[
		A \otimes E \otimes E \otimes A^\op = M_n(E) = M_r(E) \otimes (E \otimes A^\op) = M_r(E \otimes A^\op)
	\]
	Daher ist $E \otimes A^\op \simeq M_s(E)$ mit $r \cdot s = n$.
	\item Ist $A \simeq M_r(B)$, $B$ zentral einfach über $K$, dann ist $E \supseteq K$ Zerfällungskörper von $A$ genau dann, wenn $E$ Zerfällungskörper von $B$ ist. Es ist
	\[
		A \otimes E \simeq M_r(B) \otimes E \simeq M_r(K) \otimes B \otimes E \simeq M_r(B_E).
	\]
\end{enumerate}
% subsection 327 (end)

\subsection{Proposition} % (fold)
\label{sub:328}
Sei $K$ algebraisch abgeschlossen, $A$ eine endlich dimensionale $K$-Algebra. Ist $A$ nullteilerfrei, dann ist $A=K$. Insbesondere existieren keine echten endlich 
dimensionalen Schiefkörper über einem algebraisch abgeschlossenen Körper $K$.
\minisec{Beweis}
Sei $a \in A$. Da $A$ endliche $K$-Dimension hat, sind die Potenzen von $a$ linear abhängig über $K$. Sei etwa $f(a) = a^n c_1  + a^{n-1} c_2 + \ldots + c_n = 0$ mit
$c_i \in K$ und $n$ minimal. Da $K$ algebraisch abgeschlossen ist, hat $f(X) \in K[X]$ eine Nullstelle $\lambda  \in K$. Daher ist $f(X)= (X- \lambda )\cdot g(X)$ mit 
$\deg(g) = n-1$. Nach Vorraussetzung ist aber $g(a) \not= 0$, wegen $n$ minimal. Weil $A$ nullteilerfrei ist, folgt $a= \lambda  \in K$. \bewende
% subsection 328 (end)

\subsection{Korollar} % (fold)
\label{sub:329}
Sei $K$ algebraisch abgeschlossen, $A$ zentral einfache, endlich dimensionale Algebra über $K$. Dann ist $A \simeq M_n(K)$ für ein $n \in \mathds{N}$.
\minisec{Beweis}
$A \simeq M_n(D)$ für einen endlich dimensionalen Schiefkörper über $K$. Dann ist $D=K$ nach \ref{sub:328}. \bewende
\minisec{Bemerkung}
Insbesondere zerfällt jede zentral einfache, endlich dimensionale $K$-Algebra über dem algebraischen Abschluss von $K$. Ziel: Zeige Existenz endlicher Körpererweiterungen
als Zerfällungskörper. Zeige: Jeder maximale kommutative Teilkörper von $D$ ist Zerfällungskörper für $A \simeq M_N(D)$.
% subsection 319 (end)

\subsection{Satz} % (fold)
\label{sub:330}
Ist $D$ ein Schiefkörper endlicher Dimension über $K$. Dann ist eine endliche Körpererweiterung $\nicefrac{E}{K}$ ein Zerfällungskörper für $D$ genau dann, wenn $E$ 
Teilkörper einer Algebra $A= M_r(D)$ mit $C_A(E)=E$ ist.
\minisec{Beweis}
\begin{description}
	\item["$\Leftarrow$":] Sei $E \subseteq M_r(D) = A$ mit $C_A(E)=E$. Dann ist $A = \End(V)$, wobei $V$ ein $r$-dimensionaler $D^\op$-Vektorraum ist. Dann ist
	$V$ auch $D^\op \otimes _K E$-Modul, da 
	\[
		(d \otimes e)x = d \cdot e \cdot x = e \cdot d \cdot x 
	\]
	für $d \in D^\op$, $ e \in E \subseteq A$, $x \in V$. Da $D^\op \otimes E$ einfach ist, ist die Darstellung von $D^\op \otimes E$ auf $V$ treu,
	\marginnote{treu, da Kern sonst Unteralgebra} d.h. wir können $D^\op \otimes E$ mit dem entsprechenden Endomorphismenring identifizieren. Da 
	$D^\op \otimes E$ endlich dimensional und einfach über $K$ ist, ist nach den Sätzen von Wedderburn $V$ als $D^\op \otimes E$ als Modul halbeinfach und
	daher ist der Dichtheitssatz anwendbar:
	
	Es ist $A = \End_{D^\op}(V)$, also ist $\End_{D^\op \otimes E} (V) = C_A(E) = E$ \marginnote{zusätzlich auch mit $E$ vertauschen}. 
	Also können wir $V$ auch als endlich dimensionalen $E$-Vektorraum 
	betrachten. Nach dem Dichtheitssatz ist $D^\op \otimes E = \End_E(V) \simeq M_n(E)$ für $n = \dim_E V$. Daher ist $E$ Zerfällungskörper von $D^\op$ und
	damit von $D$.
	\item["$\Rightarrow$":] Ist umgekehrt $D \otimes E \simeq M_n(E)$, dann auch $D^\op \otimes E \simeq M_n(E)$. Ist $V$ ein irreduzibler 
	$D^\op \otimes E$-Modul, dann auch ein $M_n(E)$-Modul, d.h. $V$ ist $n$-dimensionaler $E$-Vektorraum und $\End_E(V) = D^\op \otimes E$. Aber $V$ ist 
	auch ein endlich dimensionaler $D^\op$-Vektorraum also etwa $r = \dim_{D^\op}V$. Dann ist $E \subseteq \End_{D^\op}(V)$. Daher ist 
	\[
		E \subseteq C_{\End_{D^\op}(V)}(E) \subseteq C_{\End_E(V)} \enbrace*{D^\op \otimes E} = E.
	\]
	Daher ist $E=C_{M_r(D)}(E)$. \bewende
\end{description}
% subsection 330 (end)

\subsection{Korollar} % (fold)
\label{sub:331}
Ist $A$ eine zentral einfache $K$-Algebra mit $A \simeq M_r(D)$, $D$ Schiefkörper. Dann ist jeder maximale kommutative Teilkörper $E$ von $D$ Zerfällungskörper für $A$.
\minisec{Beweis}
Klar: Es existiert ein maximaler kommutativer Teilkörper $E$, denn $\dim_K D$ ist endlich. Es reicht zu zeigen: $E$ ist Zerfällungskörper für $D$. (Wende Satz 
\ref{sub:330} an mit $A=D$). Sei $E' = C_D(E) \supseteq E$. Ist $E \subsetneq E'$, wähle $c \in E' \setminus E$. Dann ist der von $E$ und $c$ erzeugte Unterkörper von $D$
auch kommutativ. \light Maximalität von $E$. Also gilt $C_D(E)=E$ und die Behauptung folgt aus \ref{sub:330}. \bewende
\minisec{Beispiel}
$\mathbb{H}$ reelle Quaternionen. 
% subsection 331 (end)

\subsection{Korollar} % (fold)
\label{sub:332}
Ist $A$ zentral einfache $K$-Algebra endlicher Dimension, dann ist $\dim_K A = n^2$ für ein $n \in \mathds{N}$ der \Index{Grad} von $A$.
\minisec{Beweis}
Es ist $A \otimes E \simeq M_n(E)$ für einen Zerfällungskörper $E$ und $\dim_K A = \dim_E \enbrace[\big]{M_n(E)} = n^2 $. \bewende
% subsection 332 (end)

\subsection{Satz (Skolem-Noether)} % (fold)
\label{sub:333}
Ist $B$ eine endlich dimensionale, zentral einfache $K$-Algebra, $A \subseteq B$ einfache Unteralgebra, dann lässt sich jeder Monomorphismus $f : A \to B$ zu einem inneren
Automorphismus\footnote{Automormphismus durch Konjugation} fortsetzen, d.h. es existiert ein $b \in B$ mit $f(a) = b ^{-1} \cdot a \cdot b$.
\minisec{Beweis}
$E= A \otimes_K B^\op$ ist einfach und endlich dimensional. Nach Wedderburn ist jeder $E$-Modul halbeinfach und alle einfachen $E$-Moduln sind isomorph. Daher sind
auch endlich dimensionale $E$-Moduln derselben Dimension isomorph. Betrachte $B$ als $E$-Modul auf zwei Weisen:
\[
	\enbrace*{\sum a_i \otimes b_i} x = \sum a_i \cdot x \cdot b_i \qquad \qquad \enbrace*{\sum a_i \otimes b_i} x = \sum f(a_i) \cdot x \cdot b_i  
\]
Dies definiert eine $E$-Modulstruktur auf $B$, weil $f$ ein Homomorphismus ist. Nach Vorüberlegung sind diese beiden $E$-Moduln isomorph, d.h. es existiert ein $E$-linearer
Isomorphismus $\varphi : B \to B$ mit 
\[
	\varphi \enbrace*{\sum f(a_i) \otimes b_i} = \sum a_i \varphi(x) b_i 
\]
für $a_i \in A$, $b_i, x \in B$. Insbesondere ist dann $\varphi(x b) = \varphi(x) b$ für alle $x,b \in B$. Daher ist $\varphi(x) = \varphi(1) x = d x$ für ein 
invertierbares $d \in B$, denn $\varphi$ ist invertierbar. Wegen $\varphi \enbrace*{f(a) x} = a \cdot \varphi(x) $ für alle $a \in A, x \in B$ folgt 
$A d x = d f(a) x$ und für $x=1$ folgt $f(a) = d ^{-1} a d$, d.h. $f$ lässt sich zum inneren Automorphismus $\kappa_d : x \mapsto d ^{-1} x d$ fortsetzen. \bewende
% subsection 333 (end)

\subsection{Korollar} % (fold)
\label{sub:334}
Jeder Automorphismus einer endlich dimensionalen, zentral einfachen Algebra ist ein innerer.
\minisec{Beweis}
Setze $A=C(B) = K$ in Satz \ref{sub:333}. \bewende
% subsection 334 (end)

\subsection{Korollar} % (fold)
\label{sub:335}
\todo{Dieses Korollar verlangt noch nach etwas Erläuterung}
In einer endlich dimensionalen, zentral einfachen Algebra sind einfache Unteralgebren derselben Dimension konjugiert (und ebenso ihre Zentralisatoren).
% subsection 335 (end)

\subsection{Satz (Brauer)} % (fold)
\label{sub:336}
Sei $A$ endlich dimensional, zentral einfache $K$-Algebra, $B \subseteq A$ eine einfache Unteralgebra mit $C(B)= E \supseteq K$. Dann gilt: $B' = C_A(B)$ ist einfach mit
Zentrum $E$ und $B'' = C_A(B') = C_A\enbrace[\big]{C_A(B)} = B$. Weiter ist $C_A(E) = B \otimes B'$ und $[A : B'] = [B :K]$. Für $r := [B:K]$ ist 
\[
	A \otimes B^\op \simeq B' \otimes  M_r(K) \simeq M_r(B')
\]
\minisec{Beweis}
Sei $r= \dim_K B$ und betrachte $B$ als $M_r(K)$-Modul. Dann enthält $M_r(K)$ die Unteralgebra $\rho_B = B^\op$ und $\lambda_B  \simeq B$ als Bilder der rechts- bzw. 
linksreguläre Darstellung. Dann ist $A \otimes_K M_r(K)$ zentral einfach über $K$ und enthält Unteralgebren
\[
	B \otimes_K 1  \simeq 1 \otimes \lambda_B \simeq 1 \otimes B \simeq B.
\]
Diese Unteralgebren sind einfach und daher nach Skolem-Noether (\ref{sub:333}) in $A \otimes M_r(K)$ konjugiert mit isomorphen Zentralisatoren. Also gilt:  
\marginnote{Rechtswirkung und Linkswirkung zentralisieren sich}
\[
	B' \otimes M_r(K) \simeq A \otimes \rho_B \simeq A \otimes B^\op \simeq M_r(B')
\]
Damit folgt die letzte Behauptung. Die Dimension über $B'$ ergeben
\[
	r^2 = \dim_{B'} M_r(B') = [A:B'] \cdot [B:K],
\]
also $[A:B']=r$. Nun ist nach \ref{sub:324} $A \otimes B^\op$ einfach\marginnote{beide Teile einfach} und wieder nach \ref{sub:324} ist dann auch $B'$ einfach.
Es ist $B \subseteq B''$, $B''' = C_A(B'') = B'$ und wie eben (mit $B$ ersetzt durch $B'$) folgt $[B'' : K]=r$, d.h. $B''=B$. Sei nun $F= C(B')$, dann ist $F \supseteq E$
und wegen $B''=B$ folgt $E \supseteq F$, d.h. $E=F$. Daher sind $B$ und $B'$ zentral einfache $E$-Algebren. $B,B' \subseteq C_A(E)$ und aus Satz \ref{sub:322} folgt
$C_A(E)= B \otimes B'$. \bewende
\minisec{Bemerkung}
Falls $C(B) \supsetneq K$, dann ist $C_A(E) \subsetneq A$, denn $A$ ist zentral einfach.
% subsection 336 (end)

\subsection{Korollar} % (fold)
\label{sub:337}
Ist $A$ zentral einfach mit endlicher Dimension über $K$, $E \subseteq A$ ein Teilkörper mit $K \subseteq E$ und $B \subseteq C_A(E)$. Dann ist
\[
	A \otimes E \simeq B \otimes M_r(K)
\]
mit $r = \dim_K E$. Ist $n= \dim_K A$, dann ist $r^2 \mid n$ und $A \otimes B^\op \simeq M_{n/r}(E)$.
\minisec{Beweis}
Setze $B = E$ im Satz \ref{sub:336}, dann ist $[C_A(E) :A] = \frac{n}{r^2}  = [B :K]$. Mit $n = [A : K]= r^2 \cdot [B:K]$ gilt $C_A(E) \simeq E \otimes B$. Folglich ist
$[C_A(E) :E] = [B:K]$. \bewende
% subsection 337 (end)

\subsection{Korollar} % (fold)
\label{sub:338}
$A$ endlich dimensionale, zentral einfache $K$-Algebra, $E \subseteq A$ Teilkörper mit $K \subseteq E$. Dann sind äquivalent:
\begin{enumerate}[(i)]
	\item $E' = C_A(E) = E$
	\item $E$ ist ein maximaler kommutativer Unterring von $A$.
	\item $[A:K] = [E:K]^2$
	\item $[A:K] = [A:E]^2$
\end{enumerate}
Insbesondere gilt für maximale Unterkörper in einem Schiefkörper $D$ also 
\(
	[D:K] = [E:K]^2 = [D:E]^2
\)
\minisec{Beweis}
\begin{description}
	\item[(i)$\Leftrightarrow $(ii):] Klar.
	\item[(iii)+(iv)$\Leftrightarrow$(i)] Es ist $E' \supseteq E$ und 
	\[
		[A:K] = [A:E] \cdot [E:K]= [E:K] = [E:K] \cdot [E':K].
	\]
	Also ist $[A:E]^2 \ge [A:K] \ge [E:K]^2$ und Gleichheit gilt genau dann, wenn $E=E'$, d.h. (iii)+(iv) sind äquivalent zu (i).
\end{description}
\minisec{Bemerkung}
Unterkörper $E \subseteq A$ mit $\dim_K A  = \enbrace*{\dim_E A}^2 $ muss es nicht geben. Z.B. $M_r(K)$, $K$ algebraisch abgeschlossen, $r \ge 2$.
% subsection 338 (end)
Zum Satz von Wedderburn über Nicht-Existenz von endlichen Schiefkörpern:
\subsection{Lemma} % (fold)
\label{sub:339}
Sei $G$ eine endliche Gruppe, $H \lneq G$. Dann ist $G$ \emph{nicht} die Vereinigung der Konjugierten von $H$.
\minisec{Beweis}
Sei $\abs*{G} = \abs*{H} \cdot n$, also wenn $h= \abs*{H}$, dann ist $\nicefrac{G}{H} = \set{a_1 H , \ldots , a_n H}$. Dann ist für $b \in G$
\[
	b H b ^{-1} = \underbrace{b H}_{a_i H} \underbrace{H ^{-1} b ^{-1}}_{H ^{-1} a_i ^{-1}} = a_i \underbrace{H H ^{-1}}_{=H} a_i ^{-1} = a_i H a_i ^{-1}
\]
Daher existieren höchstens $n$ viele Konjugierte von $H$, die alle $1_G$ enthalten. Damit ist 
\[
	\abs*{\bigcup_{i=1}^n a_i H a_i ^{-1}}  \le n \cdot (h-1) + 1 < n \cdot h = \abs*{G} \bewende
\]
% subsection 339 (end)

\subsection[Satz (Wedderburn): Endliche Schiefkörper sind kommutativ]{Satz (Wedderburn)} % (fold)
\label{sub:340}
Jeder endliche Schiefkörper ist kommutativ.
\minisec{Beweis}
Sei $D$ ein endlicher Schiefkörper, $K= C(D)$, d.h. $D$ ist zentral einfach, endlich dimensional über $K$. Sei $E \subseteq D$ ein maximaler kommutativer Teilkörper.
Nach Korollar \ref{sub:338} (ii) haben alle maximalen Unterkörper denselben Grad $r$ über $K$, die maximalen Teilkörper sind nach Korollar \ref{sub:335} alle konjugiert.
Da jedes Element von $a \in D$ über $K$ einen kommutativen Teilkörper erzeugt, liegt jedes $a \in D$ in einem Konjugierten von $E$. nach Lemma \ref{sub:339} (angewandt auf 
die multiplikative Gruppe von $D$) folgt also $E=D=K$. \bewende
% subsection 340 (end)

\subsection[Definition und Satz: Die Brauergruppe]{Definition und Satz} % (fold)
\label{sub:341}
Zwei zentral einfache, endlich dimensionale $K$-Algebren $A$ und $B$ heißen \Index{ähnlich} (oder Brauer-äquivalent, $A \sim B$), falls $A\simeq M_n(D)$, $B\simeq M_k(D)$ 
für einen Schiefkörper $D$ über $K$, $n,k \in \mathds{N}$. Dabei sind $D,n,k$ nach Satz \ref{sub:229} eindeutig bestimmt. Insbesondere gibt es zu jeder zentral einfachen 
Algebra $A$ einen Schiefkörper $D$ über $K$ mit $A \sim D$. Offensichtlich ist $\sim$ eine Äquivalenzrelation. 

Für jeden Körper $K$ bildet die Menge der Äquivalenzklassen bezüglich des
Tensorprodukts eine Gruppe, die \Index{Brauergruppe} $\Br(K)$ von $K$. Dabei ist $[K]$ das Neutralelement (denn $K \sim M_n(K)$) und $A \otimes K \simeq A$ und für
$[A] \in \Br(K)$ ist $A \otimes A^\op \simeq M_n(K)$, d.h. $[A^\op] = [A]^{-1}$.
\minisec{Beweis}
Wir hatten bereits gesehen, dass das Tensorprodukt kommutativ und assoziativ ist, daher ist nur zu zeigen, dass
\(
	[A] [B] = [A \otimes B]
\)
wohldefiniert ist, d.h. sind 
\[
	M_n(D) \otimes M_k(D') \sim M_m(D) \otimes M_l(D') \sim D \otimes D'
\]
\minisec{Beispiel}
\begin{enumerate}[(i)]
	\item Ist $K$ algebraisch abgeschlossen oder endlich, dann ist $\Br(K)=1$.
	\item Ist $K=\mathds{R}$, dann folgt aus $[\mathds{C} : \mathds{R}]=2$ und $\mathds{C}$ algebraisch abgeschlossen nach Korollar \ref{sub:340} \todo{Nummerierung prüfen},
	dass ein Schiefkörper über $\mathds{R}$ Grad 2 (also Dimension 4) hat. Das heißt ist $D$ endlicher Schiefkörper über $\mathds{R}$, dann ist $\dim_\mathds{R}D=\set{1,4}$.
	Jeder echte Schiefkörper über $\mathds{R}$ ist isomorph zu den Quaternionen.
\end{enumerate}
% subsection 341 (end)
% section 3 (end)
\newpage

\section{Darstellungstheorie endlicher Gruppen} % (fold)
\label{sec:4}

\subsection[Definition: Darstellung einer Gruppe]{Definition} % (fold)
\label{sub:41}
Eine \Index{Darstellung} einer Gruppe ist (gegeben durch) einen Homomorphismus $\rho : G \to GL(V)$ für einen (endlich dimensionalen) $K$-Vektorraum. Der \Index{Grad} der 
Darstellung ist $\dim V$. Ist $B = \set{u_1, \ldots , u_n}$ eine $K$-Basis für $V$, dann ergibt jede Darstellung $\rho : G \to \GL(V)$ eine Matrix-Darstellung von $G$,
indem $\rho(g) \in \GL(V)$ durch die zugehörigen Matrix bezüglich $B$ beschrieben wird. \\
Ist $h : G \to \Sym(n)$ ein Homomorphismus, dann erhält man die zugehörige
\Index{Permutationsdarstellung} durch $\rho(g) u_i = u_{h(g)(i)}$. Dies hat man insbesondere für $\abs*{G}=n$ und 
\[
	h : G \to \Sym(G) \quad , \quad  g \mapsto \lambda_g = [x \mapsto g x].
\]
% subsection 41 (end)

\subsection{Beispiel} % (fold)
\label{sub:42}
\begin{enumerate}[(i)]
	\item Ist $G= \langle g \rangle$ zyklisch, $\abs*{G}=n$, dann erhält man für $g=(1, \ldots ,n) \in \Sym(n)$ also $\rho(g)u_i = u_{i+1}$, $\rho(g)u_n= u_1$ mit
	zugehöriger Matrixdarstellung 
	\[
		\rho : g \longmapsto \begin{pmatrix}
			0 & & & 1 \\
			1 & \ddots & 0 & \\
			 & \ddots & \ddots &\\
			 0 & & 1 & 0
		\end{pmatrix}
	\]
	Ziel ist es, die Struktur aller Darstellungen zu verstehen. Jede Darstellung $\rho : G \to \GL(V)$ hat eine eindeutige Fortsetzung auf die Gruppenalgebra:
	\[
		\tilde{\rho} : K[G] \to \End(V) \simeq M_n(K)
	\]
	Erinnerung: $K[G]$ ist die $K$-Algebra mit Basis $G = \set{g_1, \ldots, g_n}$ mit Multiplikation 
	\[
		\enbrace*{\sum_{g \in G} \alpha_g g} \enbrace*{\sum_{h \in G} \beta_h h} = \sum_{g,h} \alpha_g \beta_h (g \cdot h) \qquad, \quad \alpha_g, \beta_h \in K
	\]
	Umgekehrt induziert jede Darstellung der Gruppenalgebra $K[G] \to \End(V)$ eine Darstellung der Gruppe $G \to \GL(V)$. Durch eine
	Darstellung $\rho : K[G] \to \End_K(V)$ wird $V$ zu einem $K[G]$-Modul durch 
	\[
		\enbrace*{\sum_{g \in G} \alpha_g g} \cdot v= \rho \enbrace*{\sum_g \alpha_g g}(v).
	\]
	Das heißt endlich dimensionale Darstellungen einer Gruppe $G$ sind \enquote{das gleiche} wie $K[G]$-Moduln, die als $K$-Vektorräume endlich dimensional sind.
	
	
	\textbf{Übersetzung der Begriffe:} Ist $\rho : G \to \GL{V}$ eine Darstellung, dann ist $U \le V$ ein $K[G]$-Modul von $V$ genau dann, wenn $U$ ein $\rho(G)$-invarianter
	Untervektorraum von $V$ ist, d.h. wenn $\rho(g)U \subseteq U$ für alle $g \in G$. Ist $U$ ein $K[G]$-Untermodul, dann erhalten wir induzierte Darstellungen, nämlich 
	$\rho\big|_{U} : G \to \GL(U)$ und die Quotientendarstellung 
	\[
		\rho_{\nicefrac{V}{U}} : G \to \GL{\nicefrac{V}{U}} \quad ,\quad p_{\nicefrac{V}{U}} (g) = \benbrace*{x+ U \mapsto \rho(g)(x)+U}
	\]
	Eine Darstellung $\rho : G \to \GL(V)$ heißt irreduzibel (bzw. vollständig reduzibel oder halbeinfach), falls $V$ als $K[G]$-Modul irreduzibel (oder einfach) bzw.
	vollständig reduzibel (oder auch halbeinfach) ist. \footnote{$p \in \End_k(V)$ zentral $G$, d.h. $\forall v \in V, g \in G : \varphi(g v)= g (\varphi(v)) \iff \varphi 
	\in \End_{K[G]}(V)$}
	
	\textbf{Erinnerung:} $M$ ist halbeinfacher $R$-Modul $\iff M= \bigoplus M_i$, $M_i$ einfacher $R$-Modul $\iff$ jeder $R$-Untermodul $U \subseteq M$ hat Komplement
	$U^\bot$ in $M$, d.h. $M= U \oplus U^\bot$.
\end{enumerate}
% subsection 42 (end)

\subsection[Definition: Äquivalenz von Darstellungen]{Definition} % (fold)
\label{sub:43}
Zwei Darstellungen $\rho_i : G \to \GL(V_i)$, $i=1,2$ heißen \Index{äquivalent}, falls ein Vektorraumisomorphismus $\varphi : V_1 \to V_2$ existiert mit 
$\rho_2 = \varphi \rho_1 \varphi ^{-1}$. Dann sind offensichtlich $V_1$ und $V_2$ als $K[G]$-Moduln isomorph. 
% subsection 43 (end)

\subsection[Definition: Spurabbildung]{Definition} % (fold)
\label{sub:44}
Sei $G$ eine endliche Gruppe, $V$ ein $K[G]$-Modul (oder $G$-Modul). Die \Index{Spurabbildung} ist definiert durch 
\[
	\Tr_G : V \to V^G = \set[v \in V]{g v = v \enspace\forall g \in G}  \quad , \qquad  x \mapsto \sum_{g \in G} g x.
\]
Ist $f \colon V \to W$ ein $K$-Homomorphismus von $K[G]$-Moduln, dann setze $\Tr_G(f) \colon x \mapsto \sum_{g \in G} g f g ^{-1} x$. Dann ist $\Tr_G(f): V \to W$ ein Homomorphismus
von $K[G]$-Moduln\footnote{$\varphi : V \to W$ ist $K[G]$-Homomorphismus: $\varphi(g \cdot v) = g \cdot \varphi(v)$}.
% subsection 44 (end)

\subsection{Bemerkung} % (fold)
\label{sub:45}
Ist $f : V \to W$ bereits ein Homomorphismus von $G$-Moduln, dann ist $\Tr_G(f)= \abs*{G} \cdot f$.
% subsection 45 (end)

\subsection{Satz (Maschke)} % (fold)
\label{sub:46}
Sei $G$ eine endliche Gruppe, $K$ ein Körper. Dann ist $K[G]$ genau dann halbeinfach, wenn 
\[
	\Char(K) \nmid \abs*{G}.
\]
\minisec{Beweis}
\begin{description}
	\item["$\Leftarrow$":] Sie $n= \abs*{G}$, $\Char(K) \nmid \abs*{G}$. Dann existiert $\frac{1}{n} \in K$. Sei $V$ ein $K[G]$-Modul, $U \subseteq V$ ein $K[G]$-Untermodul.
	Wir müssen zeigen, dass $U$ ein $G$-invariantes Komplement hat. Sei $U'$ ein beliebiges Vektorraum-Komplement von $U$ in $V$, d.h. $V= U \oplus U'$ als Vektorraum (aber
	$U'$ ist vielleicht nicht $G$-invariant). Betrachte nun die Vektorraum-Projektion $\pi_U : V \to U$ auf $U$. Definiere $\varphi : V \to U$ durch
	\[
		\varphi = \frac{1}{n} \Tr_G(\pi) .
	\]
	Wegen $\pi\big|_U = \id_U$, gilt $\varphi\big|_{U} = \frac{1}{n} \Tr(\id_U) = \id_U$ nach Lemma \ref{sub:45}. Nun ist $\varphi$ ein $K[G]$-Modulhomomorphismus, d.h.
	$\ker \varphi = U^\bot$ ist ein $K[G]$-Untermodul von $V$ und daher gilt
	\[
		V = U \oplus  U^\bot
	\]
	als $K[G]$-Modul.
	\item["$\Rightarrow$":] Sei nun $\Char K \mid \abs*{G}$, $\abs*{G} = n$. Betrachte die reguläre Darstellung von $G$ auf $K[G]=V$. Dann ist $v = \sum_{g \in G} g \in V^G$
	und $V^G \le V$ ist ein $G$-invarianter Untervektorraum von $V$, also $K[G]$-Untermodul von $V$, $0 \lneq V^G \lneq V$. Ist nun $\pi : V \to V^G$ ein beliebiger 
	$K[G]$-Homomorphismus, betrachte $v= \sum g \in V^G$
	\[
		\pi(v) = \sum_g  \pi(g) = \sum_g g  \pi(1) = n \pi(1) = 0
	\]
	Daher kann es keine $K[G]$-lineare Projektion von $V$ auf $V^G$ geben und daher hat $V^G$ kein Komplement in $V$ als $K[G]$-Modul. \bewende 
\end{description}
Ist nun $G$ eine Gruppe, $\rho : G \to \GL(V)$ halbeinfache Darstellung von $G$ auf einem $K$-Vektorraum $V$ (d.h. $V$ ist als $K[G]$-Modul halbeinfach), dann zerfällt $
V$ in die direkte Summe von einfachen $K[G]$-Moduln als 
\[
	V = V_1^{(1)} \oplus \ldots \oplus V_1^{(m_1)} \oplus V_2^{(1)} \oplus \ldots \oplus V_2^{(m_2)} \oplus  \ldots  \oplus V_r^{(1)} \oplus \ldots \oplus V_r^{(m_r)}
\]
mit $V_i^{k} \simeq V_j^{l}$ als $K[G]$-Moduln genau dann, wenn $i=j$. Mit $\rho_i = \rho \big|_{V_i}$ schreibt man auch
\[
	\rho \sim m_1 \rho_1 \oplus m_2 \rho_2 \oplus \ldots \oplus m_r \rho_r.
\]
Dabei heißen die $m_i$ die Vielfachheiten der irreduziblen Komponenten $\rho_i$.
\[
	W_i = \bigoplus_{j=1}^{m_i} V_i^{(r)}
\]
heißt homogenene Komponente von $V$. Ist $G$ eine Gruppe, $H \unlhd G$ ein Normalteiler, $\sigma : H \to \GL(V)$ eine Darstellung, dann ist für $g \in G$ die 
Abbildung $\sigma^g : H \to \GL(V)$, $h \mapsto [v \mapsto \sigma(g h g ^{-1}) v]$ wieder eine Darstellung von $H$ auf $V$ ($g h g ^{-1} \in H$), die zu $\sigma$ 
konjugiert ist. Es ist $\sigma$ irreduzibel genau dann wenn $\sigma^g$ irreduzibel ist.
% subsection 46 (end)

\subsection{Satz (Clifford)} % (fold)
\label{sub:47}
Sei $G$ eine endliche Gruppe, $H \unlhd G$ ein Normalteiler, $\rho : G \to \GL(V)$ eine irreduzible Darstellung von $G$ auf $V$. Dann ist $\rho_H : H \to\GL(V)$ halbeinfach 
und alle homogenen irreduziblen Komponenten von $V$ als $K[H]$-Modul sind konjugiert und haben die gleiche Multiplizität.
\minisec{Beweis}
Sei $U$ ein irreduzibler $K[H]$-Untermodul von $V$ und betrachte $\sum_{g \in G} \rho(g) U$. Dieser Raum ist $G$-invariant, weil $V$ als $K[G]$-Modul irreduzibel ist, also 
$= V$. Seien nun $g \in G$, $h \in H$, $y \in U$. Dann folgt 
\[
	\rho(h) \underbrace{\enbrace[\big]{\rho(g)y}}_{\in \rho(g) U} = \rho(h g) y = \rho(g h^g) y = \rho(g) \underbrace{\rho(h^g)y}_{\in U} \in \rho(g) U
\]
Daher ist $\rho(g) U$ ebenfalls $K[H]$-Untermodul. Behauptung: Für alle $g \in G$ ist $\rho(g)U$ irreduzibel $K[H]$-Modul. $U$ ist als $H$-Modul irreduzibel und die Wirkung
von $H$ auf $\rho(g)U$ ist gegeben durch die Konjugationswirkung von $H$ auf einem zu $U$ isomorphen Vektorraum. Nach der Vorbemerkung sind also alle $\rho(g) U$ 
irreduzible  $H$-Moduln, d.h. ist als $H$-Modul halbeinfach. Da die irreduziblen Komponenten von der Form $\rho(g) U$ für ein $g \in G$ sind, haben alle dieselbe Dimension
und sind konjugiert.
Die $G$-Wirkung bildet isomorphe $H$-Untermoduln wieder auf solche ab. Daher permutiert $G$ die $H$-homogenen Komponenten transitiv und alle homogenen Komponenten haben 
dieselbe Multiplizität. \bewende
% subsection 47 (end)

\subsection{Korollar} % (fold)
\label{sub:48}
Ist $G$ eine endliche Gruppe, $H \unlhd G$, $\rho : G \to \GL(V)$ halbeinfach, dann ist auch $\rho \big|_H$ halbeinfach.
\minisec{Beweis}
Wende \ref{sub:47} auf die irreduziblen Komponenten an. \bewende
% subsection 48 (end)

\subsection[Definition: Doppelzentralisatoreigenschaft]{Definition} % (fold)
\label{sub:49}
Eine $K$-Algebra $A$ hat die \Index{Doppelzentralisatoreigenschaft}, falls für einen $A$-Modul $V$ und $A' = \End_A V$, $A'' = \End_{A'} V$ gilt $A= A''$.
% subsection 49 (end)

\subsection{Satz} % (fold)
\label{sub:410}
Ist $\rho : G \to \GL(V)$ eine halbeinfache Darstellung, dann hat $A =K[G]$ die Doppelzentralisatoreigenschaft.
\minisec{Beweis}
Dichtheitssatz von Jacobson \ref{sub:310}. \bewende \medskip \\
% subsection 410 (end)
Sei nun $G$ eine endliche Gruppe, $K$ ein Körper, $\Char K \nmid \abs*{G}$, $A = K[G]$. Es gibt eine Zerlegung $A = \bigoplus_{i=1}^s A_i$ in einfache Algebren. Nach
Wedderburn ist $A_i \simeq M_{n_i}(D_i)$ für einen Schiefkörper $D_i$ über $K$, $i=1, \ldots , s$, $d_i = \dim_K (D_i)$.
\subsection{Satz} % (fold)
\label{sub:411}
Sein $G$, $K$, $A_i$ wie eben, dann entsprechen die minimalen Linksideale $J_i \subseteq A_i$ genau den irreduziblen Darstellungen von $G$.  Der Grad der zu $J_i$ gehörigen 
Darstellung ist $n_i \cdot d_i$. $J_i = M_{n_i}(D_i)e_{kk}$, $A_i = \bigoplus_{k=1}^{n_i} M_{n_i}(D_i)e_{kk}$
\minisec{Beweis}
Für eine halbeinfache Algebra $A$ entsprechen die irreduziblen $A$-Moduln genau den minimalen Linksidealen der $A_i$. Ein minimales Linksideal in $M_n(D)$ ist von der 
Form $M_n(D)e_{ii}$ mit $K$-Dimension $n \cdot \dim_K D$. 

Ist $\rho_i$ die zu $J_i = A_i$ gehörige irreduzible Darstellung, dann ist für eine beliebige Darstellung 
$\rho : G \to \GL(V)$ (für $\Char K \nmid \abs*{G}$) also $V$ ein halbeinfacher $A(=K[G])$-Modul und zerfällt in einfache Komponenten 
$V= \bigoplus V_j$, $V_j$ $G$-invariant. Für jedes $j$ existiert ein $i$ mit $\rho \big|_{V_j} = \rho_i$ und wir können schreiben 
\[
	\rho \sim m_1 \rho_1 \oplus \ldots \oplus m_s \rho_s \bewende
\]
% subsection 411 (end)

\subsection{Satz} % (fold)
\label{sub:412}
Ist $\rho : G \to K[G]$ die reguläre Darstellung, dann ist die Vielfachheit von $\rho_i$ in $\rho$ genau $n_i$ in der Notation von vorher, also
\[
	\rho \sim n_1 \rho_1 \oplus \ldots \oplus n_s \rho_s
\]
\minisec{Beweis}
Die reguläre Darstellung ist die Linkswirkung von $G$ auf $A=K[G]$, d.h. $K[G]$ betrachtet als $K[G]$-Modul. Es ist $A= \bigoplus_{i=1}^s A_i$, 
\[
	A_i = M_{n_i}(D_i) = \bigoplus_{k=1}^{n_i} M_{n_i}(D_i) e_{kk}
\]
Daraus folgt die Behauptung. \bewende
% subsection 412 (end)

\subsection{Satz} % (fold)
\label{sub:413}
Sei $G$ eine endliche Gruppe, $G= C_1 \cup \ldots \cup C_r$ die Zerlegung in die Konjugiertenklassen.\marginnote{$Z(A)= \bigoplus_{i=1}^s Z(A_i)$}
Wir setzen $c_i = \sum_{h \in C_i}h \in A$, $i \le r$. Dann bilden 
$c_1, \ldots , c_r$ eine Basis für $Z(A)$.
\begin{description}
	\item[1. Behauptung: $c_i \in Z(A)$:] Für $g \in G$ gilt 
	\[
		g ^{-1} c_i g = \sum_{h \in c_i} g ^{-1} h g = \sum h = c_i \in Z(A)
	\]
	\item[2. Behauptung: $\set{c_1,\ldots ,r_r} $ linear unabhängig:] Klar.
	\item[3. Behauptung: $\langle \set{c_1, \ldots ,c_r}\rangle=Z(A)$:] Sei nun $c \in Z(A)$ beliebig, etwa $c = \sum_{g \in G} \gamma_g g$, $\gamma_g \in K$, $g \in G$.
	Sei $h \in G$, also 
	\begin{align*}
		h ^{-1}c h = \sum_g \gamma_g h ^{-1} g h = \sum_g \gamma_{h g h ^{-1}} g = \sum_g \gamma_g g =c
	\end{align*}
	Ein Koeffizientenvergleich ergibt $\gamma_g = \gamma_{h g h ^{-1}}$ für alle $g,h \in G$, d.h. es existiert $\gamma_i \in K$, so dass
	für alle $g \in C_i$ gilt $\gamma_g= \gamma_i$. Daraus folgt die Behauptung. \bewende
\end{description}
% subsection 413 (end)

\subsection{Bemerkung} % (fold)
\label{sub:414}
Daher gilt $s \le \dim Z(A)=r$. Es gilt also $\dim Z(A) = \sum_{i=1}^s [K_i :K]=r$ und daher $r=s$ genau dann, wenn $A_i$ zentral einfach für alle $i=1, \ldots ,s$.
Ist $K$ algebraisch abgeschlossen, dann ist immer $r=s$.
% subsection 414 (end)

\subsection{Satz} % (fold)
\label{sub:415}
Sei $G$ eine endliche Gruppe, $K$ algebraisch abgeschlossen, $\Char K \nmid \abs*{G}$, $s = \#$Konjugiertenklassen von $G$. Dann ist die Anzahl der irreduziblen 
Darstellungen von $G$ (bis auf Äquivalenz) genau $s$ und falls die $\rho_1, \ldots \rho_s$ wie eben sind mit $\deg \rho_i = n_i$, dann folgt 
\[
	\abs*{G} = \sum_{i=1}^s  n_i^2.
\]
\minisec{Beweis}
Die erste Aussage ist genau die Vorbemerkung \ref{sub:414}. Ist $K$ algebraisch abgeschlossen, dann ist $A= K[G] = \bigoplus_{i=1}^s M_{n_i}(K)$, daher gilt 
$\abs*{G} =\dim A = \sum_{i=1}^s n_i^2$. \bewende
% subsection 415 (end)

\subsection{Korollar} % (fold)
\label{sub:416}
Ist $G$ abelsch, dann ist über algebraisch abgeschlossenen Körpern $K$ mit $\Char K \nmid \abs*{G}$ jede irreduzible Darstellung $1$-dimensional, d.h. jede Darstellung
$G \to \GL(V)$, für $V$ $K$-Vektorraum, ist diagonalisierbar.
\minisec{Beweis}
Es ist $s= \abs*{G}$, also $n_i=1$ für $i=1,\ldots ,s$. \bewende
% subsection 416 (end)

\subsection{Beispiel} % (fold)
\label{sub:417}
Sei $n \in \mathds{N}$, $G=C_n$ zyklisch der Ordnung $n$, $A=\mathds{C}[G]$. Da $G$ abelsch ist, gibt es genau $n$ irreduzible Darstellungen, jeweils von Grad $1$. Die
entsprechenden Darstellungen sind:
\[
	g \mapsto \enbrace*{e^{\frac{2 \pi i r}{n}}} \qquad , \qquad r=1, \ldots ,n, \enspace\langle g \rangle = C_n
\]
\minisec{Erinnerung}
Ist $V$ ein $K$-Vektorraum, $L \supseteq K$ eine Körpererweiterung, dann ist $V_L = V \otimes_K L$ ein $L$-Vektorraum mit $\dim_L V_L = \dim_K V$. Jede $K$-lineare Abbildung
von $V$ lässt sich kanonisch zu einer $L$-linearen Abbildung von $V_L$ fortsetzen. Ds gilt insbesondere für Darstellungen $\rho : G \to \GL(V)$ und entsprechend
$\rho_L : G \to \GL(V_L)$.
% subsection 417 (end)

\subsection{Satz} % (fold)
\label{sub:418}
Sei $G$ eine Gruppe, $K$ ein Körper, $L \supseteq K$ eine Körpererweiterung, $V$ ein $K$-Vektorraum. Dann gilt 
\[
	\End_{L[G]} V_L = L \cdot \End_K[G] V \simeq \enbrace*{\End_{K[G]} V}_L 
\]
\minisec{Beweis}
Es ist 
\begin{align*}
	\End_{L[G]} V_L = \End_{L[G]} \enbrace*{V \otimes_K L} = \End_{K[G]} V \otimes_K L = \enbrace*{\End_{K[G]} V}_L  
\end{align*}
Klar ist auch $L \cdot \End_{K[G]} V \subseteq \End_{L[G]} V_L$. Für die andere Inklusion wähle eine Basis $\set[\lambda_i]{i \in I}$ für $L$ über $K$. Jedes 
$\alpha \in \End_L V_L$ lässt sich eindeutig schreiben als $\alpha = \sum_{i \in I} \lambda_i l_i$ mit $l_i \in \End_K V$. Es ist $\alpha \in \End_{L[G]} V_L$ genau dann,
wenn $\alpha$ mit der $G$-Wirkung vertauscht, genau dann wenn jedes $l_i$ mit der $G$-Wirkung vertauscht, d.h. wenn $l_i \in \End_{K[G]} V$. \bewende\smallskip\\
Ist $U \subseteq V$ ein $G$-invariante Unterraum, dann ist auch $U \otimes L \subseteq V_L$ $G$-invariant. Daher ist falls $\rho_L$ irreduzibel ist, ach $\rho$ irreduzibel.
Aber die Umkehrung gilt \emph{nicht}: 
\minisec{Beispiel}
Sie $G = C_4 = \langle g \rangle$, $K=\mathds{R}$, $V=\mathds{R}^2$, $\rho : G \to \GL_2(\mathds{R})$, $\rho(g) = \begin{psmallmatrix}
	0 & 1 \\
	-1 & 0
\end{psmallmatrix}$, Drehung um 90°. In $V$ gibt es keine $G$-invarianten Unterräume, d.h. $\rho$ ist irreduzibel. Aber über $\mathds{C}$ hat eine abelsche Gruppe keine
irreduzible Darstellung in Dimension $ >1$. Bei geeigneter Wahl einer Basis für $V_\mathds{C}$ wird $\rho_\mathds{C}(G)$ als Diagonalmatrix dargestellt.
\minisec{Bemerkung}
$\rho :G \to \GL_n(K)$ irreduzibel, $L \supseteq K$ $\not\Rightarrow$ $\rho_L : G \to \GL_n(L)$ irreduzibel. $\Leftrightarrow$ gilt aber!
% subsection 418 (end)

\subsection[Definition: Absolut irreduzible Darstellung]{Definition} % (fold)
\label{sub:420}
\todo[noline]{Hier stimmt die Nummerierung noch nicht!}%
Eine Darstellung $\rho : G \to \GL(V)$ für einen $K$-Vektorraum $V$ heißt \Index{absolut irreduzibel}, falls $\rho_L$ irreduzibel ist für jeden Körper $L \supseteq K$. 
% subsection 420 (end)

\subsection{Satz} % (fold)
\label{sub:421}
Sei $G$ eine endliche Gruppe, $\rho :G  \to \GL(V)$ eine Darstellung über $K$ mit $\Char K \nmid \abs*{G}$. Setze $A := K[G]$ und $A' = \End_A(V)= C_{\End_K(V)}(A)$. Dann
ist $\rho$ irreduzibel genau dann, wenn $A'$ ein Schiefkörper ist und absolut irreduzibel, wenn $A'=K$ gilt.
\minisec{Beweis}
Der erste Teil der Behauptung ist Schurs Lemma (\ref{sub:226}): Wenn $\rho$ irreduzibel ist, dann ist $A'$ ein Schiefkörper. Ist umgekehrt $\rho$ nicht irreduzibel, dann ist
$V= B \oplus C$ für $K[G]$-Untermoduln $B$ und $C$. Die Projektion $\pi_B :V \to B$ ist $K[G]$-linear, d.h. $\pi_B \in A'$, aber $\pi_B \not= 0$ und $\pi_B$ ist nicht 
invertierbar, d.h. $A'$ ist kein Schiefkörper.

Ist nun $\rho$ absolut irreduzibel und $K \subsetneq A'$, dann sei $c \in A' \setminus K$. $f $ sei das Minimalpolynom von $c$ über $K$, d.h. $\deg f > 1$ und $f$ ist
irreduzibel, weil $A'$ ein Schiefkörper ist. Setze
\[
	L= \nicefrac{K[X]}{(f)}
\]
und betrachte $A'_L$. Dann hat $f$ eine Nullstelle in $L$, d.h. $A'_L$ ist kein Schiefkörper, weil $A'_L$ Nullteiler enthält \light $\rho$ absolut irreduzibel. Ist umgekehrt
$A'=K$, dann ist für jeden Körper $L \supseteq K$ natürlich $A'_L = K \otimes L = L$ ein Körper, also $\rho$ absolut irreduzibel. \bewende
% subsection 421 (end)

\subsection[Definition: Zerfällungskörper einer Gruppe]{Definition} % (fold)
\label{sub:422}
Sei $G$ eine endliche Gruppe, $K$ ein Körper. Dann heißt $K$ \Index{Zerfällungskörper} für $G$, wenn jede Darstellung von $G$ über $K$ absolut irreduzibel ist. 
% subsection 422 (end)

\subsection{Satz} % (fold)
\label{sub:423}
Sei $G$ eine endliche Gruppe, $K$ ein Körper mit $\Char K \nmid \abs*{G}$. Dann ist $K$ Zerfällungskörper für $G$ genau dann, wenn $K[G]$ die direkte Summe von 
Matrizengruppen $M_n(K)$ ist.
\minisec{Beweis}
Zerlege $A=K[G]= \bigoplus A_i$, wobei $A_i$ einfach ist. Sei $J_i \subseteq A_i$ ein minimales Linksideal in $A_i$. Die irreduziblen Darstellungen von $G$ über $K$
entsprechen genau den Darstellungen $\rho_i$ von $G$ auf $J_i$. Also ist $K$ genau dann Zerfällungskörper, wenn alle $\rho_i$ absolut irreduzibel sind. Nach Satz 
\ref{sub:422} ist dies genau dann der Fall, wenn 
\[
	D_i' = \End_A J_i = K \quad \text{ für alle }i.
\]
Wegen der Irreduzibilität von $J_i$ ist $D_i'$ ein Schiefkörper und wegen der Doppelzentralisatoreigenschaft (Satz \ref{sub:410}) gilt
\[
	\End_{D_i'} J_i = \rho_i(A) \simeq M_{n_i}(D_i)
\]
für $D_i = (D_i')^\op$ und $n_i = \dim_{D_i'} J_i$. Da die $A_i$ einfach sind, ist $\rho_i(A)=A_i$, d.h. $A_i \simeq M_{n_i}(D_i)$. Wenn $K$ ein Zerfällungskörper ist, dann
sind die $\rho_i$ absolut irreduzibel, d.h. $D_i' = K$ und $A_i \simeq M_{n_i}(K)$.

Ist $A_i \simeq M_{n_i}(K)= M_{n_i}(D_i)$, dann folgt $K=D_i$, insbesondere $D_i'=K$ und nach Satz \ref{sub:421} sind die $\rho_i$ absolut irreduzibel, d.h.
$K$ ist Zerfällungskörper. \bewende
% subsection 423 (end)

\subsection[Definition: Charakter]{Definition} % (fold)
\label{sub:424}
Eine $1$-dimensionale Darstellung $G \to \GL_1(K)=K^*$ heißt \Index{linearer Charakter}. Ist $\rho : G \to \GL(V)$  eine Darstellung, dann heißt $\chi_\rho : G \to K$,
$\chi_\rho(g)= \Tr \enbrace[\big]{\rho(g)}$ der zu $\rho$ gehörige \Index{Charakter}, d.h. ist$p(x)= \enbrace*{\rho_{ij}(x)} $ so ist
\[
	\Tr \enbrace[\big]{ p(x)} = \sum_{i=1}^{\dim V} \rho_{ii} (x)
\]
\minisec{Bemerkung}
$\chi_\rho$ ist wohldefiniert, weil $\Tr$ nicht von der Basiswahl abhängt. Jede $1$-dimensionale Darstellung ist ihre eigener Charakter. Für abelsche Gruppen über
algebraisch abgeschlossenen Körpern sind die Charaktere genau die irreduziblen Darstellungen.
% subsection 424 (end)

\subsection{Proposition} % (fold)
\label{sub:425}
\begin{enumerate}[(i)]
	\item Äquivalente Darstellungen haben denselben Charakter.
	\item Jeder Charakter ist eine Klassenfunktion, d.h. konstant auf Konjugationsklassen von $G$.
	\item Ist $\Char K =0$, dann $\chi_\rho(1)= \dim V$.
	\item Sind $\rho_1, \rho_2$ Darstellungen mit Charakteren $\chi_i$ und $\chi_2$, dann sind $\rho_1 \oplus \rho_2$, $\rho_1 \otimes \rho_2$ Darstellungen  mit Charakter
	$\chi_1 + \chi_2$ bzw. $\chi_1 \cdot \chi_2$
	\item Ist $K= \mathds{C}$, $x \in G$ mit $x^n= 1_G$, dann ist $\chi(x)$ Summe von $n$-ten Einheitswurzeln in $\mathds{C}$.
\end{enumerate}
\minisec{Beweis}
\begin{enumerate}[(i)]
	\item Klar: Sind $\rho_1, \rho_2$ äquivalent via dem Vektorraumisomorphismus $\varphi$, dann ist auch $\rho_2(g)= \varphi \rho_1(g) \varphi ^{-1}$, also 
	$\Tr \rho_2(g) = \Tr \rho_1(g)$.
	\item Für $g,h \in G$ ist 
	\[
		\Tr \enbrace*{ \rho \enbrace*{h g h ^{-1}} } = \Tr \enbrace*{\rho(h) \rho(g) \rho(h) ^{-1}} = \Tr \enbrace[\big]{\rho(g)}   
	\]
	\item Es ist $\rho(1_G)= \id_V$
	\item Es ist $\Tr(A \oplus B)= \Tr(A) + \Tr(B)$ und $\Tr(A \otimes B)= \Tr(A) \cdot \Tr(B)$
	\item Ist $x^n = 1_G$ für $x \in G$, dann gilt für $A= \rho(x)$ ebenfalls $A^n= \id_V$, d.h. das Minimalpolynom von $A$ teilt das Polynom $X^n-1$, mit Wurzeln
	$\omega_1, \ldots , \omega_n \in \mathds{C}$. Daher ist $A$ eine Diagonalmatrix mit Einträgen aus $\set{\omega_1, \ldots , \omega_n}$. \bewende
\end{enumerate}
% subsection 425 (end)

\subsection{Beispiel} % (fold)
\label{sub:426}
\begin{enumerate}[(i)]
	\item Sei $\rho : G \to \GL_1(V)$ die triviale Darstellung mit $\rho(g)=1$, dann heißt $\chi_\rho : G \to K^*$, $g  \mapsto 1$  \Index{trivialer Charakter}.
	\item Sei $G$ endlich, $\abs*{G}=n$, $\rho : G \to \GL_n(K)$ die reguläre Darstellung durch Permutationsmatrizen. Dann ist 
	\[
		\chi_\rho(g) = \begin{cases}
			n, &\text{ falls }g= 1_G\\
			0 , &\text{ sonst}
		\end{cases}
	\]
\end{enumerate}
% subsection 426 (end)

\subsection{Satz (Orthogonalitätsrelationen für irreduzible Darstellungen)} % (fold)
\label{sub:427}
Sei $G$ eine endliche Gruppe, $K$ der Zerfällungskörper für $G$, $\Char K = 0$\marginnote{in der \enquote{Praxis}: $K= \mathds{C}$}. Sind $\rho, \sigma$ irreduzible 
Darstellungen der Dimension $c$ bzw. $d$, dann gilt für alle $i,j \le c$ und $p,q \le d$
\[
	\sum_{g \in G} \rho_{ij}(g) \cdot \sigma_{pq}(g ^{-1}) = \begin{cases}
		0, &\text{ falls }\rho, \sigma \text{ nicht äquivalent}\\
		\frac{\abs*{G} }{d} \cdot \delta_{jp} \cdot \delta_{iq} , &\text{ falls } \rho = \sigma 
	\end{cases}
\]
wobei $\rho(g)= \enbrace[\big]{\rho_{ij }(g)}_{i,j \le c} $, $\sigma(g)= \enbrace[\big]{\sigma_{pq}(g)}_{p,q \le d}$.
\minisec{Beweis}
Seien $\rho : G \to \GL(U)$, $\sigma : G \to \GL(V)$, $(u_1, \ldots , u_c)$ bzw. $(v_1, \ldots , v_d)$ Basen von $U$ bzw. $V$. Sei weiter 
\[
	\alpha_{jp} : U \to V \quad , \quad u_i \mapsto\delta_{ij} \cdot v_p = \begin{cases}
		u_j \mapsto v_p, &\text{ für } i=j\\
		u_i \mapsto 0 & \text{ für } i \not= j
	\end{cases}
\]
$\alpha_{jp}$ ist $K$-linear und $\Tr(\alpha_{jp}): U \to V$ ist $K[G]$-linear. Sind $\rho$ und $\sigma$ nicht äquivalent, dann folgt $\Tr(\alpha_{jp})=0$ für alle $j\le c$,
$p \le d$. (siehe 2. Satz von Wedderburn, \ref{sub:231}). Sind $\rho \sim \sigma$, also $U \simeq V$ $K[G]$-Moduln, dann ist $\Tr(\alpha_{jp})= \lambda_{jp} \in K$, denn
$U$ ist einfach und $\End_{K[G]} U = K$.

Es ist $(\alpha_{jp})_{iq} = \delta_{ij} \cdot \delta_{pq}$, also 
\[
	\enbrace[\big]{\Tr(\alpha_{jp})}_{iq} = \sum_{\substack{h \le c, r \le d \\ g \in G}} \rho_{ij}(g) \cdot \delta_{hj} \cdot \delta_{pr} \cdot \sigma_{rq}(g) ^{-1}
	= \sum_{g \in G} \rho_{ij}(g) \cdot \sigma_{pq}(g) ^{-1}
\]
Falls $\rho \not\sim \sigma$, folgt die Behauptung. Sei nun $\rho = \sigma$. Dann ist
\begin{align*}
	\sum_g \rho_{ij}(g) \cdot \rho_{pq}(g) ^{-1} = \lambda_{jp} \cdot \delta_{iq} 
\end{align*}
Für $i=q$ ergibt sich als Summe $i=1, \ldots ,d =c$
\begin{align*}
	d \lambda_{jp} = \sum_i \sum_g \rho_{pi}(g) ^{-1} \cdot \rho_{ij}(g) = \sum_g \rho_{pj}(1) = \abs*{G} \cdot \delta_{jp} 
\end{align*}
Daher gilt $\lambda_{jp}= \frac{\abs*{G} }{d} \delta_{jp}$ und damit 
\[
	\sum_{g} \rho_{ij}(g) \cdot \sigma_{pq}(g) ^{-1} = \frac{\abs*{G} }{d} \cdot \delta_{jp} \cdot \delta_{iq}. \bewende
\]
% subsection 427 (end)

\subsection{Korollar} % (fold)
\label{sub:428}
Sei $G$ eine endliche Gruppe, $K$ ein Zerfällungskörper, $\Char K  \nmid \abs*{G}$, $\set{\rho_1, \ldots , \rho_s}$ Repräsentanten für die Äquivalenzklassen der 
irreduziblen Darstellung über $K$. Für jedes $i=1, \ldots , s$ sei $\rho^{(i)}$ eine Matrixdarstellung für $\rho_i$. Dann gilt 
\[
	\sum_g \rho_{ij}^{(k)}(g) \cdot \rho_{pq}^{(l)}(g) ^{-1} = \begin{cases}
		0, &\text{ falls }k \not= l\\
		\frac{\abs*{G} }{\deg \rho_k} \cdot \delta_{jp} \cdot \delta_{iq} , &\text{ falls } k=l 
	\end{cases}
\]
% subsection 428 (end)

\subsection{Korollar} % (fold)
\label{sub:429}
Sind $\chi_i$, $i=1, \ldots ,s$ die zu $\rho_i$ gehörigen Charaktere, dann gilt
\[
	\sum_g \chi_i(g) \cdot \chi_j(g) ^{-1} = \begin{cases}
		0, &\text{ falls }i \not= j\\
		\abs*{G}, &\text{ falls } i=j 
	\end{cases}
\]
Ist speziell $K=\mathds{C}$, dann folgt wegen Proposition \ref{sub:425}(v)
\[
	\sum_g \chi_i(g) \cdot \chi_j(g) ^{-1} = \sum_g \chi_i(g) \cdot \overline{\chi_j(g)} = \begin{cases}
		0, &\text{ falls }i \not= j\\
		\sum_g \chi_i(g) \cdot \overline{\chi_i(g)} = \sum_g \abs*{\chi_i(g)^2} = \abs*{G} , &\text{ falls } i=j   
	\end{cases}
\]
Betrachte $K[G]= K^G$ als $K$-Vektorraum aller $K$-wertigen Funktionen auf $G$. Definiere inneres Produkt auf $K[G]$ durch 
\[
	\enbrace*{\sum_{x \in G} \alpha_x x, \sum_{y \in G} \beta_y y} = \frac{1}{\abs*{G}} \sum_{x \in G} \alpha_{x ^{-1}} \beta_x \in K ,
\]
falls $\Char K \nmid \abs*{G}$. Für $K=\mathds{C}$ erhalten wir eine hermitesche Form durch
\[
	(\varphi, \psi) = \frac{1}{\abs*{G}}  \sum_g \overline{\varphi(g)} \psi(g) .
\]
Dann ist 
\begin{align*}
	(\varphi, \varphi) = \frac{1}{\abs*{G}} \sum_g \overline{\varphi(g)} \varphi(g) = \frac{1}{\abs*{G}} \sum_g \abs*{\varphi(g)}^2 \ge 0   
\end{align*}
und $(\varphi, \varphi)=0 \iff \varphi=0$, d.h. $(\_, \_)$ ist positiv definit und regulär. Nach Korollar \ref{sub:429} gilt $(\chi_i, \chi_j)= \delta_{ij}$, d.h. die
Charaktere der irreduziblen Darstellung $\chi_i , \ldots \chi_s$ bilden eine Orthonormalbasis des $\mathds{C}$-Vektorraumes der Klassenfunktionen.

Ist nun $\rho : G \to \GL(\mathds{C})$ eine beliebige Darstellung mit Charakter $\chi$, dann ist $\rho \sim m_1 \rho_i \oplus \ldots \oplus m_s \rho_s$ und daher
\[
	\chi = m_i \cdot \chi_i \ldots + m_s \cdot \chi_s.
\]
Daher ist $(\chi_i, \chi_i)= m_i$, $i=1, \ldots ,s$.
% subsection 429 (end)

\subsection{Korollar} % (fold)
\label{sub:430}
Komplexe Darstellungen einer endlichen Gruppe sind genau dann äquivalent, wenn sie denselben Charakter haben.
% subsection 410 (end)

\subsection{Beispiel} % (fold)
\label{sub:431}
Bestimme alle irreduziblen Charaktere (bzw. Darstellungen) von $S_3$. Es gilt $\abs*{S_3}=6$. $S_3$ hat drei Konjugationsklassen mit Vertretern 
\[
	\id, \begin{matrix}
		(1  \enspace 2) \\
		(1 \enspace 3) \\
		(2 \enspace 3)
	\end{matrix} ,
	\begin{matrix}
		(1 \enspace 2 \enspace 3) \\
		(2 \enspace 3 \enspace 1)
	\end{matrix}
\]
Es gibt zwei $1$-dimensionale Darstellung von $S_3$ mit Charakter $\chi_i(g)=1$ und $\chi_2(g)= \sgn(g)$. Wegen $6= 1^2 + 1^2 + 2^2$ existiert noch eine $2$-dimensionale
Darstellung, d.h. es gilt $\chi_3(1)=2$. Nun muss $\chi_3$ auf $\chi_1$ und $\chi_2$ senkrecht sein, d.h. 
\begin{align*}
	(\chi_1, \chi_3) &= \frac{1}{6} \enbrace[\big]{2+3 \chi_3( (1 \enspace2)) + 2 \chi_3((1 \enspace 2 \enspace 3))} = 0 \\
	(\chi_2, \chi3) &= \frac{1}{6} \enbrace[\big]{2 - 3 \chi_3((1 \enspace 2)) + 2 \chi_3((1 \enspace 2 \enspace 3))} =0  
\end{align*}
$\Rightarrow \chi_3(1 \enspace 2) = 0$ und $\chi_3(1 \enspace 2 \enspace 3)= -1$. Damit ergibt sich die Charaktertafel
\begin{table}[H]
	
\end{table}
Und für größere $n$ die Young-Tableaus.
\section*{Dualität für endliche abelsche Gruppen} % (fold)
\label{sec:dualitaet}

\subsection{Definition} % (fold)
\label{sub:432}
Sei $G$ eine abelsche Gruppe, dann heißt $G^* = \Hom(G,\mathds{C}^*)$ die \Index{Charaktergruppe} (oder duale Gruppe)  von $G$
unter der Multiplikation: $\varphi, \psi \in G^*$, 
\[
	\varphi \cdot \psi : G \to \mathds{C}^* , g \mapsto \varphi(g) \cdot \psi(g)
\]
\minisec{Bemerkung}
Ist $G= \nicefrac{\mathds{Z}}{n \mathds{Z}}= C_n$ zyklisch mit $G= \langle s \rangle$, $\chi : G \to \mathds{C}^*$, dann ist $\chi(s)$ eine $n$-te Einheitswurzel.
Ist $\omega^* \in \mu_n$ eine $n$-te Einheitswurzel, dann ist $\chi(s^i) = \omega^i$ ein Charakter. Daher ist hier $C_n^* = \mu_n \simeq C_n$ Gruppe der $n$-ten 
Einheitswurzeln.
% subsection 432 (end)

\subsection{Satz} % (fold)
\label{sub:433}
Jede endliche abelsche Gruppe $A$ ist isomorph zur dualen Gruppe $A^*$.
\minisec{Beweis}
Nach dem Hauptsatz über endlich erzeugte Moduln über Hauptidealringen ist $A \simeq A_1 \oplus  \ldots \oplus A_r$ mit $A_i$ zyklisch. 
Offensichtlich gilt
\[
	\hom(B \times C, \mathds{C}^*) \cong \hom(B, \mathds{C}^*) \times \hom(C,\mathds{C}^*).
\]
Daher ist $(B \times C)^* \cong B^* \times C^*$. Also 
\[
	A^* \simeq A_1^* \times \ldots \times A_r^* \simeq A_1 \times \ldots \times A_r = A. \bewende
\]
% subsection 433 (end)
% section dualitaet (end)
% subsection 431 (end)
% section 4 (end)










\cleardoubleoddemptypage
\appendix
\section{Anhang} % (fold)
\label{sec:anhang}
\subsection{Alternative Definition von Gruppenwirkungen} % (fold)
\label{sub:anh_1}
Eine \Index{Gruppenwirkung} von $G$ auf $X$ ist (gegeben durch) eine Abbildung $G\times X \to X$, $(g,x)\mapsto g\cdot x$, mit folgenden Eigenschaften:
\begin{itemize}
	\item $e_G \cdot x= x$, und
	\item $g\cdot(h\cdot x)=(gh)\cdot x$,
\end{itemize}
für alle $x \in X$ und alle $g,h\in G$.
\minisec{Rechts-Gruppenwirkung}
Die oben angegebenen Definitionen beschreiben eine \bet{linksseitige Gruppenwirkung}\index{Gruppenwirkung!linksseitige}.\\
Eine \bet{rechtsseitige Gruppenwirkung}\index{Gruppenwirkung!rechtsseitige} ist (gegeben durch) eine Abbildung $X\times G \to X$, $(x,g)\mapsto x\cdot g$, 
wobei für alle $x\in X$ und alle $g,h\in G$ die Bedingungen $x\cdot e=x$ und $(x\cdot g) \cdot h =x \cdot (gh)$ erfüllt sind.\\
Alternativ kann eine rechtsseitige Gruppenwirkung durch einen Gruppenhomomorphismus $\varphi : G \to \Sym(X)^\op$ beschrieben werden, 
wobei $\Sym(X)^\op:=(\Sym(X),*)$ mit $f*g:=g\circ f$ ist. (vgl. Definition \ref{sub:228}).
% subsection anh_1 (end)

\subsection{Beweis für \ref{sub:26}(a) von Blatt 3} % (fold)
\label{sub:anh_2}
Sei $M$ ein zyklischer $R$-Modul. Sei $0 \not= x \in M$. Dann gilt $\langle x \rangle = M$. Betrachte den Homomorphismus $\varphi : R \to M$, $r \mapsto r \cdot x$.
$\varphi$ ist surjektiv, da $\langle x \rangle= M$. Dann gilt
\[
	\nicefrac{R}{\ker \varphi} \simeq M
\]
Dann ist $J := \ker \varphi$ ein Ideal in $R$. Zu zeigen: $J$ ist maximal. Annahme: Es existiert ein echtes Ideal $I \subset R$ mit $J \subsetneq I$. Dann ist 
$\nicefrac{I}{J}$ ein Ideal in $\nicefrac{R}{J} \simeq M$. Dann ist $\nicefrac{I}{J}$ ein Untermodul von $M$. Da $M$ zyklisch ist und $\nicefrac{I}{J}$ ein nichtriviales 
Element enthält, gilt $\nicefrac{I}{J}=M$. Dann folgt aber $I=R$ und wir erhalten einen Widerspruch. \bewende
% subsection anh_2 (end)
% section anhang (end)
\newpage
\pagenumbering{Alph}
\setcounter{page}{1}

\printindex
\listoffigures
\todototoc
\listoftodos[Todo's und andere Baustellen]
\end{document}
